
\documentclass[11pt]{article}

\usepackage[margin=1.02in]{geometry}
\usepackage{amsmath,amssymb,amsthm,mathtools,mathrsfs}
\usepackage{bm}
\usepackage{enumitem}
\usepackage{booktabs,array,tabularx}
\usepackage{xcolor}
\usepackage{microtype}
\usepackage[colorlinks=true,linkcolor=blue,citecolor=blue,urlcolor=blue]{hyperref}
\usepackage[nameinlink,capitalise]{cleveref}
\setlength{\emergencystretch}{4em}
\raggedbottom
\hbadness=10000

% -----------------------------------------------------------------------------
% theorem environments
% -----------------------------------------------------------------------------
\newtheorem{theorem}{Theorem}[section]
\newtheorem{proposition}[theorem]{Proposition}
\newtheorem{lemma}[theorem]{Lemma}
\newtheorem{corollary}[theorem]{Corollary}
\newtheorem{countertheorem}[theorem]{Countertheorem}
\newtheorem{example}[theorem]{Example}
\theoremstyle{definition}
\newtheorem{definition}[theorem]{Definition}
\newtheorem{construction}[theorem]{Construction}
\newtheorem{assumption}[theorem]{Assumption}
\theoremstyle{remark}
\newtheorem{remark}[theorem]{Remark}
\newtheorem{interpretation}[theorem]{Interpretation}
\newtheorem{scope}[theorem]{Scope statement}

% -----------------------------------------------------------------------------
% notation
% -----------------------------------------------------------------------------
\newcommand{\N}{\mathbb N}
\newcommand{\Z}{\mathbb Z}
\newcommand{\R}{\mathbb R}
\newcommand{\C}{\mathbb C}
\newcommand{\one}{\mathbf 1}
\newcommand{\id}{\mathrm{id}}
\newcommand{\dd}{\mathrm d}
\newcommand{\Tr}{\operatorname{Tr}}
\newcommand{\tr}{\operatorname{tr}}
\newcommand{\rank}{\operatorname{rank}}
\newcommand{\Ker}{\operatorname{Ker}}
\newcommand{\Ran}{\operatorname{Ran}}
\newcommand{\Span}{\operatorname{span}}
\newcommand{\End}{\operatorname{End}}
\newcommand{\Hom}{\operatorname{Hom}}
\newcommand{\supp}{\operatorname{supp}}
\newcommand{\spec}{\operatorname{spec}}
\newcommand{\diag}{\operatorname{diag}}
\newcommand{\Aut}{\operatorname{Aut}}
\newcommand{\U}{\mathrm U}
\newcommand{\SU}{\mathrm{SU}}
\newcommand{\Cl}{\mathrm{Cl}}
\newcommand{\HS}{\mathrm{HS}}
\newcommand{\norm}[1]{\lVert#1\rVert}
\newcommand{\abs}[1]{\lvert#1\rvert}
\newcommand{\ip}[2]{\left\langle#1,#2\right\rangle}
\newcommand{\A}{\mathcal A}
\newcommand{\Hh}{\mathcal H}
\newcommand{\M}{\mathcal M}
\newcommand{\Kk}{\mathcal K}
\newcommand{\Oalg}{\mathcal O}
\newcommand{\GT}{\boldsymbol{\mathbb T}_{\mathrm{Grand}}}
\newcommand{\GTX}{\boldsymbol{\mathbb T}_{\mathrm{Grand},X}}
\newcolumntype{L}[1]{>{\raggedright\arraybackslash}p{#1}}

\title{\textbf{Spacetime--Gauge Commutant Duality: Finite Rigidity and Cofinal Stability}}

\author{%
  {\small $^{\dagger}$Correspondence: \href{mailto:aurelien.pelissier.38@gmail.com}{aurelien.pelissier.38@gmail.com}}%
}
\date{}

\begin{document}
\maketitle

\begin{abstract}
Renewal Geometry reconstructs Lorentzian spacetime and spectral geometry as distinct effective descriptions of a common predictive structure. Here we study their joint finite realization and show that spacetime action and matter multiplicity form an exact double-centralizer pair. The duality is not between the bare spacetime Clifford algebra and the Standard-Model gauge algebra, but between the complete spacetime--gauge--incidence action and the residual multiplicity invisible to that action. This relation admits a positive finite rigidity certificate and remains stable under controlled regulator refinement.
%
On a certified active branch, removing the complete spacetime source leaves an internal structure whose multiplicities reconstruct the Standard-Model nonabelian sectors. A canonical determinant incidence selects $
S(\U(3)\times\U(2))
\cong
(\SU(3)\times\SU(2)\times\U(1))/\Z_6,
$ while anomaly cancellation fixes the hypercharge ratios. A distinct rank-three residual provides a three-generation carrier, and a single grading-odd operator contains the finite Yukawa and Majorana sectors in a common generation frame.
%
The construction is structural rather than numerical: couplings, masses, flavour parameters, renormalization-group flow, and the fermionic measure phase remain independent physical data.
\end{abstract}


\medskip
\noindent\textbf{Keywords.}
Renewal Geometry; Standard Model; spacetime; gauge symmetry; commutant duality;
double centralizer; Howe duality; noncommutative geometry; Clifford algebra;
generations; hypercharge.

% =============================================================================
\section{Introduction}
\label{sec:introduction}
% =============================================================================

The Standard Model and general relativity organize physical structure in very
different ways.  General relativity geometrizes gravity by identifying the
metric with the dynamical field of spacetime.  The Standard Model instead
starts from an internal gauge group, chiral representations, a scalar doublet,
and Yukawa couplings, all defined over an already existing spacetime.  Under
the standard analytic $S$-matrix hypotheses, the Coleman--Mandula theorem
formalizes the familiar separation of continuous spacetime and internal
symmetries \cite{ColemanMandula1967}.  A long-standing structural question is
nevertheless whether the two sectors, while remaining distinct as physical
symmetries, can arise from and constrain one another through a common
underlying organization.

Noncommutative geometry gives one of the most developed geometric answers to
this question.  Early Connes--Lott models already formulated electroweak and
Standard-Model sectors as noncommutative gauge theories
\cite{ConnesLott1991}; in the spectral formulation, geometry is encoded by an
algebra, a Hilbert-space representation, and a Dirac operator
\cite{Connes1994,ConnesReality1995}.  Almost-commutative geometries combine an
ordinary spin manifold with a finite noncommutative space, and the spectral
action recovers the Standard Model coupled to gravity in the corresponding
framework
\cite{ConnesGravityMatter1996,ChamseddineConnes1997,
ChamseddineConnesMarcolli2007,ChamseddineConnes2008,vanSuijlekom2025}.
In that construction the continuum spacetime factor and the finite internal
factor are specified as parts of the spectral datum.  The finite sector itself
has a substantial classification theory.  Krajewski diagrams encode finite
spectral triples in terms of gauge-multiplet vertices and Yukawa links, and
subsequent classifications impose irreducibility and dynamical
nondegeneracy \cite{Krajewski1998,IochumSchuckerStephan2004}.  The
almost-commutative formulation and its particle-physics interpretation are
reviewed systematically in \cite{vanDenDungenVanSuijlekom2012,vanSuijlekom2025};
relaxing spectral-triple conditions can also enlarge the internal symmetry,
for example to Pati--Salam type models
\cite{ChamseddineConnesVanSuijlekom2013}.  Related algebraic reformulations
have sought to make the relation among differential forms, spinors, and the
finite Standard-Model data more intrinsic \cite{BoyleFarnsworth2018}.  In the
conventional Standard-Model finite triple the observed family multiplicity is
implemented by taking three copies of the fermionic finite Hilbert space rather
than being derived from the spacetime factor
\cite{vanDenDungenVanSuijlekom2012,vanSuijlekom2025}.

There is also a directly relevant commutant literature inside noncommutative
geometry.  Morita and Hodge formulations ask when the fermionic Hilbert space
is an equivalence bimodule for the spectral Clifford algebra, leading to
commutant identities for the Clifford action
\cite{DAndreaDabrowski2016,DabrowskiDAndreaSitarz2018}; in the Standard Model,
fermion masses and mixing constrain this Hodge property
\cite{DabrowskiSitarz2019}.  These results are close in algebraic form to the
duality studied below, but the two statements should not be conflated.  Their
commutant is built from the left/right spectral Clifford structures of an
already specified finite spectral triple.  Here the mutually centralizing pair
is instead constructed only after the independently reconstructed Lorentzian
Clifford action, typed internal action, and physical matter-incidence maps are
placed on one predictive carrier; its opposite member is the residual
flavour/generation multiplicity algebra.  Gauge symmetry of finite
noncommutative spaces has likewise been studied intrinsically at the
operator-algebraic level \cite{BhowmickEtAl2014}.

Lorentzian versions of the noncommutative Standard Model provide another
important comparison.  Barrett formulated a Lorentzian finite geometry for
the Standard Model \cite{Barrett2007}, while twisted spectral triples can
produce a Krein structure and Lorentzian signature from the spectral data
\cite{DevastatoEtAl2018}.  Outside spectral NCG, Clifford-algebraic approaches
have also sought structural relations between spacetime and internal gauge
symmetry.  Trayling and Baylis derive Standard-Model gauge transformations from
a $Cl_7$ geometric-algebra setting while keeping track of exterior spacetime
and interior actions \cite{TraylingBaylis2001}; related Clifford constructions
recover the one-generation gauge content in $Cl_{1,5}$ or from minimal ideals
of complex Clifford algebras \cite{Daviau2017,Gresnigt2020}.  These are direct
precedents for deriving internal structure from an enlarged Clifford carrier,
but they do not formulate the support-exact mutual-commutant pair studied here.
The present construction differs in logical order: the Lorentzian spacetime
packet and the spectral/internal packet are separate companion outputs of one
upstream predictive object before their joint commutant is formed.

Renewal Geometry takes a different logical route.  Its microscopic entrance is
a finite operational process of composable histories and conditional future
statistics.  Histories that have identical statistics under every admissible
future continuation are identified, producing a canonical minimum predictive
carrier.  Complete positive multitime statistics reconstruct a finite history
algebra and its intrinsic positive moment hierarchy, together forming the
Renewal Grand Tensor.  The companion spacetime work
\cite{PelissierSpacetime2026} showed that, on a specified minimal faithful
branch, calibrated renewal statistics reconstruct a Lorentzian $3+1$ ADM
geometry and a common-action stationary limit satisfying the vacuum Einstein
equation.  Independently, the companion spectral work
\cite{PelissierSpectral2026} showed that a compatible first-order variation of
the reconstructed history algebra produces Connes-type spectral geometry and
that spectralization is generally a many-to-one reduction of the underlying
predictive dynamics.

Those results leave a natural joint question.  Once the Lorentzian and
spectral sectors are represented on the \emph{same} source-minimal predictive
carrier, what algebraic relation connects spacetime, gauge structure, and
matter multiplicity?  A tempting answer would be to identify the Standard
Model gauge algebra with the commutant of the spacetime Clifford algebra.  The
first result of the present paper is that this cannot be correct.  A unital
copy of $M_4(\C)$ acting on a finite Hilbert space has a full matrix algebra as
its commutant on each multiplicity block, whereas the typed internal
Standard-Model action is reducible.  The duality therefore cannot be
\[
\text{spacetime Clifford algebra}
\quad\longleftrightarrow\quad
\text{Standard-Model gauge algebra}.
\]

\noindent
Throughout this paper, ``spacetime--gauge duality'' therefore refers to a
\emph{joint action--multiplicity double-centralizer relation}.  It never means
that the bare spacetime Clifford algebra and the physical gauge algebra are
ordinary commutants of one another.

The correct duality appears one layer later.  The external Clifford action, the
typed internal action, and the matter-incidence maps generate a \emph{joint
typed action algebra}.  Its commutant is precisely the residual multiplicity
algebra, and finite bicommutant closure gives the converse.  In the terminology
of invariant theory this is a finite Howe-type situation: two commuting
operator algebras determine one another as mutual commutants on the represented
carrier \cite{Howe1989}.  Mutual-centralizer ideas also have direct precedents
in quantum field theory.  Doplicher--Roberts reconstruction recovers a compact
gauge group and field algebra from superselection structure
\cite{DoplicherRoberts1989,DoplicherRoberts1990}, and Todorov exhibited
observable--gauge dual pairs in the sense of Howe in four-dimensional
conformal field-theory models \cite{Todorov2010}.  Our pair is neither a
superselection-category reconstruction nor an observable--gauge dual pair:
both members are finite algebras reconstructed on the same source-complete
spacetime--matter carrier.  We therefore use the term \emph{Howe-type duality}
only in this finite mutual-centralizer sense; no oscillator representation or
reductive-dual-pair hypothesis is required.  The terminology is also unrelated
to holographic gauge/gravity duality in the AdS/CFT sense
\cite{Maldacena1998,Witten1998}.

The first main theorem is therefore
\begin{equation}
\boxed{
\A_{{\rm act},X}'
=
\M_{{\rm type},X},
\qquad
\M_{{\rm type},X}'
=
\A_{{\rm act},X}.}
\label{eq:intro-main-duality}
\end{equation}
The statement is support exact.  Singular or rectangular matter-incidence maps
are treated by their canonical polar supports, positive support metrics, and
partial isometries; no inverse Yukawa matrix is introduced.  If the support
projections fail to commute, the same data generate a canonical
operator-space quiver whose endomorphism algebra is the surviving
flavour/generation algebra.  On the connected invertible branch, all
multiplicity spaces can be transported to one root fibre and the dual algebra
is generated there by positive polar metrics and cycle holonomies.

A second result makes the duality testable at finite cutoff.  If
$c_1,\ldots,c_s$ generate the action algebra, the commutant Laplacian
\[
\mathscr L_{\bm c}
=
\sum_{j=1}^{s}\operatorname{ad}_{c_j}^*
                 \operatorname{ad}_{c_j}
\]
has kernel exactly the commutant.  Restricting this positive form to the
orthogonal complement of a proposed protected multiplicity algebra gives a
\emph{relative Howe Gram}.  Positive definiteness is equivalent to exactness
of the proposed multiplicity algebra, while its least eigenvalue is a
quantitative rigidity margin.  A single positive finite anchor plus a
controlled generator tail gives an explicit perturbative stability estimate.
Nonperturbatively, transported Mosco convergence together with asymptotic
compactness of the commutator forms implies norm-resolvent convergence,
convergence of the isolated commutant spectral projection, and cofinal locking
of the protected commutant with a positive limiting gap.

The third result identifies an explicitly certified finite Standard-Model realization
inside this duality.  The tetrahedral spacetime construction carries twelve
oriented $A_3$ roots.  The complete spacetime source is removed by its canonical
positive Schur short.  On the explicit orthogonal active realization this short
leaves a faithful twelve-dimensional internal residual, while on a general
common carrier the same residual is measured by the positive shorted Gram
$S_{\rm int}^*(I-P_{\rm ST})S_{\rm int}$.  The categorical residual decomposes
under tetrahedral permutations and root reversal as
\begin{equation}
\C_{\rm type}^{+}
\oplus W_+
\oplus (V_2)_+
\oplus W_-
\oplus (W\otimes{\rm sgn})_- .
\label{eq:intro-residual-decomposition}
\end{equation}
The direct tetrahedral Burnside blocks are external.  The internal factors are
controlled instead by the multiplicities of these external irreducibles in the
physical word carrier.  This gives a finite selection hierarchy: the populated
degree-one carrier cannot support an internal $M_3(\C)$ factor; a colour factor
can first occur only after an external-trivial multiplicity of at least three
has appeared; and, once the multiplicity census is $(3,2,1,1)$, one connected
colour bridge and one noncommuting weak router are individually source-minimal.
A module-recognition theorem distinguishes invariant endomorphisms of a regular
word space from operations induced by the historical system algebra.  On a
common irreducible mediator, reciprocal source routes moreover force a colour
bridge through a finite positive return Gram, so the bridge need not be an
independent primitive on that route class.  Finally, a central-separation test
excludes a possible diagonal identification of the two scalar blocks before the
full relative commutant is asserted.  Under these support-separation hypotheses
the four bridge/router deficit dimensions are exactly $9,11,13,15$, with the
terminal value giving
\[
M_3(\C)\oplus M_2(\C)\oplus\C\oplus\C.
\]
Only after this multiplicity-side algebra has landed is the determinant
incidence tested.  Its nonzero alternating component has stabilizer
$S(\U(3)\times\U(2))$.  This global form of the Standard-Model group has a
long independent history \cite{Hucks1991,Tong2017}; in ordinary field theory
the Lie algebra alone does not experimentally select the full $\Z_6$ quotient,
whereas the determinant seed below selects it on the present structural branch.
Its central one-parameter subgroup is then fixed, within the displayed
tensor-type packet, by anomaly cancellation, yielding the usual hypercharge
ratios.  This is consistent with the classical anomaly/charge-quantization
analysis \cite{GengMarshak1989,MinahanRamondWarner1990} and with the relation
between unimodularity and anomaly cancellation in the noncommutative Standard
Model \cite{AlvarezGraciaBondiaMartin1995}.

The reversal-odd standard module $W_-\cong\C^3$ gives the generation carrier;
the orthogonal sign-twisted standard module is retained as loop provenance and
is not counted as another generation copy.  Algebraic explanations of the
three-family structure have been pursued by very different mechanisms,
including division-algebraic models \cite{Furey2018}, centralizer
characterizations of Standard-Model symmetry \cite{Krasnov2021}, and a broad
recent survey of algebraic particle-model constraints \cite{FureyRoadmap2025}.
Very recent constructions include a Clifford-algebra realization with an
intrinsic $S_3$ family action \cite{Gresnigt2026} and an $E_7$ decomposition
into three Standard-Model generation modules \cite{Baez2026}.  The result here is different: the triplet is not selected
from an ambient exceptional algebra or inserted as three copies, but appears
as the reversal-odd endpoint source that remains after the complete spacetime
source has been shorted on one joint Gram.

The generation statement is therefore not a dimension-count coincidence.  A
copy of the same three-dimensional tetrahedral representation has already been
used by the spacetime branch.  The internal copy appears only \emph{after}
shorting the complete geometric source on one joint Gram.  What survives is a
distinct reversal-odd standard source.  Its rank is exactly three on the
positive branch and is protected by the same finite/cofinal lower-frame tests
used in the duality theorem.

Finally, we show that the duality is compatible with the common-action
construction of the spacetime paper.  If $S_{\rm g}$ and $S_{\rm i}$ are the
gravitational and internal variation syntheses on one history, with Grams
$G_{\rm g}$ and $G_{\rm i}$ and mixed block $C$, the genuinely new internal
variation is measured by the Schur quotient
\begin{equation}
\boxed{
G_{{\rm i}\mid{\rm g}}
=
G_{\rm i}-CG_{\rm g}^{-1}C^*.}
\label{eq:intro-schur}
\end{equation}
Its rank is the exact additional internal source cost after gravity has been
accounted for.  On this common carrier the usual finite gauge, Higgs,
Dirac--Yukawa, Ward, BRST, and stress identities are derivatives of one
action.  This finite compatibility is the only dynamical input needed here;
the analytic construction of a full Einstein--Standard-Model continuum is left
to a separate work.

The scope is deliberately structural.  The paper reconstructs the finite
internal algebra, physical quotient gauge group, anomaly-fixed hypercharge
ratios, a rank-three generation carrier, one primitive weak doublet, and a
single finite Dirac/Yukawa operator on the selected finite branch.  It
does not determine the numerical Yukawa matrices, CKM or PMNS data, gauge
couplings, Higgs parameters, renormalization-group running, matching scales,
or the determinant/Pfaffian phase of the quantum fermion measure.  Nor does the
finite commutant theorem alone prove an interacting four-dimensional quantum
continuum limit.

\paragraph{Logical dependence on the companion constructions.}
The results below separate the abstract finite mathematics from its
Renewal-Geometric realization.  The finite commutant duality, its reciprocal
Wedderburn, support-polar and quiver descriptions, the positive finite
certificates, and the structural Standard-Model classification are proved from
the explicit finite hypotheses stated in this paper.  Their proofs do not
require the reader to reconstruct the Renewal Grand Tensor, the ADM root-race
calibration, or the spectral Hodge calculus.  The companion works
\cite{PelissierSpacetime2026,PelissierSpectral2026} supply the independently
reconstructed Lorentzian and spectral ingredients.  The system-bearing mixed
incidence and provenance data required by the structural branch are treated by
finite occurrence tests.  In particular, when the grading, type and incidence
compressions are word-native operations of the historical flat algebra, their
mixed ancestry Grams are already finite inserted moments of the Grand Tensor;
no independent cross-Hankel measurement is then required.  The remaining
provenance question is whether those typing operations themselves belong to the
historical algebra.  We also show why coarse record/environment data, or an
external shadow not represented in that algebra, cannot determine the acted-on
system structure.  The final common-action subsection is only a
compatibility statement and does not attempt the independent analytic
construction of an Einstein--Standard-Model continuum.  Accordingly, the
finite results may equally be read as theorems about any joint packet satisfying
Definitions~\ref{def:joint-packet} and~\ref{def:action-algebra}.

The conceptual structure of the paper can be summarized as
\begin{equation}
\boxed{
\begin{aligned}
\text{predictive Renewal Geometry}
&\longrightarrow
\text{Lorentzian + spectral joint carrier}\\
&\longrightarrow
\text{typed action algebra}\\
&\overset{\rm commutant}{\longleftrightarrow}
\text{matter multiplicity}\\
&\longrightarrow
\text{structural Standard Model}.
\end{aligned}}
\label{eq:intro-chain}
\end{equation}

Section~\ref{sec:joint-interface} states the finite joint-carrier inputs
axiomatically, records how the two companion constructions realize the relevant
external and spectral ingredients, and proves the factor-commutant obstruction
that rules out the naive identification of the Standard-Model algebra with the
full Clifford commutant.
Section~\ref{sec:finite-duality} proves the exact finite duality in its
support-polar, multiplicity-quiver, and metric--holonomy forms.
Section~\ref{sec:structural-stability-action} reconstructs the structural
Standard-Model branch through a multiplicity-birth and router/bridge hierarchy,
adds source-native reciprocal-return and central-separation criteria, proves
generic exactness away from a Howe enhancement discriminant, identifies
source-minimal incidence skeletons, gives finite and cofinal positive
certificates for the duality, and records only the source-minimal common-action
compatibility needed for later gravity--matter dynamics.
Section~\ref{sec:discussion} discusses interpretation and limitations.
Additional representation-theoretic and stability details are collected in
the appendices.

% =============================================================================
\section{Axiomatic joint carrier and the commutant obstruction}
\label{sec:joint-interface}
% =============================================================================

\subsection{Abstract Lorentzian Clifford input and its Renewal realization}

The spacetime construction of Ref.~\cite{PelissierSpacetime2026} supplies one
Renewal-Geometric realization of the finite external data used here.  At
finite cutoff $X$, it first reconstructs the minimum persistent predictive
carrier.  On the minimal faithful spacetime branch, two logically independent
finite selectors identify the same tetrahedral relation cell $K_4$: exact
endpoint--cycle source-rank balance among connected simple regular relations,
and a stronger alternating/minimal-source selector within the
pair-homogeneous simplex category.  Both yield the spatial standard module
\[
W_4
=
\left\{x\in\R^4:\sum_{i=1}^{4}x_i=0\right\}
\cong\R^3.
\]
A calibrated $A_3$ root race then reconstructs the ADM metric, lapse, and
shift.

The sign and Clifford layers occur strictly later.  The positive predictive
word Grams and the root-race bracket cannot create a negative direction.  A
protected binary orientation supplies one independent signed line, and the
minimal nondegenerate extension of the positive three-dimensional spatial form
therefore has Lorentzian inertia $(1,3)$, up to overall sign, \emph{before} a
Clifford representation is chosen.  The represented orientation involution is
\[
J=P_+-P_-.
\]
A minimal complex Clifford realization may then be chosen with generators
$\gamma^\mu$ satisfying
\[
(\gamma^0)^2=-I,
\qquad
(\gamma^a)^2=I,
\qquad
\gamma^\mu\gamma^\nu+\gamma^\nu\gamma^\mu=0
\quad(\mu\ne\nu),
\]
and hence
\begin{equation}
\boxed{
\A_{{\rm ext},X}
:=
C^*(\gamma^0,\gamma^1,\gamma^2,\gamma^3)
\cong M_4(\C).}
\label{eq:external-clifford}
\end{equation}
The corresponding principal symbol squares to the already reconstructed
Lorentzian ADM quadratic form.  Thus the matrix factor entering the duality is
a faithful representation of a signed geometry selected upstream; it is not an
algebra introduced to manufacture the spacetime signature.

For the finite theorems of
Sections~\ref{sec:joint-interface}--\ref{sec:structural-stability-action}, the
only required spacetime input is the faithfully represented Clifford factor in
\eqref{eq:external-clifford}, together with the explicitly stated tetrahedral
source data used later in the structural branch.  Predictive reconstruction,
the two $K_4$ selectors, root-race inversion, ADM calibration, the one-sign
Lorentzian extension, and the Einstein handoff are therefore realization
results rather than hidden premises of the finite duality.

\subsection{Abstract spectral input and its Renewal realization}

The spectral companion paper \cite{PelissierSpectral2026} gives an independent
Renewal-Geometric realization of the same finite-carrier philosophy.  Starting
from the predictive Grand Tensor it reconstructs a finite history algebra and,
from a compatible first-order variation, a real even finite spectral triple.
None of that reconstruction is needed to prove the finite commutant theorems
below: Definition~\ref{def:joint-packet} states the represented algebraic data
axiomatically.

The spectral construction becomes relevant only for the cofinal stability
problem, where it provides one mechanism producing compact Hodge/current
screens.  Theorem~\ref{thm:cofinal-duality} is nevertheless stated purely in
terms of Mosco convergence, compact resolvent control, and protected kernel
projections; it is independent of how such screens are produced.

The joint problem is not solved by taking the direct sum of the two previously
constructed sectors.  A direct sum would erase precisely the mixed words that
record matter incidence, gauge transport, grading change, and common-source
orientation.  The new object is therefore a common represented packet.

\subsection{Source-complete joint packets}

\begin{definition}[Finite source-complete joint packet]
\label{def:joint-packet}
At cutoff $X$, a finite source-complete spacetime--gauge packet consists of:
\begin{enumerate}[label=\textup{(J\arabic*)},leftmargin=3em]
\item a finite Hilbert carrier $\Hh_X$ on which the reconstructed external
      Clifford algebra $\A_{{\rm ext},X}\cong M_4(\C)$ acts faithfully;
\item a commuting finite-dimensional typed internal $C^*$-algebra
      $\A_{{\rm int},X}$;
\item a central type decomposition
      \[
      \Hh_X
      \cong
      \bigoplus_{\lambda\in\Lambda_X}
      \C^4\otimes V_{\lambda,X}\otimes M_{\lambda,X},
      \]
      such that
      \[
      \A_{{\rm int},X}
      =
      \bigoplus_{\lambda}
      I_4\otimes\mathcal B(V_{\lambda,X})\otimes I_{M_{\lambda,X}};
      \]
\item a finite directed incidence graph
      $\Gamma_X=(\Lambda_X,E_X)$ and represented incidence histories
      \[
      Y_{e,X}:
      \C^4\otimes V_{s(e),X}\otimes M_{s(e),X}
      \longrightarrow
      \C^4\otimes V_{t(e),X}\otimes M_{t(e),X};
      \]
\item the mixed history words needed to determine products and adjoints of
      $\A_{{\rm ext},X}$, $\A_{{\rm int},X}$, and the $Y_{e,X}$ on one common
      carrier.
\end{enumerate}
All marks used in a claimed mixed relation are required to occur on one
future-minimal readable packet.  No inverse incidence map is part of the
definition.
\end{definition}

\begin{definition}[Joint typed action algebra]
\label{def:action-algebra}
For a joint packet define
\begin{equation}
\boxed{
\A_{{\rm act},X}
:=
C^*\!\left(
\A_{{\rm ext},X},
\A_{{\rm int},X},
Y_{e,X},Y_{e,X}^*:e\in E_X
\right).}
\label{eq:joint-action-algebra}
\end{equation}
The notation ``action'' refers here to represented action on the common finite
carrier; the variational action functional is introduced later.
\end{definition}

\begin{remark}[Why the joint packet is not a direct sum]
The mixed incidence words are not decorative.  They are exactly the data that
reduce the large commutant of the bare external factor to the physically
surviving multiplicity algebra.  Replacing the common carrier by a direct sum
of a spacetime packet and an internal packet would make
\eqref{eq:joint-action-algebra} block diagonal and would destroy the duality
proved below.
\end{remark}

\subsection{Source identifiability on the acted-on carrier}
\label{sec:source-identifiability}

The source-complete packet is a represented system object.  It cannot in
general be reconstructed from a classical outcome record or from the algebra
of an auxiliary environment alone.  The following finite statement separates
what such coarse data determine from what is needed for the duality.

\begin{countertheorem}[Classical record data do not determine the system algebra]
\label{thm:record-nonidentifiability}
There exist two finite normalized instruments on the same input and output
Hilbert spaces with the same resolved classical record channel at every word
length, the same one-step effects, the same one-letter Hilbert--Schmidt Gram,
and the same nonzero coarse Choi spectrum, but with nonisomorphic represented
system coefficient algebras and commutants.  Consequently no theorem using
only those record/environment invariants can identify the joint action algebra
of Definition~\ref{def:action-algebra}.
\end{countertheorem}

\begin{proof}
An explicit pair is constructed in Appendix~\ref{app:source-identifiability}.
One coefficient family is simultaneously diagonal and generates
$\C^{20}$; the other generates $M_4(\C)\otimes I_5$.  Each resolved branch is
a common scalar multiple of a unitary, so every uninterrupted classical word
has the same effect in the two realizations, while the coefficient algebras
and their commutants are different.  The displayed Gram and Choi identities in
Appendix~\ref{app:source-identifiability} complete the finite witness.
\end{proof}

Let $V:\Hh\to\Hh\otimes\mathcal E$ be an isometry and choose an orthonormal
environment basis $(|\nu\rangle)$.  Write
\[
 V=\sum_\nu C_\nu\otimes|\nu\rangle,
 \qquad
 C(z)=(I\otimes\langle z|)V,
\]
and let $\mathcal P_*$ be the complementary channel.  Its Heisenberg pullback
obeys
\begin{equation}
 \mathcal P^*(|z\rangle\langle w|)=C(z)^*C(w).
 \label{eq:complementary-pullback}
\end{equation}

\begin{theorem}[Identity-pivot reconstruction of the system action]
\label{thm:pullback-reconstruction}
Suppose the occurring coefficient algebra satisfies
$\mathcal A_{\rm sys}=C^*(C_\nu:\nu)=C^*(A_1,\ldots,A_s)$ and there are environment vectors
$z_0,z_1,\ldots,z_s$ with
\[
 C(z_0)=\alpha I,\qquad C(z_j)=\beta_jA_j,
 \qquad \alpha\beta_j\ne0.
\]
Then
\begin{equation}
 \boxed{
 C^*\!\left(\mathcal P^*(\mathcal B(\mathcal E))\right)
 =\mathcal A_{\rm sys},
 \qquad
 [\mathcal P^*(\mathcal B(\mathcal E))]'=\mathcal A_{\rm sys}'.}
 \label{eq:pullback-reconstruction}
\end{equation}
In particular, a complete operator-valued complementary pullback recovers the
represented system action whenever the physical source contains such an
identity pivot; no fitted remainder algebra is required.
\end{theorem}

\begin{proof}
Every pullback $C(z)^*C(w)$ belongs to $C^*(C_\nu)$ and hence to
$\mathcal A_{\rm sys}$.  Conversely,
\[
 \mathcal P^*(|z_0\rangle\langle z_j|)
 =\overline\alpha\,\beta_j A_j,
\]
so every generator $A_j$ lies in the pullback algebra.  Equality of the
commutants follows immediately.
\end{proof}

\begin{corollary}[Finite coherent system export]
\label{cor:finite-system-export}
For the normalized six-edge source class of
Appendix~\ref{app:source-identifiability}, ten fixed Hermitian environment
quadratures suffice to reconstruct all six edge coefficients on the acted-on
carrier.  With Hilbert--Schmidt coefficient error $e_K$ as defined in
\eqref{eq:appendix-reconstruction-error}, the associated star-closed
commutator map changes in operator norm by at most $2\sqrt2\,e_K$.  Hence an
independently verified protected commutant complement with measured least
singular value $\widehat s>2\sqrt2\,e_K$ has no additional commutant direction.
\end{corollary}

\begin{remark}[Logical role of source identification]
Theorem~\ref{thm:record-nonidentifiability} is a stopping theorem, not a
weakening of the double-centralizer result.  Once the represented system
operators are supplied or reconstructed, Theorem~\ref{thm:main-duality} is
exact.  What cannot be inferred from coarse record data is which system
operators occurred on the common carrier in the first place.
\end{remark}

\subsection{Historical operation recognition and word-native source compilation}
\label{sec:historical-operation-recognition}

The preceding reconstruction problem has a second layer.  Even after a finite
history algebra has been reconstructed, an invariant endomorphism of a word
Hilbert space need not be an operation induced by the historical system
algebra.  The following module criterion separates these notions.

Let $\A\subseteq\mathcal B(\Hh)$ be a finite unital $C^*$-algebra generated by
the represented historical operations on one exposed carrier, and let a finite
group $G$ act by unitaries $u_g$ with $u_g\A u_g^*=\A$.  Write
$\alpha_g(a)=u_gau_g^*$ and
\[
 \A\cong\bigoplus_\beta M_{n_\beta}(\C),\qquad
 D=\dim_\C\A=\sum_\beta n_\beta^2,
\]
with normalized regular trace
\[
 \tau_{\rm reg}(a)=D^{-1}\sum_\beta n_\beta\Tr(a_\beta).
\]
On $\mathscr W=L^2(\A,\tau_{\rm reg})$ let
$L_ax=ax$, $R_bx=xb$ and $U_gx=\alpha_g(x)$.

\begin{theorem}[Word-module recognition of historical internal actions]
\label{thm:word-module-recognition}
The operations on the regular word carrier induced by invariant historical
system elements are exactly
\begin{equation}
 \boxed{
 R(\A)'\cap U(G)'=L(\A^G),
 \qquad
 \A^G=\A\cap\{u_g:g\in G\}'.}
 \label{eq:word-module-recognition}
\end{equation}
Equivalently, $T\in\End(\mathscr W)$ is left multiplication by a unique
invariant historical element if and only if it commutes with every right
multiplication and every $U_g$; the coefficient is $T1$.
\end{theorem}

\begin{proof}
Every $L_a$ commutes with every $R_b$.  Conversely, if $[T,R_b]=0$ for all
$b$, then $T(b)=T(R_b1)=R_bT1=(T1)b$, hence $T=L_{T1}$.  Moreover
$U_gL_aU_g^{-1}=L_{\alpha_g(a)}$, so faithfulness of left multiplication gives
$[U_g,L_a]=0$ exactly when $a\in\A^G$.
\end{proof}

\begin{proposition}[Exact module and symmetry defects]
\label{prop:word-module-defect}
Let $(e_\nu)_{\nu=1}^D$ be a $\tau_{\rm reg}$-orthonormal basis of $\A$ and
put $a=T1$.  Define
\[
 \delta_R(T)^2=\sum_{\nu=1}^D\norm{[T,R_{e_\nu}]1}_{\tau}^2,
 \qquad
 \delta_G(T)^2=\frac1{|G|}\sum_g\norm{[U_g,T]1}_{\tau}^2.
\]
If $E_G(a)=|G|^{-1}\sum_g\alpha_g(a)$, then
\begin{align}
 \norm{T-L_a}_{\HS(\mathscr W)}^2&=\delta_R(T)^2,
 \label{eq:word-module-defect}\\
 \norm{T-L_{E_G(a)}}_{\HS(\mathscr W)}
 &\le \delta_R(T)+\sqrt{D/2}\,\delta_G(T).
 \label{eq:word-module-invariant-defect}
\end{align}
Thus both defects vanish exactly on the algebra in
\eqref{eq:word-module-recognition}.
\end{proposition}

\begin{proof}
For each basis element,
$(T-L_a)e_\nu=T(R_{e_\nu}1)-R_{e_\nu}T1=[T,R_{e_\nu}]1$; summation gives
\eqref{eq:word-module-defect}.  Since $U_g1=1$,
$[U_g,T]1=\alpha_g(a)-a$.  Orthogonal group averaging gives
$|G|^{-1}\sum_g\norm{\alpha_g(a)-a}_\tau^2=2\norm{a-E_G(a)}_\tau^2$.
Finally $\norm{L_b}_{\HS(\mathscr W)}^2=D\norm b_\tau^2$, and the triangle
inequality proves \eqref{eq:word-module-invariant-defect}.
\end{proof}

\begin{corollary}[Word-space multiplicity firewall]
\label{cor:word-multiplicity-firewall}
Multiplicity of the conjugation representation on $L^2(\A,\tau_{\rm reg})$
need not equal multiplicity on the exposed physical carrier and cannot by
itself be interpreted as an executable internal factor.  Any such inference
must also pass Theorem~\ref{thm:word-module-recognition}.  In particular, the
multiplicity-birth analysis of Section~\ref{sec:multiplicity-birth} always
refers to an actual reachable--observable physical carrier, not merely to the
regular word representation.  Appendix~\ref{app:historical-source-selection}
contains an explicit twelve-root example with an invariant $M_7(\C)$ word-space
block but no physical internal $M_3(\C)$ factor.
\end{corollary}

We next identify a class of mixed source quantities that require no independent
measurement once the relevant typing operations are already historical.

\begin{definition}[Word-native linear shadow]
\label{def:word-native-shadow}
A linear map $\mathscr L:\A\to\A$ is \emph{word native} if
\begin{equation}
 \boxed{\mathscr L(X)=\sum_{\alpha=1}^m c_\alpha A_\alpha X B_\alpha,}
 \label{eq:word-native-shadow}
\end{equation}
where $A_\alpha,B_\alpha\in\A$ are reconstructed historical-algebra elements
and the coefficients $c_\alpha$ are fixed by the declared shadow before the
target conclusion is evaluated.  Finite compressions, conjugations, finite
group averages and finite functional-calculus operations are included whenever
their defining elements belong to $\A$.
\end{definition}

\begin{theorem}[Inserted word-moment compiler]
\label{thm:inserted-word-moments}
Let $x,y,z\in\A$ be represented words and let $\mathscr L,\mathscr M$ be
word-native maps, with
$\mathscr M(Z)=\sum_\beta d_\beta C_\beta ZD_\beta$.  Then
\begin{align}
 \langle y,\mathscr L(x)\rangle_{\HS}
 &=\sum_\alpha c_\alpha\,
   \tau_{\rm reg}(y^*A_\alpha xB_\alpha),
 \label{eq:inserted-one-moment}\\
 \langle\mathscr L(x),\mathscr M(z)\rangle_{\HS}
 &=\sum_{\alpha,\beta}\overline{c_\alpha}d_\beta\,
   \tau_{\rm reg}(B_\alpha^*x^*A_\alpha^*C_\beta zD_\beta).
 \label{eq:inserted-two-moments}
\end{align}
Hence every ordinary--shadow and shadow--shadow Gram block is a finite linear
combination of ordinary Grand-Tensor word moments.  If the historical word
algebra is already flat, the reconstructed multiplication table reduces all
inserted products and no deeper physical histories are required.
\end{theorem}

\begin{proof}
The formulas are the Hilbert--Schmidt inner product followed by substitution of
\eqref{eq:word-native-shadow}.  Cyclicity of $\tau_{\rm reg}$ rewrites each
summand as the inner product of two represented algebra elements.  At flatness
the finite word span is a $*$-algebra, so every inserted product reduces in its
known multiplication table.
\end{proof}

\begin{corollary}[Table-native ancestry short]
\label{cor:table-native-ancestry}
Let $W_r$ synthesize a flat historical word bank, let
$S=W_rB_S$, $T=\mathscr L W_rB_T$, and define
\[
 K^{00}=W_r^*W_r,\qquad
 K^{0L}=W_r^*\mathscr LW_r,\qquad
 K^{LL}=W_r^*\mathscr L^*\mathscr LW_r.
\]
Then
\begin{equation}
 \boxed{
 G=B_S^*K^{00}B_S,
 \quad C=B_S^*K^{0L}B_T,
 \quad H=B_T^*K^{LL}B_T,}
 \label{eq:table-native-GCH}
\end{equation}
and the ancestry short
\begin{equation}
 \boxed{R_{T\mid S}=H-C^*G^\dagger C}
 \label{eq:table-native-short}
\end{equation}
is a deterministic function of the historical word Gram, multiplication table
and the fixed word-native compiler.  No independently executed mixed
cross-Hankel panel is required on this branch.
\end{corollary}

\begin{countertheorem}[Word moments cannot manufacture an external shadow]
\label{thm:external-shadow-no-go}
There are conservative extensions with identical historical word algebra,
multiplication table, complete word-Gram hierarchy and resolved classical
cylinder probabilities but different values of a proposed grading-changing or
typed-incidence cross block whose defining shadow is not represented in the
historical algebra.  Therefore no function of the old word moments alone can
recover such an external shadow.
\end{countertheorem}

\begin{proof}
Tensor a faithful historical realization with an unread two-level spectator on
which every old operation acts trivially.  One extension contains no coherent
cross-support link; another admits a future-visible spectator link and its
typed compression.  All old word products and moments agree, whereas the
proposed external cross block is zero in the first extension and nonzero in the
second.
\end{proof}


% =============================================================================
\subsection{Why the naive spacetime--gauge commutant fails}
\label{sec:obstruction}
% =============================================================================

\subsubsection{The factor obstruction}

The simplest possible duality would identify the internal gauge algebra with
the full commutant of the external Clifford factor.  Finite-dimensional
representation theory rules this out.

\begin{theorem}[Factor-commutant obstruction]
\label{thm:factor-obstruction}
Let a unital copy of $M_n(\C)$ act on a finite-dimensional Hilbert space
$\Hh$.  Then for some finite multiplicity space $\mathcal K$,
\begin{equation}
\Hh\cong\C^n\otimes\mathcal K,
\qquad
M_n(\C)=M_n(\C)\otimes I_{\mathcal K},
\qquad
M_n(\C)'=I_n\otimes\mathcal B(\mathcal K).
\label{eq:factor-commutant}
\end{equation}
Consequently a reducible internal algebra with nontrivial centre cannot in
general equal the full commutant of the external Clifford factor.
\end{theorem}

\begin{proof}
Every finite-dimensional unital representation of the simple algebra
$M_n(\C)$ is a finite direct sum of copies of its defining representation.
Combining the copies gives $\Hh\cong\C^n\otimes\mathcal K$.  An operator
commutes with every matrix unit on the first factor if and only if it is the
identity there and arbitrary on the multiplicity factor.
\end{proof}

For the spacetime branch $n=4$.  Thus the raw Clifford commutant is always a
full matrix algebra on the unresolved multiplicity space.  The Standard Model
cannot appear merely by writing
$\A_{\rm SM}=\A_{\rm ext}'$.  Incidence data are essential.

\subsubsection{External, internal, and multiplicity factors}

\begin{theorem}[External--internal--multiplicity decomposition]
\label{thm:trinity}
Let $\A_{\rm ext}\cong M_n(\C)$ and let $\A_{\rm int}$ be a commuting
finite-dimensional unital $C^*$-algebra on $\Hh$.  Then there are finite
spaces $V_\lambda$ and $M_\lambda$ such that
\begin{equation}
\boxed{
\Hh\cong
\bigoplus_{\lambda\in\Lambda}
\C^n\otimes V_\lambda\otimes M_\lambda,}
\label{eq:trinity-H}
\end{equation}
\begin{equation}
\boxed{
\A_{\rm int}
=
\bigoplus_\lambda
I_n\otimes\mathcal B(V_\lambda)\otimes I_{M_\lambda}.}
\label{eq:trinity-Aint}
\end{equation}
If
$\mathcal C=C^*(\A_{\rm ext},\A_{\rm int})$, then
\begin{align}
\mathcal C
&=
\bigoplus_\lambda
\mathcal B(\C^n\otimes V_\lambda)\otimes I_{M_\lambda},
\label{eq:trinity-C}\\
\mathcal C'
&=
\bigoplus_\lambda
I_{\C^n\otimes V_\lambda}\otimes\mathcal B(M_\lambda).
\label{eq:trinity-Cprime}
\end{align}
In particular, the residual spaces $M_\lambda$, rather than the raw internal
representation spaces $V_\lambda$, carry flavour and generation multiplicity.
\end{theorem}

\begin{proof}
Apply Theorem~\ref{thm:factor-obstruction} to write
$\Hh=\C^n\otimes\mathcal K$ and
$\A_{\rm int}=I_n\otimes\widetilde\A_{\rm int}$.
The finite-dimensional representation theorem gives
\[
\mathcal K
\cong
\bigoplus_\lambda V_\lambda\otimes M_\lambda,
\qquad
\widetilde\A_{\rm int}
=
\bigoplus_\lambda
\mathcal B(V_\lambda)\otimes I_{M_\lambda}.
\]
On each minimal central block the external and internal full matrix factors
generate $\mathcal B(\C^n\otimes V_\lambda)$.  Schur's lemma then gives the
commutant \eqref{eq:trinity-Cprime}.
\end{proof}

\begin{interpretation}
The theorem identifies three logically distinct structures.  The external
factor carries the Lorentzian Clifford action.  The spaces $V_\lambda$ carry
the typed internal gauge representations.  The spaces $M_\lambda$ are what the
combined external--internal action cannot see.  Matter incidence can reduce
this residual freedom, and the surviving reduction is precisely the dual
algebra reconstructed in the next section.
\end{interpretation}

% =============================================================================
\section{Exact finite spacetime--gauge duality}
\label{sec:finite-duality}
% =============================================================================

\subsection{Typed incidence and the residual multiplicity algebra}

\begin{definition}[Typed multiplicity algebra]
\label{def:typed-multiplicity}
Let $\Gamma=(\Lambda,E)$ be the incidence graph of a joint packet.  Suppose
the typed carrier-intertwiner space for each edge $e:\lambda\to\mu$ is one
dimensional and choose a nonzero typed intertwiner $D_e$.  Write
\begin{equation}
Y_e=D_e\otimes F_e,
\qquad
F_e:M_\lambda\to M_\mu.
\label{eq:incidence-factorization}
\end{equation}
The \emph{typed multiplicity algebra} is
\begin{equation}
\boxed{
\M_{\rm type}
=
\left\{
(R_\lambda)_\lambda:
R_\mu F_e=F_eR_\lambda,\ 
R_\lambda F_e^*=F_e^*R_\mu
\ \forall e:\lambda\to\mu
\right\}.}
\label{eq:typed-multiplicity}
\end{equation}
It acts on the full carrier as
$\bigoplus_\lambda I_{\C^4\otimes V_\lambda}\otimes R_\lambda$.
\end{definition}

\begin{lemma}[Multiplicity-one incidence factorization]
\label{lem:multiplicity-factorization}
If the external--internal symmetry acts trivially on the multiplicity spaces
and
\[
\dim_\C
\Hom_{\rm typed}
(\C^4\otimes V_\lambda,\C^4\otimes V_\mu)=1,
\]
then every equivariant incidence map has a unique factorization
\eqref{eq:incidence-factorization}.
\end{lemma}

\begin{proof}
The equivariant hom-space splits as
\[
\Hom_{\rm typed}
((\C^4\otimes V_\lambda)\otimes M_\lambda,
(\C^4\otimes V_\mu)\otimes M_\mu)
\cong
\Hom_{\rm typed}(\C^4\otimes V_\lambda,\C^4\otimes V_\mu)
\otimes
\Hom(M_\lambda,M_\mu).
\]
The first factor is the line $\C D_e$.
\end{proof}

\begin{theorem}[Exact edge-commutant reconstruction]
\label{thm:edge-commutant}
For
\[
\A_{\rm act}
=
C^*(\A_{\rm ext},\A_{\rm int},Y_e,Y_e^*:e\in E),
\]
one has
\begin{equation}
\boxed{\A_{\rm act}'=\M_{\rm type}.}
\label{eq:edge-commutant}
\end{equation}
\end{theorem}

\begin{proof}
By Theorem~\ref{thm:trinity}, an operator commuting with
$\A_{\rm ext}$ and $\A_{\rm int}$ has the form
\[
X=
\bigoplus_\lambda
I_{\C^4\otimes V_\lambda}\otimes R_\lambda.
\]
For $Y_e=D_e\otimes F_e$,
\[
XY_e-Y_eX
=
D_e\otimes(R_\mu F_e-F_eR_\lambda),
\]
and the commutator with $Y_e^*$ gives the adjoint intertwining equation.
Because $D_e\ne0$, the full commutators vanish exactly when
\eqref{eq:typed-multiplicity} holds.
\end{proof}

The theorem already gives the finite duality in one direction.  The reverse
direction is automatic by finite bicommutant closure, but the raw edge formula
does not yet reveal how singular residues behave.  We therefore give two
intrinsic presentations.

\subsection{Support-polar form: singular residues are allowed}

Let $F:E\to H$ have polar decomposition
\[
F=UP,
\qquad
P=(F^*F)^{1/2},
\qquad
p=U^*U=\supp P,
\qquad
q=UU^*=\supp(FF^*).
\]

\begin{lemma}[Singular polar-edge equivalence]
\label{lem:polar-edge}
For $R_E\in\mathcal B(E)$ and $R_H\in\mathcal B(H)$, the equations
\[
R_HF=FR_E,
\qquad
R_EF^*=F^*R_H
\]
are equivalent to
\begin{equation}
\boxed{
[R_E,P^2]=0,
\qquad
R_HU=UR_E,
\qquad
R_EU^*=U^*R_H.}
\label{eq:polar-edge-equivalence}
\end{equation}
No invertibility assumption is required.
\end{lemma}

\begin{proof}
The two incidence equations imply
$F^*R_HF=F^*FR_E=R_EF^*F$, hence $[R_E,P^2]=0$.
Functional calculus gives commutation with $P$ and its support.
Multiplication by the Moore--Penrose inverse of $P$ on its support gives the
partial-isometry transport equations.  The converse follows by multiplying
the transport equations by $P$ and using $[R_E,P]=0$.
\end{proof}

On the total multiplicity carrier
$\mathscr M=\bigoplus_\lambda M_\lambda$, let $Z_\lambda$ denote the vertex
projections and write $F_e=U_eP_e$.

\begin{definition}[Support-polar action algebra]
\label{def:support-polar}
Define
\begin{equation}
\boxed{
\Oalg_F^{\rm sup}
=
C^*(Z_\lambda,F_e^*F_e,U_e,U_e^*:\lambda,e)
\subseteq\mathcal B(\mathscr M).}
\label{eq:support-polar}
\end{equation}
\end{definition}

\begin{theorem}[Exact support-polar duality]
\label{thm:support-polar}
For arbitrary rectangular, singular, rank-deficient, or zero incidence maps,
\begin{equation}
\boxed{
\M_{\rm type}
=
(\Oalg_F^{\rm sup})',
\qquad
\Oalg_F^{\rm sup}
=
\M_{\rm type}'.}
\label{eq:support-polar-duality}
\end{equation}
Moreover,
\begin{equation}
\boxed{
C^*(Z_\lambda,F_e,F_e^*:\lambda,e)
=
\Oalg_F^{\rm sup}.}
\label{eq:incidence-support-polar-equality}
\end{equation}
Thus a zero singular value remains a genuine support boundary and is never
replaced by an inverse residue.
\end{theorem}

\begin{proof}
Commutation with the vertex projections makes an operator block diagonal.
Lemma~\ref{lem:polar-edge} shows that commutation with
$F_e^*F_e,U_e,U_e^*$ is exactly the pair of incidence intertwining equations
in Definition~\ref{def:typed-multiplicity}.  This gives the first equality,
and finite bicommutant closure gives the second.

For \eqref{eq:incidence-support-polar-equality}, $F_e^*F_e$ belongs to the edge
algebra.  Since its spectrum is finite, interpolate the function
$g(0)=0$, $g(t)=t^{-1/2}$ for $t>0$ on that spectrum by a polynomial.
Then
\[
U_e=F_eg(F_e^*F_e)
\]
belongs to the incidence algebra.  Conversely
$F_e=U_e(F_e^*F_e)^{1/2}$.
\end{proof}

\subsection{Noncommuting supports and the multiplicity quiver}

Support projections themselves can fail to commute.  The correct residual
description is then a finite operator-space quiver, in the standard
representation-theoretic sense of a quiver with linear arrow spaces
\cite{Schiffler2014}.

\begin{definition}[Multiplicity quiver]
\label{def:multiplicity-quiver}
Let
\[
\mathcal S_F
=
C^*(Z_\lambda,\supp(F_e^*F_e),\supp(F_eF_e^*):\lambda,e)
\]
and write its canonical finite isotypic decomposition as
\[
\mathscr M
\cong
\bigoplus_{a\in\mathfrak A}V_a\otimes N_a,
\qquad
\mathcal S_F
=
\bigoplus_a\mathcal B(V_a)\otimes I_{N_a}.
\]
Let $\bm T$ be any finite $*$-closed family of further reconstructed support
metrics, polar arrows, cycle products, or action generators.  For a block
$T_{ba}:V_a\otimes N_a\to V_b\otimes N_b$, define
\[
\mathfrak B_{ba}(T)
=
\{(\varphi\otimes\id)(T_{ba}):
\varphi\in\mathcal B(V_a,V_b)^*\}
\subseteq\mathcal B(N_a,N_b).
\]
The multiplicity quiver $\mathfrak Q_F(\bm T)$ has vertex spaces $N_a$ and
arrow spaces
\[
\mathfrak B_{ba}
=
\Span_{T\in\bm T}\mathfrak B_{ba}(T).
\]
Its endomorphism algebra is
\[
\End\mathfrak Q_F
=
\left\{(X_a)_a:
X_bB=BX_a
\ \text{for every }B\in\mathfrak B_{ba}\right\}.
\]
\end{definition}

\begin{theorem}[Noncommuting-support quiver theorem]
\label{thm:quiver}
With the preceding notation,
\begin{equation}
\boxed{
C^*(\mathcal S_F,\bm T)'
\cong
\End\mathfrak Q_F(\bm T).}
\label{eq:quiver-commutant}
\end{equation}
If
\[
(T_j)_{ba}
=
\sum_\alpha E_\alpha^{ba}\otimes B_{j,\alpha}^{ba}
\]
in Hilbert--Schmidt orthonormal bases of
$\mathcal B(V_a,V_b)$, then the ambient commutator form restricts to
\begin{equation}
\boxed{
q_{\bm T}[(X_a)]
=
\sum_{j,b,a,\alpha}
\norm{B_{j,\alpha}^{ba}X_a-X_bB_{j,\alpha}^{ba}}_{\HS}^2,}
\label{eq:quiver-form}
\end{equation}
whose kernel is exactly $\End\mathfrak Q_F$.
\end{theorem}

\begin{proof}
An operator commuting with $\mathcal S_F$ has the form
$\bigoplus_a I_{V_a}\otimes X_a$.  Expanding its commutator with every
$T_j$ in the Hilbert--Schmidt basis gives precisely the quiver intertwining
relations.  Orthogonality gives \eqref{eq:quiver-form}.
\end{proof}

\subsection{The connected invertible branch: polar metric and holonomy}

When every residue map is invertible, the duality has a particularly geometric
form.

\begin{theorem}[Polar metric--holonomy commutant]
\label{thm:polar-holonomy}
Assume $\Gamma$ is connected, all multiplicity spaces have common dimension
$g$, and every edge map is invertible:
\[
F_e=U_eP_e,
\qquad
U_e\in\U(g),
\qquad
P_e\succ0.
\]
Choose a root $o$ and a spanning tree.  Let
$Q_\lambda:M_o\to M_\lambda$ be the ordered product of the $U_e$ along the
tree path.  For $e:\lambda\to\mu$, define on the root fibre
\[
K_e=Q_\lambda^*P_e^2Q_\lambda,
\qquad
W_e=Q_\mu^*U_eQ_\lambda,
\]
and put
\begin{equation}
\boxed{
\Oalg_F
=
C^*(K_e,W_e:e\in E)\subseteq M_g(\C).}
\label{eq:root-holonomy-algebra}
\end{equation}
Then
\begin{equation}
\boxed{
\M_{\rm type}^{(o)}
=
\Oalg_F',
\qquad
\Oalg_F
=
(\M_{\rm type}^{(o)})'.}
\label{eq:root-holonomy-duality}
\end{equation}
The multiplicity family associated with
$R\in\Oalg_F'$ is
\[
R_\lambda=Q_\lambda RQ_\lambda^*.
\]
\end{theorem}

\begin{proof}
The polar-edge equations give
$[R_\lambda,P_e^2]=0$ and
$R_\mu=U_eR_\lambda U_e^*$.
Along the spanning tree these force
$R_\lambda=Q_\lambda RQ_\lambda^*$ from the root value $R=R_o$.
The metric condition becomes $[R,K_e]=0$, while a non-tree edge gives the
cycle-consistency condition $[R,W_e]=0$.  Thus the residual family is
equivalent to one root operator commuting with the algebra
\eqref{eq:root-holonomy-algebra}.  Finite bicommutant closure gives the reverse
equality.
\end{proof}

\begin{corollary}[Flat and irreducible extremes]
\label{cor:polar-extremes}
Under Theorem~\ref{thm:polar-holonomy}:
\begin{enumerate}[label=\textup{(\roman*)},leftmargin=2.8em]
\item $\M_{\rm type}^{(o)}=M_g(\C)$ if and only if every $K_e$ and $W_e$ is
      scalar;
\item $\M_{\rm type}^{(o)}=\C I$ if and only if $\Oalg_F=M_g(\C)$;
\item if every $F_e$ is unitary, then $K_e=I$ and the obstruction reduces to
      relative cycle holonomy;
\item even on a tree, a non-scalar positive polar metric can reduce the
      residual multiplicity algebra.
\end{enumerate}
\end{corollary}

\subsection{The main finite duality theorem}

\begin{theorem}[Joint action--multiplicity double-centralizer theorem]
\label{thm:main-duality}
For every finite source-complete joint packet of
Definition~\ref{def:joint-packet},
\begin{equation}
\boxed{
\A_{{\rm act},X}'
=
\M_{{\rm type},X},
\qquad
\M_{{\rm type},X}'
=
\A_{{\rm act},X}.}
\label{eq:main-duality}
\end{equation}
On the total multiplicity carrier the same dual pair has the support-exact
presentation
\[
\M_{{\rm type},X}
=
(\Oalg_{F,X}^{\rm sup})',
\qquad
\Oalg_{F,X}^{\rm sup}
=
\M_{{\rm type},X}'.
\]
If the support projections do not commute, the residual algebra is
canonically represented as
\[
\M_{{\rm type},X}
\cong
\End\mathfrak Q_{F,X}.
\]
On a connected invertible branch, choosing one root multiplicity fibre gives
\[
\M_{{\rm type},X}^{(o)}
=
\Oalg_{F,X}',
\qquad
\Oalg_{F,X}
=
(\M_{{\rm type},X}^{(o)})'.
\]
Hence the exact duality is
\begin{equation}
\boxed{
\text{joint typed spacetime--gauge--incidence action}
\quad\longleftrightarrow\quad
\text{residual gauge--matter multiplicity}.}
\label{eq:dual-pair-words}
\end{equation}
\end{theorem}

\begin{proof}
The first commutant equality is Theorem~\ref{thm:edge-commutant}; the second is
finite bicommutant closure.  The remaining presentations are
Theorems~\ref{thm:support-polar}, \ref{thm:quiver}, and
\ref{thm:polar-holonomy}.
\end{proof}

\begin{corollary}[Reciprocal Wedderburn structure and common centre]
\label{cor:reciprocal-wedderburn}
For every finite joint packet satisfying Theorem~\ref{thm:main-duality}, there
are a finite index set $\mathfrak I_X$ and finite Hilbert spaces
$H^{\rm act}_{\alpha,X}$ and $H^{\rm mult}_{\alpha,X}$ such that, after one
unitary decomposition of the represented carrier,
\begin{equation}
\boxed{
\Hh_X
\cong
\bigoplus_{\alpha\in\mathfrak I_X}
H^{\rm act}_{\alpha,X}\otimes H^{\rm mult}_{\alpha,X},}
\label{eq:wedderburn-H}
\end{equation}
with
\begin{align}
\A_{{\rm act},X}
&\cong
\bigoplus_{\alpha\in\mathfrak I_X}
\mathcal B(H^{\rm act}_{\alpha,X})\otimes I_{H^{\rm mult}_{\alpha,X}},
\label{eq:wedderburn-action}\\
\M_{{\rm type},X}
&\cong
\bigoplus_{\alpha\in\mathfrak I_X}
I_{H^{\rm act}_{\alpha,X}}\otimes
\mathcal B(H^{\rm mult}_{\alpha,X}).
\label{eq:wedderburn-multiplicity}
\end{align}
In particular,
\begin{equation}
\boxed{
Z(\A_{{\rm act},X})
=
Z(\M_{{\rm type},X})
=
\A_{{\rm act},X}\cap\M_{{\rm type},X}.}
\label{eq:common-centre}
\end{equation}
The common minimal central projections therefore label the same represented
sectors on both sides, while the two matrix sizes are reciprocal action and
multiplicity dimensions inside each sector.  Any unitary equivalence of
source-complete joint packets transports the two algebras and their common
centre simultaneously.
\end{corollary}

\begin{proof}
Apply the finite-dimensional structure theorem to the unital $C^*$-algebra
$\A_{{\rm act},X}$.  In a unitary decomposition of its representation it has
the form
$\bigoplus_\alpha\mathcal B(H^{\rm act}_{\alpha,X})\otimes I$.
Its commutant is therefore exactly the algebra in
\eqref{eq:wedderburn-multiplicity}, which equals $\M_{{\rm type},X}$ by
Theorem~\ref{thm:main-duality}.  The centre of either algebra is generated by
the same minimal central block projections, proving
\eqref{eq:common-centre}.  Simultaneous unitary covariance is immediate from
the defining commutant identities.
\end{proof}


\subsection{Incidence essentiality and source-minimal skeletons}
\label{sec:incidence-skeleton}

Fix the non-incidence action
$\mathcal A_0=C^*(\mathcal A_{\rm ext},\mathcal A_{\rm int})$ and a finite
incidence bank $E$.  For $S\subseteq E$ set
\[
 \M(S)=C^*(\mathcal A_0,Y_e,Y_e^*:e\in S)',
 \qquad \M=\M(E),
\]
and define the residual test space
$V=\mathcal A_0'\cap\M^\perp$, with $n=\dim_\C V$.  For $X\in V$ let
\begin{equation}
 D_eX=2^{-1/2}([Y_e,X],[Y_e^*,X]),
 \qquad G_e=D_e^*D_e,
 \qquad G_S=\sum_{e\in S}G_e.
 \label{eq:incidence-forms}
\end{equation}
The full bank is exact precisely when $G_E\succ0$ on $V$.

\begin{theorem}[Incidence rank polymatroid]
\label{thm:incidence-polymatroid}
Let $R_e=\Ran D_e^*\subseteq V$ and
\begin{equation}
 f(S)=\dim_\C\!\left(\sum_{e\in S}R_e\right).
 \label{eq:incidence-rank}
\end{equation}
Then
\begin{equation}
 \boxed{
 f(S)=\rank G_S
     =n-\dim_\C(\M(S)\cap V),
 \qquad
 \M(S)=\M\oplus\Ker G_S}
 \label{eq:incidence-kernel}
\end{equation}
as orthogonal vector spaces.  The function $f$ is normalized, monotone and
submodular:
\begin{equation}
 f(S)+f(T)\ge f(S\cup T)+f(S\cap T).
 \label{eq:incidence-submodular}
\end{equation}
An incidence $e\in E$ is essential for the full duality exactly when
\begin{equation}
 \boxed{f(E\setminus\{e\})<n}
 \label{eq:essential-incidence}
\end{equation}
or, equivalently, when deleting $e$ strictly enlarges the commutant.
\end{theorem}

\begin{proof}
The quadratic form of $G_S$ is
$\sum_{e\in S}\|D_eX\|^2$, hence
\[
 \Ker G_S=\bigcap_{e\in S}\Ker D_e
 =\left(\sum_{e\in S}\Ran D_e^*\right)^\perp.
\]
Every element of $\M(S)$ already lies in $\mathcal A_0'$.  Orthogonal
projection onto $\M$ therefore leaves a residual vector in $V$, and the
incidence constraints on that residual are exactly $G_SX=0$.  This proves
\eqref{eq:incidence-kernel}.  If $R(S)=\sum_{e\in S}R_e$, then
$R(S)+R(T)=R(S\cup T)$ and
$R(S\cap T)\subseteq R(S)\cap R(T)$.  The dimension identity for two
subspaces yields submodularity.  The deletion criterion is the kernel formula
with $S=E\setminus\{e\}$.
\end{proof}

\begin{corollary}[Finite source-minimal incidence skeleton]
\label{cor:minimal-incidence-skeleton}
Every finite exact incidence bank contains an inclusion-minimal subset
$S\subseteq E$ with $\M(S)=\M(E)$.  Such a subset is obtained by repeatedly
deleting nonessential incidences and has at most $n$ members.  In general
these minimal subsets need not be unique and need not have equal cardinality;
the correct grouped rank object is the polymatroid $f$, not a matroid on the
incidence labels.
\end{corollary}

\begin{proof}
Finite reverse deletion terminates.  In an inclusion-minimal sufficient bank,
adding its members one by one strictly increases the span of the corresponding
constraint ranges, so at most $n$ additions are possible.  The remaining
claims follow from the grouped nature of the subspaces $R_e$; an explicit
unequal-cardinality example is given in Appendix~\ref{app:incidence-details}.
\end{proof}

\begin{interpretation}[Meaning of spacetime--gauge Howe-type duality]
The phrase ``Howe-type duality'' is used below for the mutual-centralizer
content of Theorem~\ref{thm:main-duality} together with the reciprocal block
structure of Corollary~\ref{cor:reciprocal-wedderburn}.  The statement does not
exchange a bare spacetime algebra with a bare gauge algebra.  It says that,
once spacetime, typed gauge action, and physical matter incidence are retained
on one represented carrier, the complete visible action and the invisible
matter multiplicity determine one another exactly.
\end{interpretation}

\begin{remark}[Relation to earlier commutant dualities]
Classical Howe duality concerns mutually centralizing actions with strong
multiplicity correspondence properties \cite{Howe1989}.  Gauge/observable
Howe pairs also occur in four-dimensional conformal field theory
\cite{Todorov2010}.  In noncommutative Standard-Model geometry, the Hodge or
Morita condition instead compares the spectral Clifford algebra with its
opposite/commutant \cite{DAndreaDabrowski2016,DabrowskiDAndreaSitarz2018}.
Theorem~\ref{thm:main-duality} is a different finite operator-algebraic
analogue: its content is the exact double-centralizer relation between the
\emph{joint} spacetime--gauge--incidence action and the residual typed
multiplicity algebra on the source-complete predictive carrier.  We do not
assume that the pair arises from a classical reductive dual pair.
\end{remark}

% =============================================================================
% =============================================================================
\section{Structural Standard Model, coercive stability, and common-action compatibility}
\label{sec:structural-stability-action}
% =============================================================================

\subsection{The minimal structural Standard-Model branch}
\label{sec:structural-sm}
% =============================================================================

We now identify one certified finite branch on which the residual multiplicity
algebra and its incidence data have the structural Standard-Model content used
below, together with a rank-three generation factor.  The construction is
intentionally separated from the duality theorem: the latter holds for every
source-complete typed packet, whereas the present section classifies one
particular finite active residual by explicit occurrence and positivity
conditions.

\subsubsection{Geometry shorting and the twelve-dimensional residual}

Let $\Phi$ be the twelve oriented $A_3$ roots associated with the tetrahedral
cell and let
\[
P=\R^\Phi,
\qquad
P_0=\left\{f\in P:\sum_{\alpha\in\Phi}f(\alpha)=0\right\}.
\]
Let $\mathscr T=\{c,\ell\}$ be two accepted categorical types and set
\[
t_0=2^{-1/2}(\mathbf e_c+\mathbf e_\ell),
\qquad
t_1=2^{-1/2}(\mathbf e_c-\mathbf e_\ell).
\]
After removal of the overall accepted mean, define
\begin{equation}
\boxed{
E_{\rm acc}^0
=
(P_0\otimes t_0)
\oplus
(\R\one_\Phi\otimes t_1)
\oplus
(P_0\otimes t_1).}
\label{eq:accepted-residual}
\end{equation}
The three summands have dimensions $11,1,11$.

Write
\[
E_{\rm ST}^{\rm cat}:=P_0\otimes t_0,
\qquad
E_{\rm int}^{\rm cat}
:=(\R\one_\Phi\otimes t_1)\oplus(P_0\otimes t_1).
\]
Let
\[
S_{\rm ST}:E_{\rm ST}\to\Hh_{\rm acc},
\qquad
S_{\rm int}:E_{\rm int}^{\rm cat}\to\Hh_{\rm acc}
\]
denote respectively the complete represented spacetime source and the accepted
internal categorical source on one common carrier.  The spacetime source may
contain, in addition to $E_{\rm ST}^{\rm cat}$, independently reconstructed
waiting-time, calibration, bracket, shift, or geometric-circulation directions.
Put
\[
G_{\rm ST}=S_{\rm ST}^*S_{\rm ST},
\qquad
G_{\rm int}=S_{\rm int}^*S_{\rm int},
\qquad
C=S_{\rm int}^*S_{\rm ST},
\]
and let $P_{\rm ST}$ be the orthogonal projection onto $\Ran S_{\rm ST}$.

\begin{theorem}[Universal geometry short and active residual]
\label{thm:geometry-short}
For every source-complete common realization, shorting the accepted internal
source by the complete represented spacetime source gives
\begin{equation}
\boxed{
G_{{\rm int}\mid{\rm ST}}
=
G_{\rm int}-CG_{\rm ST}^{\dagger}C^*
=
S_{\rm int}^*(I-P_{\rm ST})S_{\rm int}
\succeq0.}
\label{eq:geometry-short-general}
\end{equation}
The source-minimal internal categorical carrier is therefore
\begin{equation}
\boxed{
E_{{\rm int}\mid{\rm ST}}^{\min}
=
E_{\rm int}^{\rm cat}/\Ker G_{{\rm int}\mid{\rm ST}},
\qquad
\dim E_{{\rm int}\mid{\rm ST}}^{\min}
=
\rank G_{{\rm int}\mid{\rm ST}}.}
\label{eq:geometry-short-quotient}
\end{equation}

On the explicit orthogonal active realization, the accepted contrast synthesis
has scalar Gram
\[
S_{\rm acc}^*S_{\rm acc}=\vartheta I_{E_{\rm acc}^0},
\qquad \vartheta>0,
\]
and the additional spacetime scores are orthogonal to the accepted categorical
source.  Hence the mixed block vanishes on $E_{\rm int}^{\rm cat}$ and
\begin{equation}
\boxed{
G_{{\rm int}\mid{\rm ST}}=\vartheta I_{12}.}
\label{eq:geometry-short}
\end{equation}
On the tetrahedrally equivariant branch, the categorical residual carries the
multiplicity-free representation
\begin{equation}
\boxed{
E_{\rm int}^{\rm cat}\otimes\C
\cong
\C_{\rm type}^{+}
\oplus W_+
\oplus(V_2)_+
\oplus W_-
\oplus(W\otimes{\rm sgn})_-.}
\label{eq:residual-S4}
\end{equation}
Here $\dim W=3$ and $\dim V_2=2$.  The odd standard block $W_-$ is the
internal endpoint module, while the odd sign-twisted block is the distinct loop
module.  If the restriction of $G_{{\rm int}\mid{\rm ST}}$ to the endpoint
block is strictly positive, that block survives the source-minimal quotient at
complex rank three.
\end{theorem}

\begin{proof}
The first identity is the Schur/Anderson--Trapp short of the joint source Gram;
equivalently, it is the Gram of the orthogonal residual
$(I-P_{\rm ST})S_{\rm int}$.  Positivity and the quotient/rank statement follow
immediately.  On the explicit active product carrier the accepted categorical
summands are orthogonal with common weight $\vartheta$, and all additional
spacetime scores lie in the independently adjoined relational port, giving
\eqref{eq:geometry-short}.

The representation decomposition is the
$S_4\times\langle\iota\rangle$ decomposition of the twelve oriented roots:
the reversal-even six-dimensional sector is
$\mathbf1\oplus W\oplus V_2$, while the reversal-odd sector is
$W\oplus(W\otimes{\rm sgn})$.  Removing the even root mean and adjoining the
scalar type contrast gives \eqref{eq:residual-S4}.  On an equivariant common
carrier the spacetime range, its orthogonal projection, and hence the shorted
Gram are equivariant, so the inequivalent endpoint and loop irreducibles remain
reducing.  Strict positivity on the endpoint block prevents its collapse in
\eqref{eq:geometry-short-quotient}.
\end{proof}

The essential point is that the rank-three internal endpoint block is defined
\emph{after} the complete geometric source has been removed on the same joint
Gram.  It is therefore neither a second name for the spatial $W_4$ carrier nor
a dimension count copied from the external tetrahedral representation.  The
scalar active realization provides a particularly transparent nonempty branch,
but the source-short formulation itself is not tied to an isotropic accepted
metric.

\subsubsection{External multiplicity census and the first internal-factor obstruction}
\label{sec:multiplicity-birth}

The external tetrahedral action and the internal gauge action must be kept on
opposite sides of the centralizer.  Let $\mathcal V_r$ be the future-minimal
reachable--observable residual \emph{physical carrier} spanned by represented
source words of length at most $r$, after the complete spacetime and
physical-null shorts.  This is not the regular word Hilbert space
$L^2(\A,\tau_{\rm reg})$: by Corollary~\ref{cor:word-multiplicity-firewall},
large invariant multiplicities of the latter do not by themselves create
physical internal factors.  Decompose the external $S_4$ action as
\begin{equation}
\mathcal V_r
\cong
\bigoplus_{\lambda\in\widehat{S_4}}
V_\lambda\otimes M_{\lambda,r},
\qquad
m_{\lambda,r}:=\dim_\C M_{\lambda,r}.
\label{eq:word-depth-isotypic}
\end{equation}

\begin{theorem}[External isotypic algebra and multiplicity census]
\label{thm:active-isotypic}
At every finite word depth,
\begin{align}
\A_{{\rm ext},r}
&=
C^*(\rho_r(S_4))
\cong
\bigoplus_\lambda
M_{\dim V_\lambda}(\C)\otimes I_{M_{\lambda,r}},
\label{eq:external-isotypic}\\
\A_{{\rm ext},r}'
&\cong
\bigoplus_\lambda
I_{V_\lambda}\otimes M_{m_{\lambda,r}}(\C).
\label{eq:external-isotypic-commutant}
\end{align}
Consequently an internal simple factor $M_n(\C)$ can occur on the relative
commutant side only if some external irreducible occurs with multiplicity at
least $n$.

On the categorical degree-one carrier,
\[
\mathcal V_{\rm cat}
\cong
\C_{\rm type}\oplus(W\otimes\C^2)\oplus V_2\oplus W^{\rm sgn},
\]
so the multiplicities are $(1,2,1,1)$ and
\begin{equation}
\boxed{
\A_{{\rm ext},1}'
\cong
\C\oplus M_2(\C)\oplus\C\oplus\C.}
\label{eq:degree-one-commutant}
\end{equation}
If one additional external-trivial private contrast line is retained, the
multiplicity vector becomes $(2,2,1,1)$ and the relative commutant is
\begin{equation}
M_2(\C)\oplus M_2(\C)\oplus\C\oplus\C.
\label{eq:degree-one-private-commutant}
\end{equation}
Neither algebra contains a unital copy of $M_3(\C)$.
\end{theorem}

\begin{proof}
Finite-dimensional isotypic decomposition gives
\eqref{eq:word-depth-isotypic}.  Burnside's theorem gives the full matrix
algebra on each irreducible external factor and the identity on its
multiplicity factor; Schur's lemma gives the commutant.  A unital
$*$-representation of $M_n(\C)$ is a direct sum of $n$-dimensional defining
representations, so no simple commutant block of size smaller than $n$ can
contain it.  The two displayed finite censuses follow from the stated
multiplicities.
\end{proof}

\begin{countertheorem}[Degree-one internal-colour obstruction]
\label{thm:degree-one-colour-no-go}
The populated categorical degree-one residual cannot by itself support a
commuting internal colour factor $M_3(\C)$.  The same remains true after
adjoining only one external-trivial private contrast line.
\end{countertheorem}

\begin{proof}
This is the final assertion of Theorem~\ref{thm:active-isotypic}: every simple
summand of the corresponding relative commutant has size at most two.
\end{proof}

Let $K_r$ denote the labelled source-word Gram whose quotient span is
$\mathcal V_r$, and define the external-trivial multiplicity
\begin{equation}
m_{0,r}:=\dim_\C\Hom_{S_4}(\mathbf1,\mathcal V_r).
\label{eq:trivial-multiplicity}
\end{equation}
The first colour-eligible depth is
\begin{equation}
r_{\rm col}^{\rm birth}
:=
\min\{r:m_{0,r}\ge3\},
\label{eq:colour-birth-depth}
\end{equation}
with value $\infty$ if no such depth occurs before the word hierarchy becomes
flat.

\begin{theorem}[Flat-depth multiplicity-birth obstruction]
\label{thm:flat-depth-colour}
Suppose
\begin{equation}
\rank K_r=\rank K_{r+1}
\label{eq:word-flatness}
\end{equation}
and $m_{0,r}\le2$.  Then $\mathcal V_s=\mathcal V_r$ for every $s\ge r$ and
no later physical source word can produce an external-trivial internal
$M_3(\C)$ factor.  In particular $r_{\rm col}^{\rm birth}=\infty$ on that
source algebra.
\end{theorem}

\begin{proof}
By construction $\mathcal V_{r+1}$ is the span of $\mathcal V_r$ and all
one-letter images of $\mathcal V_r$.  Equality of Gram ranks therefore gives
$\mathcal V_{r+1}=\mathcal V_r$, so $\mathcal V_r$ is invariant under every
source letter.  Induction gives $\mathcal V_s=\mathcal V_r$ for all later
$s$.  The multiplicity census is consequently frozen, and
Theorem~\ref{thm:active-isotypic} excludes an internal $M_3(\C)$ block.
\end{proof}

A positive multiplicity birth is necessary but not sufficient for colour: it
creates only a coherent multiplicity carrier.  The physical internal word
algebra must still connect its minimal projections.

\subsubsection{Source-minimal routers and exact internal landing}
\label{sec:internal-landing}

\begin{theorem}[Connected-router full-matrix theorem]
\label{thm:connected-router}
Let
\[
\mathcal D_n=\Span\{p_1,\ldots,p_n\}\cong\C^n
\]
be the diagonal algebra of mutually orthogonal rank-one projections on
$\C^n$, and let $u\in M_n(\C)$.  Form the undirected support graph with an
edge $i-j$ whenever $p_iup_j\ne0$ or $p_jup_i\ne0$.  Then
\begin{equation}
\boxed{
C^*(\mathcal D_n,u)=M_n(\C)
\quad\Longleftrightarrow\quad
\text{the support graph of $u$ is connected}.}
\label{eq:router-full-matrix}
\end{equation}
If the graph is disconnected, every connected component is a reducing
subspace and the generated algebra is a proper block-diagonal subalgebra.
\end{theorem}

\begin{proof}
For every support edge, $p_iup_j$ is a nonzero scalar multiple of the matrix
unit $E_{ij}$; adjoints give $E_{ji}$.  Products along a graph path generate
all matrix units between vertices in the same connected component.  Thus a
connected graph generates all of $M_n(\C)$.  Conversely the sum of the
$p_i$ over any connected component is a nontrivial central reducing
projection when the graph is disconnected.
\end{proof}

\begin{corollary}[One-router source minimality]
\label{cor:one-router-minimality}
Once a coherent $n$-dimensional multiplicity carrier and its diagonal minimal
projections have occurred, one cyclic nonnormal router is sufficient to
generate $M_n(\C)$.  Without any off-diagonal word the generated algebra
remains commutative.  Thus one router is source-minimal at fixed multiplicity
carrier.
\end{corollary}

Let $T\cong\C$ be the surviving external-trivial type line and let
$H_{\rm priv}\cong\C^2$ be a faithful external-trivial private plane on the
same source carrier.  Let $p_T$ and $p_H$ denote the corresponding
projections.  On the standard-isotypic multiplicity plane let $p_{\rm t}$ and
$p_{\rm h}$ be rank-one projections onto unit vectors $t$ and $h$.

\begin{lemma}[Weak multiplicity router]
\label{lem:weak-M2}
If
\begin{equation}
0<\abs{\langle t,h\rangle}<1,
\label{eq:weak-overlap}
\end{equation}
then
\begin{equation}
\boxed{C^*(p_{\rm t},p_{\rm h})=M_2(\C).}
\label{eq:weak-M2}
\end{equation}
With only one rank-one projection the generated algebra is $\C^2$.
\end{lemma}

\begin{proof}
The two projections have no common one-dimensional reducing subspace, so
Burnside's theorem gives $M_2(\C)$.  A single projection generates only its
two spectral projections.
\end{proof}

For an actual operation word $u$ on $T\oplus H_{\rm priv}$ define
\begin{equation}
\omega_{\rm col}(u)
=
\norm{p_Hup_T}_{\HS}^2+\norm{p_Tup_H}_{\HS}^2.
\label{eq:colour-bridge}
\end{equation}

\begin{theorem}[Necessary and sufficient colour-bridge criterion]
\label{thm:colour-bridge}
Assume the represented private support words generate $M_2(H_{\rm priv})$.
Then
\begin{equation}
\boxed{
C^*(\C p_T\oplus M_2(H_{\rm priv}),u)
=
M_3(\C)
\quad\Longleftrightarrow\quad
\omega_{\rm col}(u)>0.}
\label{eq:colour-M3}
\end{equation}
In the zero-bridge case the algebra is exactly
$\C\oplus M_2(\C)$ and has a nontrivial central projection.
\end{theorem}

\begin{proof}
A nonzero corner between the type line and the private $2\times2$ block,
together with its adjoint and the private matrix units, generates every
rectangular matrix unit and hence $M_3(\C)$.  If both cross corners vanish,
the line and plane remain reducing and the algebra is precisely the direct
sum $\C\oplus M_2(\C)$.
\end{proof}

The bridge need not be a primitive operation when the historical route table
contains reciprocal access through one common mediator.  Let $G$ be the
external finite symmetry, let $V$ be an irreducible unitary $G$-module of
dimension $d$, and let the mediator carrier be
$\mathcal E=V\otimes\mathcal N$ with $n=\dim\mathcal N$.  Suppose actual entry
and reverse-exit operations have the covariant form
\[
 A_g=(v_g\otimes I)A:T\to\mathcal E,
 \qquad
 B_h^*=B^*(v_h^*\otimes I):\mathcal E\to H_{\rm priv},
\]
and that the actual mediator dynamics contains
$D=I_V\otimes h$ with $h=h^*$.  Put
\[
 r_{g,k}=B^*(v_g\otimes h^k)A,
 \qquad 0\le k<n,
\]
\[
 \rho_A=\Tr_V(AA^*),\qquad
 \rho_B=\Tr_V(BB^*),\qquad
 W_A=\sum_{k=0}^{n-1}h^k\rho_Ah^k.
\]

\begin{theorem}[Finite reciprocal-return criterion]
\label{thm:reciprocal-return}
For the declared route bank,
\begin{equation}
 \boxed{
 \mathfrak m_{\rm ret}
 :=\frac1{|G|}\sum_{g\in G}\sum_{k=0}^{n-1}
       \norm{r_{g,k}}_{\HS}^2
 =\frac1d\Tr(\rho_BW_A)\ge0.}
 \label{eq:reciprocal-return-mass}
\end{equation}
The following are equivalent:
\begin{enumerate}[label=\textup{(\roman*)},leftmargin=2.8em]
\item an admitted word in the route bank has a nonzero
      $H_{\rm priv}\leftarrow T$ corner;
\item some $r_{g,k}$ is nonzero;
\item $\mathfrak m_{\rm ret}>0$;
\item the $h$-cyclic multiplicity spaces generated by $\Ran\rho_A$ and
      $\Ran\rho_B$ are not orthogonal.
\end{enumerate}
If these conditions fail, an explicit reducing projection separates $T$ from
$H_{\rm priv}$ for the entire declared bank, so arbitrary repeated excursions
cannot create a colour bridge within that bank.  If they hold, a witness has
at most $n+1$ route letters.
\end{theorem}

\begin{proof}
Schur averaging on the irreducible $V$ factor gives
\[
 \frac1{|G|}\sum_g
 (v_g\otimes I)AA^*(v_g^*\otimes I)
 =\frac{I_V}{d}\otimes\rho_A.
\]
Taking the Hilbert--Schmidt norm of $B^*(v_g\otimes h^k)A$ and summing yields
\eqref{eq:reciprocal-return-mass}.  Since it is a finite sum of squared norms,
positivity is equivalent to a nonzero returned word.  The range of $W_A$ is
$\Span\{h^k\Ran\rho_A:0\le k<n\}$ by Cayley--Hamilton.  Thus
$\Tr(\rho_BW_A)=0$ exactly when the two cyclic multiplicity spaces are
orthogonal.  In that case the projection onto
$T\oplus(V\otimes\Ran W_A)$ commutes with the route generators and their
adjoints while annihilating $H_{\rm priv}$, excluding all longer returns.  The
finite horizon $k<n$ gives the stated word-length bound.
\end{proof}

\begin{corollary}[Port-relative reachability forces colour]
\label{cor:reciprocal-return-force}
If $W_A\succeq\kappa I_{\mathcal N}$ with $\kappa>0$ and $B\ne0$, then
\begin{equation}
 \boxed{
 \mathfrak m_{\rm ret}\ge\frac{\kappa}{d}\norm B_{\HS}^2,
 \qquad
 \max_{g,k<n}\norm{r_{g,k}}_{\HS}^2
 \ge\frac{\kappa}{dn}\norm B_{\HS}^2.}
 \label{eq:reciprocal-return-margin}
\end{equation}
Hence the colour bridge follows from the upstream route structure and
Theorem~\ref{thm:colour-bridge}.  In particular, for a multiplicity-one
mediator, two nonzero reciprocal ports force a two-leg colour return and no
independent direct bridge is required.
\end{corollary}

Blockwise matrix generation still does not by itself guarantee that all
central blocks are independent.  The only ambiguity after the nonisomorphic
$M_3$ and $M_2$ blocks have landed concerns the two scalar sectors.

\begin{theorem}[Central separation of the landed internal blocks]
\label{thm:central-separation}
Let
$\mathcal B\subseteq M_3(\C)\oplus M_2(\C)\oplus\C\oplus\C$
be a unital $C^*$-subalgebra whose projection onto each displayed factor is
surjective.  Then exactly one of the following occurs:
\begin{align}
 \mathcal B&=M_3(\C)\oplus M_2(\C)\oplus\C\oplus\C,
 &\dim_\C\mathcal B&=15,
 \label{eq:central-separated-15}\\
 \mathcal B&=\{(A,B,z,z):A\in M_3(\C),\ B\in M_2(\C),\ z\in\C\},
 &\dim_\C\mathcal B&=14.
 \label{eq:central-locked-14}
\end{align}
If $\chi_1,\chi_2$ are the scalar characters and $(b_j)$ is any generating
bank, then the full product occurs exactly when
\begin{equation}
 \boxed{\eta_{\rm cen}:=\sum_j|\chi_1(b_j)-\chi_2(b_j)|^2>0.}
 \label{eq:central-separation-margin}
\end{equation}
In particular, represented central projectors onto the two scalar sectors
discharge the condition automatically.
\end{theorem}

\begin{proof}
A unital subdirect product of finite simple matrix algebras can correlate only
isomorphic simple factors; a proof is given in
Appendix~\ref{app:historical-source-selection}.  The $M_3$ and $M_2$ factors
are therefore independent and only the two scalar factors may be identified.
They are independent exactly when their characters differ on the generated
algebra, which is equivalent to \eqref{eq:central-separation-margin}.
\end{proof}

Define $\A_{{\rm int},r}^{\rm word}$ to be the $C^*$-algebra generated by
all represented physical word compressions at depth $r$ that commute with the
external tetrahedral action.

\begin{theorem}[Relative-commutant saturation]
\label{thm:internal-seed-saturation}
Suppose that at one finite word depth the relevant external multiplicity
vector is
\begin{equation}
(m_0,m_W,m_{V_2},m_{W^{\rm sgn}})=(3,2,1,1),
\label{eq:SM-multiplicity-census}
\end{equation}
that the external-trivial three-space is coherently realized as
$T\oplus H_{\rm priv}$, and that:
\begin{enumerate}[label=\textup{(I\arabic*)},leftmargin=3em]
\item the private support words generate $M_2(H_{\rm priv})$ and either one
      represented word has $\omega_{\rm col}>0$, or an admitted reciprocal
      mediator bank satisfies $\mathfrak m_{\rm ret}>0$ in
      Theorem~\ref{thm:reciprocal-return};
\item the standard multiplicity plane contains two distinct nonorthogonal
      represented rank-one projections;
\item the $V_2$ and $W^{\rm sgn}$ external blocks have multiplicity one;
\item the two scalar sectors are centrally separated, equivalently the
      generated algebra contains their independent central supports or passes
      $\eta_{\rm cen}>0$ in Theorem~\ref{thm:central-separation}.
\end{enumerate}
Then
\begin{equation}
\boxed{
\A_{{\rm int},r}^{\rm word}
=
\A_{{\rm ext},r}'
\cong
M_3(\C)\oplus M_2(\C)\oplus\C\oplus\C.}
\label{eq:active-internal-algebra}
\end{equation}
Its complex dimension is $15$, its centre has dimension four, and the two
nonabelian simple factors are intrinsically distinguished by their matrix
sizes.  Hence the labels ``colour'' and ``weak'' are canonical up to inner
automorphism once the multiplicity packet has landed.
\end{theorem}

\begin{proof}
Condition~\textup{(I1)} and Theorem~\ref{thm:colour-bridge} generate the
$M_3$ block; when the reciprocal route alternative is used,
Theorem~\ref{thm:reciprocal-return} supplies the required represented bridge.
Condition~\textup{(I2)} and Lemma~\ref{lem:weak-M2} generate the $M_2$ block.
The multiplicity-one sectors give surjective scalar compressions, and
condition~\textup{(I4)} excludes their diagonal identification by
Theorem~\ref{thm:central-separation}.  Thus the generated internal algebra is
the full product, which is exactly the relative commutant determined by the
multiplicity census in Theorem~\ref{thm:active-isotypic}.
\end{proof}

\begin{theorem}[Exact internal-algebra deficit alternatives]
\label{thm:internal-deficit-alternatives}
On a coherent carrier with multiplicity census $(3,2,1,1)$, assume the colour
three-space initially carries $M_2(\C)\oplus\C$, the weak multiplicity plane
carries one rank-one projection, and the displayed colour, weak and two scalar
central supports occur independently in the represented internal algebra.
Then exactly the following dimensions occur:
\begin{center}
\begin{tabular}{L{0.47\textwidth}c}
\toprule
represented internal words & $\dim_\C\A_{\rm int}^{\rm word}$\\\midrule
neither colour bridge nor weak router & $9$\\
weak router but no colour bridge & $11$\\
colour bridge but no weak router & $13$\\
both bridge and router & $15$\\
\bottomrule
\end{tabular}
\end{center}
The terminal value $15$ is equivalent to saturation of the complete relative
commutant in Theorem~\ref{thm:internal-seed-saturation}.  Each failure is
therefore an exact finite source-word rank deficit rather than a discarded
branch.  If independent central supports have not yet been established,
Theorem~\ref{thm:central-separation} must be applied before componentwise
dimensions are added; in particular, the terminal blockwise-surjective branch
may still be the $14$-dimensional scalar-locked algebra
\eqref{eq:central-locked-14}.
\end{theorem}

\begin{proof}
Without the bridge the colour block is $M_2\oplus\C$, of dimension five; with
the bridge it is $M_3$, of dimension nine.  One weak rank-one projection
generates $\C^2$, of dimension two; a noncommuting second projection generates
$M_2$, of dimension four.  The two multiplicity-one sectors add two scalar
dimensions, giving $9,11,13,15$.
\end{proof}

\begin{corollary}[Generic saturation at fixed multiplicity census]
\label{cor:generic-router-saturation}
Fix the multiplicity census $(3,2,1,1)$ and the diagonal minimal projections
on the $3$- and $2$-dimensional multiplicity carriers.  In the affine parameter
space of internal router operators on those carriers, the set that generates
$M_3(\C)\oplus M_2(\C)$ is Zariski open and dense.  Its complement is a
finite union of proper linear subspaces corresponding to disconnected support
graphs.  This genericity is conditional on the multiplicity carriers having
occurred; it does not assert that a microscopic source samples the full router
parameter space.
\end{corollary}

\begin{proof}
For fixed $n$, failure of connectedness means that there exists a nontrivial
partition of $\{1,\ldots,n\}$ across which all matrix entries of the router
vanish in both directions.  For each partition this is a proper linear
subspace of $M_n(\C)$, and there are finitely many partitions.  The complement
is therefore Zariski open and dense.  Apply this independently to $n=3$ and
$n=2$ and use Theorem~\ref{thm:connected-router}.
\end{proof}

\begin{remark}[Selection content of the internal landing]
The multiplicity census, bridge and router conditions are not names for a
preselected gauge group.  Theorem~\ref{thm:degree-one-colour-no-go} excludes
colour at the populated degree-one census,
Theorem~\ref{thm:flat-depth-colour} gives a terminating no-go when multiplicity
birth never occurs, and Theorem~\ref{thm:internal-deficit-alternatives} returns
the exact smaller algebra whenever one of the two connecting words is absent.
Conversely, once the required multiplicity carriers exist, one bridge and one
router are source-minimal for the nonabelian blocks, and on reciprocal-route
branches the bridge itself is forced by the positive return criterion rather
than separately postulated.  Full product saturation additionally requires the
finite central-separation test of Theorem~\ref{thm:central-separation}.  The
determinant test below is applied only after this multiplicity-side landing;
applying it to the external Burnside blocks would act on the wrong side of the
centralizer.
\end{remark}

\begin{remark}[Relation to the finite algebra of almost-commutative geometry]
The algebra in \eqref{eq:active-internal-algebra} is the complex relative
commutant of the present joint action.  It should not be identified with the
real finite coordinate algebra used in the standard almost-commutative
description of the Standard Model.  The two constructions organize the
internal sector differently even though they recover the same physical
nonabelian dimensions and quotient gauge group on the branch described below.
\end{remark}

\subsubsection{The determinant seed and the global gauge group}

The global form of the Standard-Model gauge group is not fixed by its Lie
algebra alone.  In conventional field theory one may quotient
$\SU(3)\times\SU(2)\times\U(1)$ by different subgroups of the common
$\Z_6$ centre, with physically distinct line-operator spectra
\cite{Hucks1991,Tong2017}.  The statement below is therefore branch-specific:
the determinant seed reconstructs the faithful $\Z_6$ quotient rather than
assuming it as the definition of the Standard Model.

Let $C\cong\C^3$ and $W_2\cong\C^2$ denote the nonabelian multiplicity
carriers after the relative-commutant landing of
Theorem~\ref{thm:internal-seed-saturation}.  The determinant relation is not
inserted as an independent group constraint, nor is it applied to the external
tetrahedral Burnside blocks.  It is the nonzero branch of a canonical
alternating projection of the represented incidence tensor on the landed
multiplicity carriers.

\begin{proposition}[Canonical determinant incidence]
\label{prop:main-det-incidence}
Consider the represented incidence space
\[
\Hom(C\otimes W_2\otimes W_2,\Lambda^2C^*).
\]
Its component that is fully alternating in the three colour and two weak
arguments is one-dimensional and canonically isomorphic to
\[
(\det C\otimes\det W_2)^*.
\]
Consequently every represented incidence tensor has a canonical alternating
projection $\tau$.  If $\tau\neq0$, it defines a determinant seed uniquely
up to a nonzero scalar and, after unit normalization, up to phase.  All nonzero
normalized choices have the same unitary stabilizer.  If the alternating
projection vanishes, the determinant-preserving branch is absent rather than
being supplied by completion.
\end{proposition}

\begin{proof}
Complete alternation is the equivariant projection onto the unique
$(\det C\,\det W_2)^{-1}$ symmetry line.  The explicit incidence map and
normalization are recorded in Appendix~\ref{app:anomaly}; in particular,
Proposition~\ref{prop:det-incidence} identifies this line isomorphically with
$(\det C\otimes\det W_2)^*$.
\end{proof}

\begin{definition}[Determinant seed]
On the nonzero branch of Proposition~\ref{prop:main-det-incidence}, normalize
its alternating projection to a unitary isomorphism
\begin{equation}
\tau:\det C\otimes\det W_2\xrightarrow{\sim}\C.
\label{eq:determinant-seed}
\end{equation}
Only the line and its stabilizer are physically relevant; multiplying $\tau$
by an overall phase leaves the stabilizer unchanged.  Thus $\tau$ is a
reconstructed incidence shadow on this branch, not an independently chosen
Standard-Model datum.
\end{definition}

\begin{theorem}[Gauge-group reconstruction]
\label{thm:gauge-group}
The unitary automorphism group of the determinant seed is
\begin{equation}
\boxed{
\Aut_{\U}(C,W_2,\tau)
=
S(\U(3)\times\U(2))
\cong
\frac{\SU(3)\times\SU(2)\times\U(1)}{\Z_6}.}
\label{eq:SM-gauge-group}
\end{equation}
\end{theorem}

\begin{proof}
A pair $(A,B)\in\U(C)\times\U(W_2)$ preserves $\tau$ exactly when
$\det A\,\det B=1$.  The map
\[
\Phi(g_3,g_2,z)=(z^{-2}g_3,z^3g_2)
\]
has image in this subgroup and kernel
\[
\{(z^2I_3,z^{-3}I_2,z):z^6=1\}\cong\Z_6.
\]
\end{proof}

\begin{proposition}[Positive operational source criterion for the determinant group]
\label{prop:operational-determinant-source}
Let $\chi(A,B)=\det A\det B$ on $\U(C)\times\U(W_2)$.  Suppose an actual typed
incidence instrument contains two marked coherent amplitude sectors on which
this group acts by the characters $1$ and $\chi^{-1}$.  Let $J_{01}$ denote
the corresponding Choi cross block and put
\begin{equation}
 \boxed{h_{\det}:=\norm{J_{01}}_{\HS}^2.}
 \label{eq:determinant-cross-mass}
\end{equation}
If $h_{\det}>0$, the stabilizer of the marked positive subpacket inside
$\U(3)\times\U(2)$ is exactly
\begin{equation}
 \boxed{\ker\chi=S(\U(3)\times\U(2)).}
 \label{eq:operational-determinant-group}
\end{equation}
If $h_{\det}=0$, the two diagonal character-sector masses alone do not fix
their relative phase and therefore do not imply the determinant subgroup.
\end{proposition}

\begin{proof}
Under $(A,B)$ the cross block transforms by the relative character
$\chi(A,B)^{-1}$.  A nonzero marked cross block is therefore fixed exactly
when $\chi(A,B)=1$.  If the cross block vanishes, independent phase rotation of
the two diagonal sectors leaves their positive masses unchanged.
\end{proof}

\begin{remark}[Determinant incidence is genuinely higher arity]
\label{rem:determinant-arity}
The determinant character cannot be manufactured by arbitrary temporal depth
inside a single native carrier containing only the defining colour and weak
blocks and neutral sectors.  Appendix~\ref{app:historical-source-selection}
proves a sharp charged-slot bound: the determinant mode is absent below five
total defining tensor slots and occurs at five.  The incidence space in
Proposition~\ref{prop:main-det-incidence} has precisely three colour and two
weak charged legs, so its arity is minimal for this character.  This is an
occurrence statement, not an assertion that every microscopic source supplies
the required five-leg coherent port.
\end{remark}

\begin{proposition}[Faithful matter realization of the determinant quotient]
\label{prop:faithful-SM-quotient}
Normalize the abelian charge by $y=6Y$ and write all fermions as left-handed
Weyl fields.  On one generation the represented matter packet is
\[
Q:(\mathbf3,\mathbf2)_1,
\qquad
u^c:(\overline{\mathbf3},\mathbf1)_{-4},
\qquad
d^c:(\overline{\mathbf3},\mathbf1)_2,
\]
\[
L:(\mathbf1,\mathbf2)_{-3},
\qquad
e^c:(\mathbf1,\mathbf1)_6,
\qquad
H:(\mathbf1,\mathbf2)_3,
\qquad
[\nu^c:(\mathbf1,\mathbf1)_0].
\]
The common kernel of these representations of
$\widetilde G=\SU(3)\times\SU(2)\times\U(1)_y$ is exactly
\begin{equation}
\boxed{
K
=
\left\langle
\left(e^{2\pi i/3}I_3,-I_2,e^{i\pi/3}\right)
\right\rangle
\cong\Z_6.}
\label{eq:faithful-SM-kernel}
\end{equation}
Consequently the represented matter packet is faithful precisely for
$\widetilde G/K$, which is the same global group selected intrinsically by the
determinant seed in Theorem~\ref{thm:gauge-group}.
\end{proposition}

\begin{proof}
A common-kernel element must be central, so write it as
$(\omega_3^aI_3,(-1)^bI_2,e^{i\theta})$.  The charged singlet $e^c$ gives
$\theta=m\pi/3$ modulo $2\pi$; the lepton doublet gives
$b\equiv m\pmod2$; and the colour antitriplet $d^c$ gives
$a\equiv m\pmod3$.  Substitution in the remaining representations shows that
precisely the six powers of the generator in
\eqref{eq:faithful-SM-kernel} act trivially.  The optional neutral singlet is
trivial under all factors and adds no further restriction.
\end{proof}

\begin{remark}[Relative commutant versus physical gauge group]
The two scalar summands in
$M_3(\C)\oplus M_2(\C)\oplus\C\oplus\C$ are components of the complex
relative commutant; they are not two additional independent physical
$\U(1)$ gauge groups.  The physical gauge group is obtained by preserving the
represented determinant/incidence structure on the nonabelian multiplicity
carriers and then requiring faithful action on the typed matter packet.
Theorem~\ref{thm:gauge-group} ties the central phases by
$\det A\det B=1$, while Proposition~\ref{prop:faithful-SM-quotient} removes the
common $\Z_6$ kernel.  Thus the passage from the relative commutant to
$G_{\rm SM}$ is a stabilizer-and-faithfulness construction, not the operation
of taking the full unitary group of
$\A_{{\rm int},r}^{\rm word}$.
\end{remark}

\subsubsection{Anomaly-forced hypercharge}

Anomaly cancellation is well known to constrain, and under appropriate
representation assumptions essentially determine, Standard-Model charge
assignments \cite{GengMarshak1989,MinahanRamondWarner1990}.  In the
noncommutative-geometric Standard Model it is closely tied to the
unimodularity reduction of the gauge group
\cite{AlvarezGraciaBondiaMartin1995}.  Here the tensor types are already fixed
by the preceding determinant/incidence construction; the theorem determines
only the remaining central weight ratio.

Let the infinitesimal central action on $C\oplus W_2$ be
\[
iaI_C\oplus ibI_{W_2}.
\]
Define the tensor types
\begin{equation}
Q=C\otimes W_2,\qquad
u=\Lambda^2C^*,\qquad
d=C,\qquad
L=W_2^*,\qquad
e=\Lambda^2W_2^*,\qquad
H=W_2.
\label{eq:SM-types}
\end{equation}
Their central weights are
\begin{equation}
(a+b,-2a,a,-b,-2b,b).
\label{eq:central-weights}
\end{equation}

\begin{theorem}[Anomaly-forced central weights]
\label{thm:hypercharge}
Write right-handed fermions as left-handed conjugates when computing chiral
anomalies.  The local
$\SU(3)^2\U(1)$,
$\SU(2)^2\U(1)$,
gravitational--$\U(1)$,
and cubic $\U(1)$ anomalies vanish simultaneously if and only if
\begin{equation}
\boxed{3a+2b=0.}
\label{eq:anomaly-condition}
\end{equation}
Up to overall reversal of the $\U(1)$ orientation, the primitive integral
solution is
\[
(a,b)=(-2,3).
\]
With $y=6Y$, the resulting physical weights are
\begin{equation}
\boxed{
y(Q,u,d,L,e,H)=(1,4,-2,-3,-6,3).}
\label{eq:hypercharges}
\end{equation}
If a neutral right-handed singlet is retained, $y(\nu)=0$.
\end{theorem}

\begin{proof}
After charge conjugation of the right-handed fields,
\begin{align*}
\mathcal A_{\SU(3)^2Y}
&=
\mathcal A_{\SU(2)^2Y}
=
\frac12(3a+2b),\\
\mathcal A_{{\rm grav}^2Y}
&=
3(3a+2b),\\
\mathcal A_{Y^3}
&=
3(3a+2b)(3a^2+2b^2).
\end{align*}
The pure colour anomaly cancels as $2-1-1$, and the weak doublet count is even.
For a nontrivial abelian action $3a^2+2b^2>0$, so simultaneous local anomaly
cancellation is equivalent to \eqref{eq:anomaly-condition}.  The primitive
coprime integral solution is $(-2,3)$ up to sign; substitution gives
\eqref{eq:hypercharges}.
\end{proof}

\subsubsection{Matter occurrence, generations, and the finite Dirac operator}

In the usual almost-commutative Standard Model, three generations are encoded
by tripling the finite fermion carrier
\cite{vanDenDungenVanSuijlekom2012,vanSuijlekom2025}.  Other algebraic
programmes have sought a structural origin for family replication, including
complex-octonionic constructions \cite{Furey2018}, Clifford-ideal models of
one generation \cite{Gresnigt2020}, an intrinsic $S_3$ family action in a
Clifford algebra \cite{Gresnigt2026}, and the recent $E_7$ three-generation
decomposition \cite{Baez2026}; for a broader recent algebraic survey see
\cite{FureyRoadmap2025}.  The mechanism below is
neither: the rank-three family carrier is a source-theoretic residual after
spacetime shorting and remains subject to a positive lower-frame certificate.

The Clifford and internal packets must occur on one carrier.  Let
$J=P_+-P_-$ be the protected pre-central grading and let
$\sigma_0,\ldots,\sigma_3$ be Hermitian Clifford axes generating the external
matrix factor.

\begin{definition}[Clifford matter margin]
\label{def:clifford-matter}
Define
\begin{equation}
\boxed{
p_{\rm Cl}
=
\frac{1}{2\dim\Hh}
\sum_{\mu=0}^{3}
\norm{P_-\sigma_\mu P_+}_{\HS}^2.}
\label{eq:clifford-matter}
\end{equation}
\end{definition}

\begin{proposition}[Clifford occurrence test]
\label{prop:clifford-occurrence}
One has $p_{\rm Cl}>0$ exactly when at least one represented Clifford axis
changes the protected grading.  If $p_{\rm Cl}=0$, the grading reduces the
external Clifford action and its residual action lies in the normal
multiplicity commutant.
\end{proposition}

\begin{proof}
Every summand in \eqref{eq:clifford-matter} is nonnegative.  The sum vanishes
exactly when $P_-\sigma_\mu P_+=0$ for every $\mu$, and self-adjointness then
also gives $P_+\sigma_\mu P_-=0$.  Thus $J$ commutes with the generated
external algebra.
\end{proof}

On the active residual of Theorem~\ref{thm:geometry-short}, the
reversal-odd standard module $W_-$ has rank three.  We denote its supported
complexification by $G_{\rm gen}\cong\C^3$.

\begin{proposition}[Endpoint generation carrier]
\label{prop:generation-carrier}
Suppose the restriction of the shorted Gram
$G_{{\rm int}\mid{\rm ST}}$ to the reversal-odd endpoint block $W_-$ in
\eqref{eq:residual-S4} is strictly positive.  Then
\begin{equation}
\boxed{\dim_\C G_{\rm gen}=3.}
\label{eq:three-generations}
\end{equation}
The block $(W\otimes{\rm sgn})_-$ is a distinct loop-provenance triplet and
cannot be counted as a second generation source.  On the explicit orthogonal
active realization the hypothesis is automatic because
$G_{{\rm int}\mid{\rm ST}}=\vartheta I_{12}$.
\end{proposition}

\begin{proof}
The endpoint and cycle projectors canonically separate the two odd irreducible
summands in \eqref{eq:residual-S4}.  The endpoint standard representation has
complex dimension three.  Strict positivity of the restricted shorted Gram
prevents any endpoint direction from entering the source kernel, while the
sign-twisted standard block lies in the distinct cycle projector.
\end{proof}

Let $A$ be a source-saturated self-adjoint full-cell generator on the common
carrier.  Its grading-odd part is
\begin{equation}
D_F
=
P_-AP_+ + P_+AP_-.
\label{eq:finite-dirac}
\end{equation}

\begin{definition}[Odd provenance completeness]
The branch is \emph{odd-provenance complete} when every future-visible
grading-odd action of the saturated full-cell carrier lies in the represented
range of \eqref{eq:finite-dirac}; equivalently, the orthogonal odd innovation
Gram vanishes.
\end{definition}

\begin{proposition}[Canonical finite Dirac operator]
\label{prop:finite-dirac}
On an odd-provenance-complete saturated branch, $D_F$ is the unique
source-minimal self-adjoint grading-odd full-cell generator.  Its typed
contractions onto the modules in \eqref{eq:SM-types} give the finite Yukawa and
optional Majorana residue blocks in one common generation frame.
\end{proposition}

\begin{proof}
Source saturation makes the full-cell carrier reducing.  The grading
decomposition is orthogonal, so the odd component
\eqref{eq:finite-dirac} is fixed by the represented generator.  Vanishing of
the orthogonal odd innovation removes any additional grading-odd source.
Typed compression is then a deterministic operation on the same operator.
\end{proof}

The weak multiplicity carrier $W_2$ supplies one primitive doublet.  On the
structural branch we require the primitive one-leg odd sector to be
source complete; any extra weak-doublet copy would then produce a nonzero
orthogonal source innovation and is therefore excluded by the same
minimality criterion.

\begin{theorem}[Structural Standard-Model classification from a landed multiplicity packet]
\label{thm:structural-SM}
Let a finite source-complete joint packet satisfy the positive geometry short
of Theorem~\ref{thm:geometry-short}.  Then the following implications hold on
the same represented carrier.
\begin{enumerate}[label=\textup{(\roman*)},leftmargin=2.8em]
\item If one finite word depth satisfies the multiplicity census and internal
      landing hypotheses of Theorem~\ref{thm:internal-seed-saturation}, then
      the complete internal relative commutant is
      \[
      \boxed{M_3(\C)\oplus M_2(\C)\oplus\C\oplus\C.}
      \]
\item On that landed multiplicity packet, if the canonical fully alternating
      incidence projection of Proposition~\ref{prop:main-det-incidence} is
      nonzero, then the physical determinant-preserving unitary group is
      \[
      \boxed{
      G_{\rm SM}
      =S(\U(3)\times\U(2))
      \cong
      \frac{\SU(3)\times\SU(2)\times\U(1)}{\Z_6}.}
      \]
\item For the represented tensor types in \eqref{eq:SM-types}, local anomaly
      freedom is equivalent to $3a+2b=0$ and hence fixes the Standard-Model
      hypercharge ratios \eqref{eq:hypercharges} up to overall orientation.
\item If the shorted endpoint Gram on $W_-$ is strictly positive, the internal
      endpoint source contains exactly one supported rank-three generation
      factor $G_{\rm gen}\cong\C^3$, while the orthogonal sign-twisted triplet
      remains loop provenance.
\item If a represented grading-changing mixed word has $p_{\rm Cl}>0$ and the
      source-saturated full-cell generator is odd-provenance complete, its odd
      part is one source-minimal finite Dirac operator $D_F$; the
      Yukawa/Majorana residue matrices are typed contractions of that operator
      in one common generation frame.
\item If the primitive weak one-leg odd sector is source complete, the weak
      sector contains exactly one primitive doublet; any additional persistent
      copy produces a nonzero orthogonal source innovation.
\end{enumerate}
When all six clauses hold, the finite packet has the complete structural
Standard-Model content claimed here.  The nonabelian landing itself is not a
dimension match: Theorems~\ref{thm:degree-one-colour-no-go} and
\ref{thm:flat-depth-colour} give explicit no-go branches, while
Theorem~\ref{thm:internal-deficit-alternatives} gives the exact smaller
internal algebra on every bridge/router failure.
\end{theorem}

\begin{proof}
Item (i) is Theorem~\ref{thm:internal-seed-saturation}.  Item (ii) follows
from Proposition~\ref{prop:main-det-incidence} and
Theorem~\ref{thm:gauge-group}.  Item (iii) is
Theorem~\ref{thm:hypercharge}; item (iv) is
Proposition~\ref{prop:generation-carrier}; item (v) is
Propositions~\ref{prop:clifford-occurrence} and~\ref{prop:finite-dirac}; and
item (vi) is the source-minimality criterion for the primitive weak one-leg
sector.
\end{proof}

\begin{corollary}[Historical typed-compilation criterion]
\label{cor:historical-typed-compilation}
Let the relevant historical word hierarchy be flat.  Suppose the grading,
type, weak-leg and incidence source/target projections used to form the
Store--Clifford and matter-incidence shadows are represented elements, or
finite functional-calculus elements, of that same historical $*$-algebra, and
that every router or finite symmetry average entering the compression is
word native in the sense of Definition~\ref{def:word-native-shadow}.  Then the
complete mixed ancestry Gram of those typed shadows, its Schur ancestry
residual, and the alternating determinant shadow are deterministic functions
of the existing Grand-Tensor word Gram and multiplication table.  No
independent mixed cross-Hankel source row is required.

Conversely, if a required typing projector or router is external to the
historical algebra, Theorem~\ref{thm:external-shadow-no-go} prevents its
historical occurrence from being inferred from the old word moments alone.
\end{corollary}

\begin{proof}
The typed compressions are word-native linear shadows, so
Theorem~\ref{thm:inserted-word-moments} and
Corollary~\ref{cor:table-native-ancestry} reconstruct all ordinary--shadow and
shadow--shadow blocks.  Complete alternation on the already reconstructed typed
Hom space is fixed linear postprocessing, so it introduces no additional
physical source datum.  The converse is exactly
Theorem~\ref{thm:external-shadow-no-go}.
\end{proof}

\begin{remark}[Classification, obstruction, and selection]
Theorem~\ref{thm:structural-SM} does not project arbitrary data onto the
Standard Model.  Before the determinant or anomaly tests are even applied,
the internal algebra is determined by a terminating multiplicity-side
alternative.  If the word hierarchy becomes flat with trivial multiplicity at
most two, internal colour is impossible.  If a three-dimensional colour-eligible
multiplicity carrier occurs but the bridge vanishes, the colour block remains
$M_2\oplus\C$.  If the weak router is absent, the weak multiplicity algebra
remains $\C^2$.  The four exact dimensions $9,11,13,15$ distinguish all such
cases.  On a fixed landed multiplicity census, full router saturation is
generic by Corollary~\ref{cor:generic-router-saturation}; microscopic occurrence
of the required word rows remains a separate source-identification question.
\end{remark}

\begin{remark}[Occurrence certification versus microscopic source identification]
\label{rem:occurrence-vs-identification}
The structural conclusions are finite system-bearing occurrence statements,
not consequences of a classical opportunity record, an environment effect
algebra, or a matching representation dimension.  For a particular
microscopic realization, the audit is finite and ordered: identify the actual
reachable carrier and apply the module-recognition test when a regular word
model is used; stop with Theorem~\ref{thm:flat-depth-colour} if a colour-eligible
multiplicity never appears; test either a direct bridge or the reciprocal-return
mass; verify the weak router and central separation; and only then apply the
determinant and grading-changing incidence tests.  When the required typing
operations are historical and word native, Corollary~\ref{cor:historical-typed-compilation}
shows that their mixed ancestry data are already compiled by the Grand Tensor.
The unresolved provenance question is then the occurrence of those typing
operations themselves.  The coarse-data and external-shadow no-go theorems
explain why absent operations cannot be manufactured from record marginals or
old word moments, while Theorem~\ref{thm:pullback-reconstruction} gives an
exact route from a complete operator-valued pullback to the represented action.
\end{remark}

\begin{scope}[Structural, not numerical]
\label{scope:structural}
Theorem~\ref{thm:structural-SM} does not determine the numerical entries of
$D_F$, the nonzero pattern of every allowed residue, the flavour hierarchy,
CKM or PMNS parameters, gauge couplings, the Higgs self-coupling or vacuum
scale, renormalization-group running, matching, or the fermionic measure
phase.  These are coordinates on, or dynamical data over, the reconstructed
structural packet rather than consequences of the commutant theorem.
\end{scope}

% =============================================================================

\medskip
The structural realization above identifies the desired finite multiplicity
packet.  We now give a basis-free positive test for exactness of the commutant
and then show that the same duality is stable along controlled regulator
refinement.

\subsection{Finite certificates and regulator-stable duality}
\label{sec:certificate}
% =============================================================================

An exact finite commutant identity is useful only if it can be recognized
without guessing a basis and if it remains stable under refinement.  Both
requirements are controlled by one positive operator.

\subsubsection{The commutant Laplacian}

\begin{definition}[Commutant Laplacian]
\label{def:commutant-laplacian}
Let $c_1,\ldots,c_s$ span a unital $*$-algebra
$\mathcal C\subseteq\mathcal B(\Hh)$.  On Hilbert--Schmidt operator space set
\begin{equation}
\boxed{
\mathscr L_{\mathcal C}
=
\sum_{j=1}^{s}
\operatorname{ad}_{c_j}^*\operatorname{ad}_{c_j},
\qquad
\operatorname{ad}_{c}(X)=[c,X].}
\label{eq:commutant-laplacian}
\end{equation}
\end{definition}

\begin{theorem}[Terminating commutant audit]
\label{thm:commutant-audit}
The operator $\mathscr L_{\mathcal C}$ is positive and
\begin{equation}
\boxed{\Ker\mathscr L_{\mathcal C}=\mathcal C'.}
\label{eq:commutant-kernel}
\end{equation}
If $P_{\mathcal C'}$ is Hilbert--Schmidt projection onto the commutant and
$\lambda_{\mathcal C}>0$ is the least positive eigenvalue, then
\begin{equation}
\boxed{
\norm{X-P_{\mathcal C'}X}_{\HS}^2
\le
\lambda_{\mathcal C}^{-1}
\sum_j\norm{[c_j,X]}_{\HS}^2.}
\label{eq:commutant-coercivity}
\end{equation}
\end{theorem}

\begin{proof}
For the Hilbert--Schmidt product,
$\operatorname{ad}_c^*=\operatorname{ad}_{c^*}$, hence
\[
\ip{X}{\mathscr L_{\mathcal C}X}_{\HS}
=
\sum_j\norm{[c_j,X]}_{\HS}^2.
\]
The quadratic form vanishes exactly on simultaneous commutants.  The positive
spectral gap on $(\mathcal C')^\perp$ gives
\eqref{eq:commutant-coercivity}.
\end{proof}

\subsubsection{The relative Howe Gram}

\begin{definition}[Relative Howe Gram]
\label{def:howe-gram}
Let $\M\subseteq C^*(\bm c)'$ be a proposed protected multiplicity algebra and
let $E_1,\ldots,E_r$ be a Hilbert--Schmidt orthonormal basis of
$\M^\perp$.  Define
\begin{equation}
\boxed{
\mathbb G_{\rm Howe}(\bm c\mid\M)_{ab}
=
\sum_j
\left\langle[c_j,E_a],[c_j,E_b]\right\rangle_{\HS}.}
\label{eq:howe-gram}
\end{equation}
\end{definition}

\begin{theorem}[Positive certificate for the exact multiplicity algebra]
\label{thm:howe-certificate}
Under Definition~\ref{def:howe-gram},
\begin{equation}
\boxed{
\mathbb G_{\rm Howe}(\bm c\mid\M)\succ0
\quad\Longleftrightarrow\quad
C^*(\bm c)'=\M.}
\label{eq:howe-certificate}
\end{equation}
On this branch the least eigenvalue of the Howe Gram is the first positive
eigenvalue of the commutant Laplacian.
\end{theorem}

\begin{proof}
The matrix in \eqref{eq:howe-gram} is the restriction of
$\mathscr L_{C^*(\bm c)}$ to $\M^\perp$.  Since $\M$ is already contained in
the kernel, the restriction is positive definite exactly when the kernel has
no additional vector.
\end{proof}

Thus the proposed generation/flavour algebra is not accepted by dimensional
matching.  It is certified by one positive finite matrix whose null directions
are exactly unaccounted commutants.


\subsubsection{Howe discriminant and generic exactness}
\label{sec:howe-discriminant}

The positive certificate also determines the symmetry-enhancement locus of a
parameter family.  Let $\Theta\subset\R^p$ be an open parameter domain, keep a
fixed represented protected algebra $\M\subseteq\mathcal B(\Hh)$, and assume
$[c_j(\theta),\M]=0$ for every $j$ and $\theta$.  On
$\Kk:=\M^\perp$ define the stacked commutator map
\[
 D(\theta)X=([c_1(\theta),X],\ldots,[c_s(\theta),X]),
 \qquad G(\theta)=D(\theta)^*D(\theta).
\]
In an orthonormal basis of $\Kk$, $G(\theta)$ is the relative Howe Gram.
Put $r=\dim_\C\Kk$ and
\begin{equation}
 \Delta_{\M}(\theta)=\det G(\theta).
 \label{eq:howe-discriminant}
\end{equation}

\begin{theorem}[Howe enhancement discriminant and generic rigidity]
\label{thm:howe-discriminant}
Assume every matrix entry of every $c_j(\theta)$ is a complex polynomial in
the real coordinates $\theta$, of degree at most $d$.  Then
$\Delta_{\M}$ is a nonnegative real polynomial of degree at most $2dr$ and
admits the Cauchy--Binet representation
\begin{equation}
 \boxed{
 \Delta_{\M}(\theta)
 =\sum_{|I|=r}|\det D_I(\theta)|^2.}
 \label{eq:howe-minor-sos}
\end{equation}
Moreover,
\begin{equation}
 \boxed{
 \{\theta:\Delta_{\M}(\theta)=0\}
 =\{\theta:C^*(\bm c(\theta))'\supsetneq\M\}.}
 \label{eq:howe-enhancement-locus}
\end{equation}
If the proposed commutant is exact at one parameter value, then
$\Delta_{\M}$ is not the zero polynomial.  Hence its enhancement locus has
empty interior and Lebesgue measure zero, while exact duality holds on an open
dense full-measure subset of every open admissible parameter region.  The same
open-dense conclusion holds for a real-analytic family on each connected
analytic component containing one exact point.
\end{theorem}

\begin{proof}
The entries of $D(\theta)$ are polynomial of degree at most $d$.
Cauchy--Binet applied to $D^*D$ gives
\eqref{eq:howe-minor-sos}; each squared modulus is a sum of squares of real
polynomials, proving nonnegativity and the degree bound.  Since $G\succeq0$,
its determinant is positive exactly when $G\succ0$, which is equivalent to
$C^*(\bm c(\theta))'=\M$ by Theorem~\ref{thm:howe-certificate}.  A nonzero
real polynomial has a zero set of Lebesgue measure zero and empty interior;
continuity makes the complement open.  For a connected real-analytic family,
the identity theorem gives the corresponding empty-interior conclusion once
the determinant is nonzero at one point.
\end{proof}

\begin{corollary}[Determinantal nullity filtration]
\label{cor:howe-nullity}
For $1\le k\le r$, the locus on which the commutant acquires at least $k$
additional complex dimensions is the common zero set of all
$(r-k+1)\times(r-k+1)$ minors of $D(\theta)$.  Additional commutant dimension
is therefore upper semicontinuous in every continuous family.
\end{corollary}

\begin{remark}[Genericity is relative to the admissible source class]
Theorem~\ref{thm:howe-discriminant} does not assert that microscopic Renewal
dynamics generically enters a chosen structural Standard-Model family.  It
states that, once a fixed protected algebra and an admissible polynomial or
analytic source family are specified, accidental enlargement of the
commutant is confined to the corresponding enhancement locus.
\end{remark}

\subsubsection{A finite anchor locks the commutant}

\begin{theorem}[Finite-anchor stability]
\label{thm:finite-anchor}
Transport a cutoff family to one common finite screened carrier and write
$\bm c_X=(c_{1,X},\ldots,c_{s,X})$.  Suppose a fixed protected algebra $\M$
commutes with every tuple.  Set
\[
C_X=\left(\sum_j\norm{c_{j,X}}_{\rm op}^2\right)^{1/2},
\qquad
\varepsilon_{X,Y}
=
\left(\sum_j\norm{c_{j,X}-c_{j,Y}}_{\rm op}^2\right)^{1/2}.
\]
If at one cutoff $X_0$
\[
\mathbb G_{\rm Howe}(\bm c_{X_0}\mid\M)
\succeq\gamma_0I,
\qquad \gamma_0>0,
\]
and for every later $X$
\begin{equation}
4(C_X+C_{X_0})\varepsilon_{X,X_0}<\gamma_0,
\label{eq:anchor-tail}
\end{equation}
then
\begin{equation}
\boxed{
C^*(\bm c_X)'=\M,
\qquad
\lambda_{\min}^+(\mathscr L_{C^*(\bm c_X)})
\ge
\gamma_0-4(C_X+C_{X_0})\varepsilon_{X,X_0}.}
\label{eq:anchor-lock}
\end{equation}
\end{theorem}

\begin{proof}
The commutator derivations satisfy
\[
\norm{\partial_{\bm c_X}-\partial_{\bm c_Y}}_{2\to2}
\le2\varepsilon_{X,Y},
\qquad
\norm{\partial_{\bm c_X}}_{2\to2}\le2C_X.
\]
Using
$A^*A-B^*B=A^*(A-B)+(A^*-B^*)B$ gives
\[
\norm{\mathscr L_X-\mathscr L_Y}
\le
4(C_X+C_Y)\varepsilon_{X,Y}.
\]
The anchor Laplacian is at least $\gamma_0I$ on $\M^\perp$.
The perturbation estimate gives \eqref{eq:anchor-lock}.
\end{proof}

\subsubsection{Compact-Mosco cofinal duality}
\label{sec:compact-mosco}

The perturbative anchor theorem controls one finite neighborhood in generator
norm.  A distinct nonperturbative mechanism is available for a genuine cofinal
family.  We formulate it on transported screened Hilbert spaces and derive
norm-resolvent and spectral convergence from the form hypotheses rather than
assuming them.  This is the compact form of the Mosco convergence mechanism;
see Kuwae--Shioya \cite{KuwaeShioya2003}.  The specialized argument required
here is recorded in Appendix~\ref{app:compact-mosco}.

Let $\mathscr H$ be a separable Hilbert space and let
$J_X:\mathscr H_X\to\mathscr H$ be isometric embeddings of the screened
Hilbert--Schmidt operator spaces, with $P_X=J_XJ_X^*$ the corresponding
orthogonal projections.  The limiting space is embedded by $J_\infty$, and we
assume the transport is compatible in the sense that
$P_X\to P_\infty$ strongly on $\mathscr H$.  For a densely defined closed
nonnegative form $q_X$ on $\mathscr H_X$, denote
its self-adjoint operator by $\mathscr L_X$ and its transported resolvent by
\begin{equation}
R_X:=J_X(\mathscr L_X+I)^{-1}J_X^*\in\mathcal B(\mathscr H).
\label{eq:transported-resolvent}
\end{equation}

\begin{definition}[Transported Mosco convergence]
\label{def:transported-mosco}
The forms $q_X$ converge to $q_\infty$ in the transported Mosco sense when:
\begin{enumerate}[label=\textup{(M\arabic*)},leftmargin=3em]
\item whenever $J_XA_X\rightharpoonup B$ weakly in $\mathscr H$, the strong
      convergence $P_X\to P_\infty$ gives $B\in J_\infty\mathscr H_\infty$;
      writing $B=J_\infty A$,
\[
 q_\infty(A)\le\liminf_X q_X(A_X);
\]
\item for every $A\in D(q_\infty)$ there are $A_X\in D(q_X)$ such that
$J_XA_X\to J_\infty A$ strongly and
\[
 q_X(A_X)\longrightarrow q_\infty(A).
\]
\end{enumerate}
\end{definition}

\begin{definition}[Asymptotic form compactness]
\label{def:asymptotic-compactness}
The family $(q_X)$ is asymptotically compact if every sequence
$A_X\in D(q_X)$ satisfying
\begin{equation}
\sup_X\bigl(\norm{A_X}^2+q_X(A_X)\bigr)<\infty
\label{eq:asymptotic-compactness}
\end{equation}
has a subsequence for which $J_XA_X$ converges strongly in $\mathscr H$.
The same requirement is imposed on the limiting form.  In particular,
$(\mathscr L_\infty+I)^{-1}$ is compact on the limiting screened space.
\end{definition}

\begin{proposition}[Compact Mosco convergence implies norm-resolvent convergence]
\label{prop:compact-mosco-resolvent}
If $q_X\to q_\infty$ in the sense of
Definition~\ref{def:transported-mosco} and the family is asymptotically
compact in the sense of Definition~\ref{def:asymptotic-compactness}, then
\begin{equation}
\boxed{\norm{R_X-R_\infty}_{\rm op}\longrightarrow0.}
\label{eq:norm-resolvent}
\end{equation}
Consequently, if
$0\le\lambda_{1,X}\le\lambda_{2,X}\le\cdots$ denote the discrete screened
eigenvalues of $\mathscr L_X$, repeated with multiplicity, then for every
fixed index in the limiting discrete spectrum,
\begin{equation}
\boxed{\lambda_{k,X}\longrightarrow\lambda_{k,\infty}.}
\label{eq:eigenvalue-convergence}
\end{equation}
If an isolated spectral cluster of $\mathscr L_\infty$ is separated by a
positive gap, the corresponding transported Riesz projections converge in
operator norm and their ranks are eventually constant.
\end{proposition}

\begin{proof}
Transported Mosco convergence gives strong convergence of the resolvents.
The resolvent equation bounds
$\norm{u_X}^2+q_X(u_X)$ uniformly for
$u_X=(\mathscr L_X+I)^{-1}f_X$ whenever $(f_X)$ is bounded.  Asymptotic form
compactness therefore makes the transported resolvents collectively compact.
Because they are self-adjoint, collective compactness together with strong
convergence upgrades the convergence to operator norm.  Norm convergence of
compact self-adjoint resolvents gives \eqref{eq:eigenvalue-convergence}; the
Riesz-projection statement follows from the contour resolvent identity.
A detailed proof, including the varying-space bookkeeping, is given in
Appendix~\ref{app:compact-mosco}.
\end{proof}

\begin{theorem}[Compact-Mosco cofinal spacetime--gauge duality]
\label{thm:cofinal-duality}
Let
\begin{equation}
q_X(A)=\sum_j\norm{[c_{j,X},A]}_{\HS}^2
\label{eq:cofinal-form}
\end{equation}
be the closed commutator forms of a cofinal family of screened support-polar
or multiplicity-quiver generators, with associated commutant Laplacians
$\mathscr L_X$.  Assume:
\begin{enumerate}[label=\textup{(H\arabic*)},leftmargin=3em]
\item $q_X\to q_\infty$ in transported Mosco sense;
\item $(q_X)$ is asymptotically form compact;
\item for every $X$ there is a protected multiplicity algebra
      $\M_X\subseteq\Ker\mathscr L_X$ of the same finite complex dimension $m$;
\item the limiting commutant contains no unprotected direction,
      $\Ker\mathscr L_\infty=\M_\infty$ and
      $\dim_\C\M_\infty=m$.
\end{enumerate}
Then the limiting operator has compact resolvent and
\begin{equation}
\gamma_\infty
:=\lambda_{m+1}(\mathscr L_\infty)>0.
\label{eq:limit-gap}
\end{equation}
For every sufficiently late cutoff,
\begin{equation}
\boxed{
\Ker\mathscr L_X=\M_X,
\qquad
C^*(\bm c_X)'=\M_X,}
\label{eq:cofinal-kernel}
\end{equation}
and
\begin{equation}
\boxed{
\lambda_{\min}^+(\mathscr L_X)
=\lambda_{m+1}(\mathscr L_X)
\longrightarrow\gamma_\infty.}
\label{eq:cofinal-gap}
\end{equation}
If $\Pi_X$ denotes the transported spectral projection onto
$\Ker\mathscr L_X$, then
\begin{equation}
\boxed{\norm{\Pi_X-\Pi_\infty}_{\rm op}\longrightarrow0.}
\label{eq:kernel-projection-convergence}
\end{equation}
Consequently, for every $0<\gamma_*<\gamma_\infty$ and all sufficiently late
$X$,
\begin{equation}
\boxed{
\sum_j\norm{[c_{j,X},A]}_{\HS}^2
\ge
\gamma_*\norm{(I-P_{\M_X})A}_{\HS}^2.}
\label{eq:cofinal-coercivity}
\end{equation}
\end{theorem}

\begin{proof}
Proposition~\ref{prop:compact-mosco-resolvent} gives norm convergence of the
transported resolvents and hence convergence of every fixed low eigenvalue and
of the isolated zero-cluster projection.  Compactness makes the limiting
spectrum discrete near zero.  By \textup{(H4)}, zero has multiplicity exactly
$m$, so \eqref{eq:limit-gap} holds.  By \textup{(H3)}, every finite cutoff has
at least $m$ exact zero modes.  Norm convergence of the Riesz projection onto
the isolated eigenvalue $1$ of $(\mathscr L_\infty+I)^{-1}$ implies that the
corresponding finite-cutoff spectral subspace has dimension exactly $m$ for
all sufficiently late $X$.  Hence no additional zero mode is possible and
$\Ker\mathscr L_X=\M_X$.  Eigenvalue convergence gives
\eqref{eq:cofinal-gap}, projection convergence gives
\eqref{eq:kernel-projection-convergence}, and the spectral theorem gives
\eqref{eq:cofinal-coercivity}.
\end{proof}

\begin{remark}[Two distinct stability mechanisms]
Theorems~\ref{thm:finite-anchor} and~\ref{thm:cofinal-duality} are logically
independent.  The former is a quantitative perturbative statement in generator
norm around one positive finite anchor.  The latter is nonperturbative: it
allows the carrier and form domains to vary and derives spectral stability from
Mosco convergence plus compactness.  Either mechanism can lock the same finite
commutant; neither is a substitute for occurrence of the underlying
system-bearing generators.
\end{remark}

\begin{remark}[Independence of the screen construction]
Theorem~\ref{thm:cofinal-duality} depends only on the transported form
hypotheses \textup{(H1)}--\textup{(H4)}.  Any regulator family satisfying them
has the stated stable commutant, regardless of the microscopic origin of its
compactness.  Ref.~\cite{PelissierSpectral2026} supplies one Renewal-Geometric
mechanism for the required screened form control.
\end{remark}

\begin{remark}[What the cofinal theorem does and does not say]
Theorem~\ref{thm:cofinal-duality} stabilizes the algebraic duality and the
protected multiplicity spectral subspace.  It does not by itself construct an
interacting four-dimensional quantum field theory, determine running
couplings, remove fermionic determinant-line phases, or construct a continuum
measure.
\end{remark}

% =============================================================================

\medskip
The commutant duality and its coercive stability are algebraic statements.  We
finally place the same internal sector on the common variational carrier of the
published spacetime construction and isolate the source-minimal gravity--matter
coupling.

\subsection{The common gravity--internal action}
\label{sec:coupled-action}
% =============================================================================

The duality theorem is algebraic.  We now formulate the source-minimal
variational coupling abstractly on a common gravitational/internal carrier.
The spacetime companion paper \cite{PelissierSpacetime2026} supplies the
Renewal-Geometric realization of the gravitational side of this common finite
carrier.  The purpose of the subsection is to establish compatibility of the
dual pair with one gravity--matter action, not to reproduce the separate
functional-analytic problem of constructing its full continuum limit.

\subsubsection{The internal source after gravitational shorting}

Let
\[
S_{\rm g}:E_{\rm g}\to\Hh_{\rm cyl},
\qquad
S_{\rm i}:E_{\rm i}\to\Hh_{\rm cyl}
\]
be the source syntheses of the complete gravitational and internal physical
variation banks on one same-history cylinder.  Assume the gravitational bank
has already been reduced to its faithful quotient and put
\[
G_{\rm g}=S_{\rm g}^*S_{\rm g},
\qquad
G_{\rm i}=S_{\rm i}^*S_{\rm i},
\qquad
C=S_{\rm i}^*S_{\rm g}.
\]

\begin{theorem}[Universal coupled source quotient]
\label{thm:coupled-schur}
Define
\begin{equation}
\boxed{
G_{{\rm i}\mid{\rm g}}
=
G_{\rm i}-CG_{\rm g}^{-1}C^*
=
S_{\rm i}^*(I-P_{\rm g})S_{\rm i}.}
\label{eq:coupled-schur}
\end{equation}
Then
\begin{equation}
\boxed{
\begin{pmatrix}
G_{\rm g}&C^*\\
C&G_{\rm i}
\end{pmatrix}
=
\begin{pmatrix}
I&0\\
CG_{\rm g}^{-1}&I
\end{pmatrix}
\begin{pmatrix}
G_{\rm g}&0\\
0&G_{{\rm i}\mid{\rm g}}
\end{pmatrix}
\begin{pmatrix}
I&G_{\rm g}^{-1}C^*\\
0&I
\end{pmatrix}.}
\label{eq:coupled-factorization}
\end{equation}
Consequently:
\begin{enumerate}[label=\textup{(\roman*)},leftmargin=2.8em]
\item the joint variation carrier is faithful exactly when
      $G_{{\rm i}\mid{\rm g}}\succ0$;
\item the number of genuinely new internal source directions is
      $\rank G_{{\rm i}\mid{\rm g}}$;
\item the source-minimal internal variation space is
      \[
      E_{{\rm i}\mid{\rm g}}^{\min}
      =
      E_{\rm i}/\Ker G_{{\rm i}\mid{\rm g}}.
      \]
\end{enumerate}
\end{theorem}

\begin{proof}
The first equality in \eqref{eq:coupled-schur} is the Gram of
$(I-P_{\rm g})S_{\rm i}$.  Direct block multiplication gives
\eqref{eq:coupled-factorization}; the remaining statements are the standard
Schur-complement rank and positivity consequences.
\end{proof}

Thus ``adding the Standard Model'' does not mean appending every internal
variation as a new primitive.  Only the component orthogonal to the already
represented gravitational source counts as genuinely new.

\subsubsection{Finite common-action representative}

The duality theorem does not require a gauge--matter variational action.
For completeness, Appendix~\ref{app:finite-common-action} constructs one
explicit finite gauge--Higgs--Dirac representative on the reconstructed
spacetime complex and proves its exact gauge, Ward, BRST, stress, and stationary
vacuum identities.  Those formulas provide a concrete realization for the
compatibility statement below, but they are not used in the proof of the
finite double-centralizer theorem or in the structural Standard-Model
reconstruction.

\subsubsection{Classical compatibility of the stabilized dual branch}
\label{sec:classical-compatibility}

The finite common-action construction is sufficient for the purpose of the
present paper.  To state precisely what it implies, let
$\mathscr S_{{\rm SM},X}^{L}$ denote the Lorentzian signed counterpart of
\eqref{eq:finite-SM-action}, obtained from the reconstructed temporal Hodge
signs.  We do not attempt here to prove the independent geometric and PDE
estimates required to construct a general Einstein--Standard-Model continuum
from microscopic renewal data.  Instead we record the exact compatibility
implication that any such certified continuum realization must satisfy.

\begin{proposition}[Classical compatibility of the stabilized dual branch]
\label{prop:classical-compatibility}
Let a cofinal family satisfy the coercive multiplicity duality of
Theorem~\ref{thm:cofinal-duality}.  Suppose, in addition, that the gravitational
and internal fields are transported to one common determining test core such
that:
\begin{enumerate}[label=\textup{(C\arabic*)},leftmargin=3em]
\item the gravitational first-variation functionals converge to the
      Einstein--Hilbert/Palatini first variation with cosmological term;
\item the first variations of $\mathscr S_{{\rm SM},X}^{L}$ converge to the
      classical Yang--Mills--Higgs--Dirac--Yukawa first variation, including
      the metric variation defining $T_{\mu\nu}^{\rm SM}$;
\item for a determining transported gravity--matter variation bank,
      \[
      D\!\left(\mathscr S_{{\rm g},X}
      +\mathscr S_{{\rm SM},X}^{L}\right)[\dot z_X]\longrightarrow0.
      \]
\end{enumerate}
Then every limiting field configuration satisfies the classical matter Euler
equations on that test core and
\begin{equation}
\boxed{
G_{\mu\nu}+\Lambda g_{\mu\nu}
=\kappa\,T_{\mu\nu}^{\rm SM}}
\label{eq:einstein-matter}
\end{equation}
distributionally, with $\kappa$ fixed by the relative normalization of the
gravity and matter terms in the common action.  Along the same cofinal family,
the limiting matter fields carry the stabilized multiplicity algebra of
Theorem~\ref{thm:cofinal-duality}; no additional unprotected
flavour/generation direction is created by the handoff.
\end{proposition}

\begin{proof}
Take first a pure compactly supported matter variation in the determining
bank.  By hypothesis~\textup{(C2)}, the finite matter first variations converge
to the classical Yang--Mills--Higgs--Dirac--Yukawa first variation, whereas the
gravity variation vanishes.  Passing to the limit in~\textup{(C3)} therefore
gives the matter Euler equations.

Next take a pure compactly supported metric variation $k^{\mu\nu}$.  With
fixed nonzero normalizations $c_{\rm g}$ and $c_{\rm m}$ for the gravitational
and matter parts, hypotheses~\textup{(C1)}--\textup{(C2)} give
\[
 c_{\rm g}\!\int
 (G_{\mu\nu}+\Lambda g_{\mu\nu})k^{\mu\nu}\,d{\rm vol}_g
 -\frac{c_{\rm m}}2\!\int
 T_{\mu\nu}^{\rm SM}k^{\mu\nu}\,d{\rm vol}_g=0,
\]
up to the fixed stress convention used in
Theorem~\ref{thm:finite-action-identities}.  Absorbing that convention into
$\kappa$ and using that the metric test bank is determining yields
\eqref{eq:einstein-matter} as a distributional identity.

Finally, Theorem~\ref{thm:cofinal-duality} identifies the multiplicity kernel
for all sufficiently late cutoffs and gives a uniform positive commutant gap.
Hence the matter packet entering the limiting action is the same stabilized
structural packet reconstructed by the finite duality theorem, rather than a
larger multiplicity sector generated by refinement.
\end{proof}

\begin{scope}[Continuum analysis is separate from the duality theorem]
Proposition~\ref{prop:classical-compatibility} is a compatibility statement,
not a construction of a general cofinal Einstein--Standard-Model solution.
Establishing the hypotheses above from microscopic renewal dynamics requires a
separate analytic theory of geometric reconstruction, nonlinear stress
compactness, constraint propagation, boundary/domain control, and strong
coupled-source convergence.  Those problems are logically downstream of the
finite commutant duality and are not used anywhere in its proof or in the
structural Standard-Model reconstruction.
\end{scope}

% =============================================================================
\section{Discussion and outlook}
\label{sec:discussion}
% =============================================================================

The main result fixes the precise meaning of ``spacetime--gauge duality'' in
this framework.  The duality is not an identification of the Lorentzian
Clifford algebra with the Standard-Model gauge algebra, nor an assertion that
one is the ordinary commutant of the other.  The factor obstruction shows that
such a statement is algebraically impossible on the relevant reducible typed
carrier.  The exact dual pair is instead formed by the \emph{joint typed
action}---spacetime, internal representation, and physical incidence
together---and the matter multiplicity that remains invisible to that complete
action.

This distinction matters physically.  A gauge representation tells us how a
field transforms, but it does not by itself determine how many equivalent
copies of that representation occur.  In the present construction that
multiplicity is not an unrelated label added after the gauge sector has been
chosen.  It is the commutant of the complete represented action.  Conversely,
once the protected multiplicity algebra is known, finite bicommutant closure
recovers the joint action algebra.  The two sides are therefore complementary
descriptions of one represented finite packet.

Corollary~\ref{cor:reciprocal-wedderburn} makes this complementarity explicit.
The common carrier splits into action--multiplicity tensor blocks; inside each
central sector the joint action is a full matrix algebra on one factor and the
residual multiplicity is a full matrix algebra on the other.  Both sides have
the same centre.  Thus the duality is stronger than equality of two abstract
dimensions: it is a reciprocal finite factorization of the represented packet,
canonical up to simultaneous unitary equivalence.

The singular theory is important for the same reason.  Yukawa or incidence
maps can have zero modes, unequal source and target dimensions, or
noncommuting supports.  None of these cases requires inserting an inverse
matrix or completing a missing edge by hand.  The polar support algebra is
defined directly from the occurring maps, and its commutant is exact.  When
the support projections themselves carry nontrivial geometry, the canonical
multiplicity quiver records it.  The surviving flavour algebra is then the
endomorphism algebra of that quiver.

The incidence bank itself has an intrinsic minimality theory.  The grouped
commutator constraints define the submodular rank function
\eqref{eq:incidence-rank}; deleting an edge is harmless exactly when that rank
is unchanged, and every finite exact packet contains an inclusion-minimal
incidence skeleton.  Thus ``physical incidence is essential'' is a terminating
rank statement rather than a graphical heuristic.

The source-complete requirement is likewise operational rather than
terminological.  Theorem~\ref{thm:record-nonidentifiability} shows that even a
complete classical record history and matching coarse Choi data can leave the
acted-on system algebra ambiguous.  Conversely, the operator-valued
complementary pullback of Theorem~\ref{thm:pullback-reconstruction} recovers
the system action exactly when an occurring identity pivot is present.  The
duality theorem is therefore unconditional once the system-bearing packet is
specified, while microscopic provenance is a distinct source-identification
problem.

The positive Howe Gram turns this structural statement into a finite
certificate.  Instead of asserting that a proposed generation or flavour
algebra is the full commutant, one computes a positive matrix on its
orthogonal complement.  A zero eigenvector is an explicit missing symmetry;
a positive least eigenvalue is a rigidity margin.  For polynomial or analytic
source families, Theorem~\ref{thm:howe-discriminant} sharpens this statement:
commutant enhancement lies on a determinantal locus and exact duality is
generic within every admissible component containing one exact point.  Cofinal
stability is likewise a conclusion of the operator theory: transported Mosco
convergence and asymptotic compactness yield norm-resolvent convergence,
convergence of the protected commutant projection, and a uniform late-cutoff
coercive gap.

The Standard-Model result should be read at the same structural level.  The
active residual does not obtain $M_3(\C)$ and $M_2(\C)$ by noticing matrix
blocks of dimensions three and two in the tetrahedral representation: those
direct Burnside blocks are external.  Internal factors live on multiplicity
spaces, and Theorem~\ref{thm:word-module-recognition} further distinguishes
kinematic invariant endomorphisms of a word space from operations induced by
the historical system algebra.  At the populated categorical degree-one
census internal colour is impossible; if the physical word hierarchy becomes
flat before an external-trivial multiplicity of three appears, colour is ruled
out at every later depth.  Once the multiplicity census $(3,2,1,1)$ is present,
the nonabelian landing requires only a colour bridge and a noncommuting weak
router; on a reciprocal common-mediator branch the bridge is itself forced by
the positive return mass \eqref{eq:reciprocal-return-mass}.  A separate central
separator rules out the only remaining $14$-dimensional scalar-locked
subdirect product before the full
$M_3(\C)\oplus M_2(\C)\oplus\C\oplus\C$ commutant is asserted.  Under
independent central supports the bridge/router deficits are the exact
$9,11,13,15$ alternatives already displayed.  Only after this internal
landing is the determinant incidence tested.  Its stabilizer is
$S(\U(3)\times\U(2))$; a nonzero coherent determinant cross mass gives the same
group operationally, while the sharp arity bound shows that this character
cannot be manufactured by lower-order defining tensor ports.  Anomaly freedom
then fixes the remaining central weight ratio.

This provides a useful contrast with the almost-commutative formulation of
noncommutative geometry.  There, the Standard Model is encoded by a finite
spectral geometry combined with a continuum spin manifold, and the spectral
action produces the coupled field theory
\cite{ConnesGravityMatter1996,ChamseddineConnes1997,
ChamseddineConnesMarcolli2007,ChamseddineConnes2008,
vanDenDungenVanSuijlekom2012}.  Krajewski classification and the later
Hodge/Morita programme make the finite incidence and commutant structures of
that input highly nontrivial \cite{Krajewski1998,DAndreaDabrowski2016,
DabrowskiDAndreaSitarz2018}.  Here the finite spectral packet and Lorentzian
spacetime are themselves downstream outputs of the same predictive Grand
Tensor.  The present duality is therefore not a replacement of the
almost-commutative product by another chosen product.  It is a statement about
the mutual commutants that remain after both downstream sectors are realized
on one common carrier.

The generation result is perhaps the most delicate structural consequence.
Three generations do not follow merely because the spacetime construction
contains a three-dimensional module.  The geometric copy is removed first.
What remains on the internal categorical contrast is a distinct
reversal-odd standard module of rank three.  Its orthogonal sign-twisted
triplet is kept as loop provenance.  Thus the finite source calculation
contains an exact mechanism that distinguishes the generation triplet from a
second geometric or cycle copy.  Whether the physical flavour matrices select
the observed hierarchy and mixing inside that triplet is a separate question.

The same separation applies to the finite Dirac operator.  The framework does
not derive five unrelated Yukawa matrices and then place them into a larger
operator.  On the source-saturated branch, one grading-odd full-cell generator
is reconstructed first; the residue matrices are typed contractions of that
one operator in one generation frame.  Representation theory determines the
allowed slots, not their numerical values.  Flavour breaking, masses, mixing
angles, and CP phases therefore remain genuine dynamical information.

The common-action Schur quotient adds a final compatibility point without
enlarging the duality theorem.  Internal variations count as genuinely new
source directions only through the residual
$G_{{\rm i}\mid{\rm g}}=S_{\rm i}^*(I-P_{\rm g})S_{\rm i}$.  This is the
variational analogue of the algebraic construction: in both cases the internal
sector is defined relative to the already represented spacetime source rather
than appended by direct sum.  The explicit finite gauge--matter action in
Appendix~\ref{app:finite-common-action} shows that the stabilized dual packet is
compatible with later coupled dynamics; no Einstein--Standard-Model continuum
limit is needed for the finite mutual-commutant result.

Several open problems remain.  The first is microscopic source identification
and dynamical selection, but the unresolved interface is narrower than a list
of independent mixed Gram entries.  Once the historical word algebra and its
multiplication table are known, every mixed ancestry quantity associated with
a word-native typing operation is already an inserted Grand-Tensor moment.
The source-level question is therefore whether the grading, type, incidence and
router operations used by the structural compiler themselves occur in the
historical flat algebra, and whether an actual determinant-bearing typed port
is present.  If so, the ancestry and determinant tests become finite table-native
computations; if not, Theorem~\ref{thm:external-shadow-no-go} shows that the
missing operation cannot be inferred from the old moments alone.  The present
theorems do not assert that generic microscopic dynamics selects this terminal
branch.  A second problem is parameter selection.  Nothing here fixes the observed gauge
couplings, Higgs parameters, Yukawa spectra, or mixing matrices.  Finally, the
full classical Einstein--Standard-Model continuum is analytically distinct:
deriving the transported Mosco and asymptotic-compactness hypotheses, together
with the convergence assumptions of
Proposition~\ref{prop:classical-compatibility}, from a single microscopic
cofinal family requires geometric reconstruction, nonlinear stress compactness,
constraint propagation, and coupled-source estimates that are not consequences
of the commutant theorem.  Beyond that lies the quantum problem, including the
determinant/Pfaffian phase, renormalization, and the construction of an
interacting four-dimensional continuum measure.

Within these limits, the logical research hierarchy is naturally staged.  The
spacetime paper reconstructs Lorentzian geometry and gravitational dynamics,
and the spectral paper reconstructs a Connes-type effective geometry and
characterizes information lost under spectralization.  The present paper uses
those works to supply the independently reconstructed Lorentzian and spectral
inputs, while its finite duality and structural branch are autonomous
consequences of the explicit joint-packet and occurrence hypotheses stated
here.  A separate continuum analysis may then
ask when the finite common gravity--matter action converges to the classical
Einstein--Standard-Model system.  The logical content of the present paper ends
one step earlier:
\begin{equation}
\boxed{
\begin{aligned}
\text{Renewal Grand Tensor}
&\longrightarrow
\text{Lorentzian + spectral common carrier}\\
&\longrightarrow
\text{joint typed action}\\
&\overset{\rm commutant}{\longleftrightarrow}
\text{matter multiplicity}\\
&\longrightarrow
\text{structural Standard Model}.
\end{aligned}}
\label{eq:discussion-chain}
\end{equation}
The central conclusion is therefore not that internal gauge symmetry is
spacetime in disguise.  Rather, spacetime action and matter multiplicity are
complementary operator-algebraic faces of one predictive carrier once the
physical incidence between them is retained.  The finite common-action construction shows that this dual
pair is compatible with later gravity--matter dynamics, while the analytic
continuum construction is deliberately left to a separate work.

% =============================================================================
\section*{Data and formalization availability}
% =============================================================================

The paper is mathematical and uses no external empirical dataset.  The finite
algebraic statements are compatible with the formalization programme at
\url{https://github.com/AI-MathPhys/renewal-geometry}.  Publication metadata
for the two companion Renewal Geometry papers should be updated in the
bibliography when final journal details are available.

% =============================================================================
\begin{thebibliography}{99}

\bibitem{PelissierRenewalSpacetime2026}
A.~P\'elissier,
\newblock Renewal Geometry and the emergence of Lorentzian spacetime,
\newblock HAL preprint hal-05744772, 2026.
\newblock \url{https://hal.science/hal-05744772}

\bibitem{PelissierSpectral2026}
A.~P\'elissier,
\newblock From predictive dynamics to spectral geometry,
\newblock 2026.
\newblock % Replace with final publication metadata.

\bibitem{ColemanMandula1967}
S.~Coleman and J.~Mandula,
\newblock All possible symmetries of the $S$ matrix,
\newblock \emph{Physical Review} \textbf{159} (1967), 1251--1256.
\newblock doi:10.1103/PhysRev.159.1251.

\bibitem{ConnesLott1991}
A.~Connes and J.~Lott,
\newblock Particle models and noncommutative geometry,
\newblock \emph{Nuclear Physics B -- Proceedings Supplements} \textbf{18} (1991),
29--47.
\newblock doi:10.1016/0920-5632(91)90120-4.

\bibitem{Connes1994}
A.~Connes,
\newblock \emph{Noncommutative Geometry},
\newblock Academic Press, San Diego, 1994.

\bibitem{ConnesReality1995}
A.~Connes,
\newblock Noncommutative geometry and reality,
\newblock \emph{Journal of Mathematical Physics} \textbf{36} (1995), no.~11,
6194--6231.
\newblock doi:10.1063/1.531241.

\bibitem{ConnesGravityMatter1996}
A.~Connes,
\newblock Gravity coupled with matter and the foundation of non-commutative
geometry,
\newblock \emph{Communications in Mathematical Physics} \textbf{182} (1996),
no.~1, 155--176.
\newblock doi:10.1007/BF02506388.

\bibitem{ChamseddineConnes1997}
A.~H.~Chamseddine and A.~Connes,
\newblock The spectral action principle,
\newblock \emph{Communications in Mathematical Physics} \textbf{186} (1997),
731--750.
\newblock doi:10.1007/s002200050126.

\bibitem{ChamseddineConnesMarcolli2007}
A.~H.~Chamseddine, A.~Connes, and M.~Marcolli,
\newblock Gravity and the standard model with neutrino mixing,
\newblock \emph{Advances in Theoretical and Mathematical Physics}
\textbf{11} (2007), no.~6, 991--1089.
\newblock doi:10.4310/ATMP.2007.v11.n6.a3.

\bibitem{ChamseddineConnes2008}
A.~H.~Chamseddine and A.~Connes,
\newblock Why the Standard Model,
\newblock \emph{Journal of Geometry and Physics} \textbf{58} (2008), no.~1,
38--47.
\newblock doi:10.1016/j.geomphys.2007.09.011.

\bibitem{vanSuijlekom2025}
W.~D.~van Suijlekom,
\newblock \emph{Noncommutative Geometry and Particle Physics},
\newblock 2nd ed., Mathematical Physics Studies, Springer, 2025.


\bibitem{Krajewski1998}
T.~Krajewski,
\newblock Classification of finite spectral triples,
\newblock \emph{Journal of Geometry and Physics} \textbf{28} (1998), 1--30.
\newblock doi:10.1016/S0393-0440(97)00068-5.

\bibitem{IochumSchuckerStephan2004}
B.~Iochum, T.~Sch\"ucker, and C.~Stephan,
\newblock On a classification of irreducible almost commutative geometries,
\newblock \emph{Journal of Mathematical Physics} \textbf{45} (2004),
5003--5041.
\newblock doi:10.1063/1.1811372.

\bibitem{vanDenDungenVanSuijlekom2012}
K.~van den Dungen and W.~D.~van Suijlekom,
\newblock Particle physics from almost commutative spacetimes,
\newblock \emph{Reviews in Mathematical Physics} \textbf{24} (2012),
1230004.
\newblock doi:10.1142/S0129055X1230004X.

\bibitem{ChamseddineConnesVanSuijlekom2013}
A.~H.~Chamseddine, A.~Connes, and W.~D.~van Suijlekom,
\newblock Beyond the spectral Standard Model: emergence of Pati--Salam
unification,
\newblock \emph{Journal of High Energy Physics} \textbf{2013} (2013),
no.~11, 132.
\newblock doi:10.1007/JHEP11(2013)132.

\bibitem{BoyleFarnsworth2018}
L.~Boyle and S.~Farnsworth,
\newblock A new algebraic structure in the standard model of particle physics,
\newblock \emph{Journal of High Energy Physics} \textbf{2018} (2018), no.~6, 071.
\newblock doi:10.1007/JHEP06(2018)071.

\bibitem{DAndreaDabrowski2016}
F.~D'Andrea and L.~Dabrowski,
\newblock The Standard Model in noncommutative geometry and Morita equivalence,
\newblock \emph{Journal of Noncommutative Geometry} \textbf{10} (2016),
551--578.
\newblock doi:10.4171/JNCG/242.

\bibitem{DabrowskiDAndreaSitarz2018}
L.~Dabrowski, F.~D'Andrea, and A.~Sitarz,
\newblock The Standard Model in noncommutative geometry: fundamental fermions
as internal forms,
\newblock \emph{Letters in Mathematical Physics} \textbf{108} (2018),
1323--1340.
\newblock doi:10.1007/s11005-017-1036-x.

\bibitem{DabrowskiSitarz2019}
L.~Dabrowski and A.~Sitarz,
\newblock Fermion masses, mass-mixing and the almost commutative geometry of
the Standard Model,
\newblock \emph{Journal of High Energy Physics} \textbf{2019} (2019),
no.~2, 68.
\newblock doi:10.1007/JHEP02(2019)068.

\bibitem{BhowmickEtAl2014}
J.~Bhowmick, F.~D'Andrea, B.~Das, and L.~Dabrowski,
\newblock Quantum gauge symmetries in noncommutative geometry,
\newblock \emph{Journal of Noncommutative Geometry} \textbf{8} (2014),
433--471.
\newblock doi:10.4171/JNCG/161.

\bibitem{Barrett2007}
J.~W.~Barrett,
\newblock A Lorentzian version of the non-commutative geometry of the standard
model of particle physics,
\newblock \emph{Journal of Mathematical Physics} \textbf{48} (2007), 012303.
\newblock doi:10.1063/1.2408400.

\bibitem{DevastatoEtAl2018}
A.~Devastato, S.~Farnsworth, F.~Lizzi, and P.~Martinetti,
\newblock Lorentz signature and twisted spectral triples,
\newblock \emph{Journal of High Energy Physics} \textbf{2018} (2018),
no.~3, 89.
\newblock doi:10.1007/JHEP03(2018)089.

\bibitem{TraylingBaylis2001}
G.~Trayling and W.~E.~Baylis,
\newblock A geometric basis for the standard-model gauge group,
\newblock \emph{Journal of Physics A: Mathematical and General} \textbf{34} (2001),
3309--3324.
\newblock doi:10.1088/0305-4470/34/15/309.

\bibitem{Daviau2017}
C.~Daviau,
\newblock Gauge group of the Standard Model in $Cl_{1,5}$,
\newblock \emph{Advances in Applied Clifford Algebras} \textbf{27} (2017), 279--290.
\newblock doi:10.1007/s00006-015-0566-5.

\bibitem{Gresnigt2020}
N.~G.~Gresnigt,
\newblock The Standard Model particle content with complete gauge symmetries from
minimal ideals of two Clifford algebras,
\newblock \emph{European Physical Journal C} \textbf{80} (2020), 583.
\newblock doi:10.1140/epjc/s10052-020-8141-1.

\bibitem{DoplicherRoberts1989}
S.~Doplicher and J.~E.~Roberts,
\newblock A new duality theory for compact groups,
\newblock \emph{Inventiones Mathematicae} \textbf{98} (1989), 157--218.
\newblock doi:10.1007/BF01388849.

\bibitem{DoplicherRoberts1990}
S.~Doplicher and J.~E.~Roberts,
\newblock Why there is a field algebra with a compact gauge group describing
the superselection structure in particle physics,
\newblock \emph{Communications in Mathematical Physics} \textbf{131} (1990),
51--107.
\newblock doi:10.1007/BF02097680.

\bibitem{Todorov2010}
I.~Todorov,
\newblock Gauge symmetry and Howe duality in 4D conformal field theory models,
\newblock \emph{Advances in Mathematical Physics} \textbf{2010} (2010),
Article ID 509538, 12 pages.
\newblock doi:10.1155/2010/509538.

\bibitem{Hucks1991}
J.~Hucks,
\newblock Global structure of the standard model, anomalies, and charge
quantization,
\newblock \emph{Physical Review D} \textbf{43} (1991), 2709--2717.
\newblock doi:10.1103/PhysRevD.43.2709.

\bibitem{Tong2017}
D.~Tong,
\newblock Line operators in the Standard Model,
\newblock \emph{Journal of High Energy Physics} \textbf{2017} (2017),
no.~7, 104.
\newblock doi:10.1007/JHEP07(2017)104.

\bibitem{GengMarshak1989}
C.~Q.~Geng and R.~E.~Marshak,
\newblock Uniqueness of quark and lepton representations in the Standard Model
from the anomalies viewpoint,
\newblock \emph{Physical Review D} \textbf{39} (1989), 693--696.
\newblock doi:10.1103/PhysRevD.39.693.

\bibitem{MinahanRamondWarner1990}
J.~A.~Minahan, P.~Ramond, and R.~C.~Warner,
\newblock Comment on anomaly cancellation in the Standard Model,
\newblock \emph{Physical Review D} \textbf{41} (1990), 715--716.
\newblock doi:10.1103/PhysRevD.41.715.

\bibitem{AlvarezGraciaBondiaMartin1995}
E.~Alvarez, J.~M.~Gracia-Bond\'ia, and C.~P.~Mart\'in,
\newblock Anomaly cancellation and the gauge group of the Standard Model in
NCG,
\newblock \emph{Physics Letters B} \textbf{364} (1995), 33--40.
\newblock doi:10.1016/0370-2693(95)01051-3.

\bibitem{Krasnov2021}
K.~Krasnov,
\newblock SO(9) characterization of the Standard Model gauge group,
\newblock \emph{Journal of Mathematical Physics} \textbf{62} (2021), 021703.
\newblock doi:10.1063/5.0039941.

\bibitem{Furey2018}
C.~Furey,
\newblock Three generations, two unbroken gauge symmetries, and one
eight-dimensional algebra,
\newblock \emph{Physics Letters B} \textbf{785} (2018), 84--89.
\newblock doi:10.1016/j.physletb.2018.08.032.

\bibitem{FureyRoadmap2025}
N.~Furey,
\newblock An algebraic roadmap of particle theories. Part II: theoretical checkpoints,
\newblock \emph{Annalen der Physik} \textbf{537} (2025), 2400323.
\newblock doi:10.1002/andp.202400323.

\bibitem{Gresnigt2026}
N.~G.~Gresnigt,
\newblock Three fermion generations and an untriplicated gauge sector with
$S_3$ family symmetry,
\newblock \emph{Physics Letters B} \textbf{880} (2026), 140782.
\newblock doi:10.1016/j.physletb.2026.140782.

\bibitem{Baez2026}
J.~C.~Baez,
\newblock Three generations in $E_7$,
\newblock arXiv:2608.06271 [hep-th], 2026.

\bibitem{Maldacena1998}
J.~M.~Maldacena,
\newblock The large $N$ limit of superconformal field theories and supergravity,
\newblock \emph{Advances in Theoretical and Mathematical Physics} \textbf{2}
(1998), 231--252.
\newblock doi:10.4310/ATMP.1998.v2.n2.a1.

\bibitem{Witten1998}
E.~Witten,
\newblock Anti de Sitter space and holography,
\newblock \emph{Advances in Theoretical and Mathematical Physics} \textbf{2}
(1998), 253--290.
\newblock doi:10.4310/ATMP.1998.v2.n2.a2.

\bibitem{Schiffler2014}
R.~Schiffler,
\newblock \emph{Quiver Representations},
\newblock CMS Books in Mathematics, Springer, Cham, 2014.

\bibitem{KuwaeShioya2003}
K.~Kuwae and T.~Shioya,
\newblock Convergence of spectral structures: a functional analytic theory and
its applications to spectral geometry,
\newblock \emph{Communications in Analysis and Geometry} \textbf{11} (2003),
599--673.
\newblock doi:10.4310/CAG.2003.v11.n4.a1.

\bibitem{Howe1989}
R.~Howe,
\newblock Remarks on classical invariant theory,
\newblock \emph{Transactions of the American Mathematical Society}
\textbf{313} (1989), no.~2, 539--570.
\newblock doi:10.1090/S0002-9947-1989-0986027-X.

\end{thebibliography}

\clearpage
\appendix

% =============================================================================
\section{Tetrahedral representation bookkeeping}
\label{app:S4}
% =============================================================================

This appendix records the representation-theoretic decomposition used in
Theorem~\ref{thm:geometry-short}.  It is independent of the probabilistic
origin of the source Gram.

Let $X=\{1,2,3,4\}$ and let $\Phi=\{(i,j):i\ne j\}$ be the twelve oriented
edges.  Root reversal is
\[
\iota(i,j)=(j,i).
\]
Let $P_\C=\C^\Phi$.  Its even and odd reversal subspaces are
\[
P^\pm=\{f:f(j,i)=\pm f(i,j)\},
\qquad
\dim P^\pm=6.
\]

\subsection{The even sector}

The even sector is the permutation representation on the six unordered edges
of a tetrahedron.  Let $\chi_+$ be its character.  For the five conjugacy
classes of $S_4$ represented by
\[
1,\quad (12),\quad (12)(34),\quad (123),\quad (1234),
\]
the number of fixed unordered edges is
\[
6,\quad 2,\quad 2,\quad 0,\quad 0.
\]
Comparing with the $S_4$ character table gives
\begin{equation}
\boxed{
P^+
\cong
\mathbf1\oplus W\oplus V_2,}
\label{eq:appendix-even}
\end{equation}
where $W$ is the standard three-dimensional representation and $V_2$ is the
two-dimensional irreducible representation.

\subsection{The odd sector}

The odd sector is the oriented edge space
$\Lambda^2\C^4$ after identifying
$e_i\wedge e_j$ with the signed edge $i\to j$.  Write
\[
\C^4=\C u_0\oplus W,
\qquad
u_0=\frac12(1,1,1,1).
\]
Then
\[
\Lambda^2\C^4
=
u_0\wedge W\oplus\Lambda^2W.
\]
The first summand is isomorphic to $W$.  Since $\dim W=3$,
\[
\Lambda^2W\cong W\otimes{\rm sgn}.
\]
Therefore
\begin{equation}
\boxed{
P^-
\cong
W\oplus(W\otimes{\rm sgn}).}
\label{eq:appendix-odd}
\end{equation}
The boundary map sends the first summand isomorphically onto the centred
endpoint space and annihilates the second.  Hence the two odd triplets are
canonically distinguished as endpoint and cycle sectors.

Combining \eqref{eq:appendix-even} and \eqref{eq:appendix-odd}, removing the
even root mean, and adjoining the scalar type contrast gives
\[
\C_{\rm type}^{+}
\oplus W_+\oplus(V_2)_+
\oplus W_-\oplus(W\otimes{\rm sgn})_-,
\]
which is \eqref{eq:residual-S4}.

% =============================================================================
\section{Finite double-commutant and support-polar details}
\label{app:commutant}
% =============================================================================

\subsection{Bicommutant closure}

For completeness, if $\mathcal B\subseteq\mathcal B(\Hh)$ is a unital
finite-dimensional $C^*$-algebra, then
\[
\mathcal B\cong
\bigoplus_\alpha
M_{n_\alpha}(\C)\otimes I_{m_\alpha}
\]
after a unitary decomposition of $\Hh$.  Its commutant is
\[
\mathcal B'
\cong
\bigoplus_\alpha
I_{n_\alpha}\otimes M_{m_\alpha}(\C),
\]
and therefore
\[
\mathcal B''=\mathcal B.
\]
This is the only bicommutant input used in
Theorem~\ref{thm:main-duality}.

\subsection{The support-polar algebra is intrinsic}

The proof of Theorem~\ref{thm:support-polar} uses polynomial functional
calculus only because the cutoff is finite.  If
\[
F^*F=\sum_{\lambda\in\Sigma}\lambda P_\lambda,
\qquad
\Sigma\subset[0,\infty)
\]
has finite spectrum, choose a polynomial $g$ with
\[
g(0)=0,
\qquad
g(\lambda)=\lambda^{-1/2}
\quad(\lambda\in\Sigma\setminus\{0\}).
\]
Then
\[
U=Fg(F^*F)
\]
is exactly the polar partial isometry.  Hence neither support projection nor
partial isometry is an additional finite input.  The construction fails only
if one asks for a uniformly bounded inverse as singular values collapse
cofinally; the finite duality itself remains exact.

\subsection{Basis independence of the quiver}

Let
\[
T_{ba}
=
\sum_\alpha E_\alpha\otimes B_\alpha
=
\sum_\beta \widetilde E_\beta\otimes\widetilde B_\beta
\]
be two Hilbert--Schmidt orthonormal expansions of the same operator.
The two coefficient families are related by a unitary transformation.
Therefore their linear spans in $\mathcal B(N_a,N_b)$ agree.  The arrow space
$\mathfrak B_{ba}(T)$ is thus basis independent.  The endomorphism algebra of
the quiver is consequently canonical up to the unitary gauges of the
multiplicity spaces and permutation of isomorphic central blocks.

% =============================================================================
\section{Multiplicity-birth and internal-router details}
\label{app:internal-landing}
% =============================================================================

This appendix records the finite word-algebra facts behind the internal
selection hierarchy of Section~\ref{sec:multiplicity-birth}.  The statements
are independent of a particular stochastic presentation once the labelled
source-word Gram and the external representation on its quotient are given.

\subsection{Flat word depth is terminal}

Let $\mathscr L$ be the finite represented source-letter alphabet on a
source-minimal quotient and let
\[
V_r=\Span\{\ell_k\cdots\ell_1\xi:
0\le k\le r,\ \ell_j\in\mathscr L,\ \xi\in V_0\}.
\]
Let $K_r$ be the Gram matrix of a labelled spanning family for $V_r$.

\begin{lemma}[Flat extension of a finite word span]
\label{lem:flat-word-span}
If $\rank K_r=\rank K_{r+1}$, then $V_{r+1}=V_r$, every letter in
$\mathscr L$ preserves $V_r$, and $V_s=V_r$ for all $s\ge r$.
Consequently every external isotypic multiplicity is constant for $s\ge r$.
\end{lemma}

\begin{proof}
The inclusion $V_r\subseteq V_{r+1}$ is automatic.  Equality of Gram ranks
makes the two finite-dimensional quotient spans equal.  Since every vector
$\ell v$ with $v\in V_r$ occurs among the words of length at most $r+1$, one
has $\ell V_r\subseteq V_{r+1}=V_r$ for each letter $\ell$.  Induction on word
length proves the remaining assertions.
\end{proof}

\begin{corollary}[Terminating multiplicity obstruction]
\label{cor:flat-multiplicity-obstruction}
At a flat word depth, failure of the multiplicity inequality
$m_{\lambda,r}\ge n$ is a permanent obstruction to an internal $M_n(\C)$
factor carried by the $\lambda$-multiplicity space.
\end{corollary}

\begin{proof}
Apply Lemma~\ref{lem:flat-word-span} and
Theorem~\ref{thm:active-isotypic}.
\end{proof}

\subsection{Router failure loci and generic saturation}

For the diagonal algebra $\mathcal D_n$ of
Theorem~\ref{thm:connected-router}, define the router failure set
\[
\mathfrak F_n
=
\{u\in M_n(\C):C^*(\mathcal D_n,u)\ne M_n(\C)\}.
\]
For every nonempty proper subset $S\subset\{1,\ldots,n\}$ let
\[
L_S
=
\{u:p_iup_j=p_jup_i=0\ \text{for all }i\in S,
\ j\notin S\}.
\]

\begin{proposition}[Algebraic structure of the router failure set]
\label{prop:router-failure-locus}
One has
\begin{equation}
\boxed{
\mathfrak F_n
=
\bigcup_{\varnothing\ne S\subsetneq\{1,\ldots,n\}}L_S.}
\label{eq:router-failure-locus}
\end{equation}
Each $L_S$ is a proper complex linear subspace of $M_n(\C)$.  Hence
$\mathfrak F_n$ is a proper Zariski-closed set and its complement is Zariski
open, Euclidean open, and dense.
\end{proposition}

\begin{proof}
The support graph is disconnected exactly when some nontrivial union $S$ of
connected components has no support edge to its complement, which is precisely
membership in $L_S$.  Each $L_S$ is defined by homogeneous linear equations
on matrix entries and is proper because, for example, a cyclic shift matrix
has connected support.  A finite union of closed algebraic subsets is closed,
and a finite union of proper linear subspaces cannot fill $M_n(\C)$.
\end{proof}

\begin{proposition}[Minimal bridge and weak-router residuals]
\label{prop:bridge-router-residuals}
On $T\oplus H_{\rm priv}$ define
\[
\delta_{\rm col}(u)
:=
\norm{p_Hup_T}_{\HS}^2+\norm{p_Tup_H}_{\HS}^2.
\]
On a weak multiplicity plane define
\[
\delta_{\rm wk}(p,q):=\norm{[p,q]}_{\HS}^2.
\]
Then $\delta_{\rm col}=0$ exactly on the reducing colour branch
$M_2\oplus\C$, while $\delta_{\rm col}>0$ gives $M_3$.  For rank-one
projections $p,q$ on $\C^2$, one has
\begin{equation}
\delta_{\rm wk}(p,q)
=2c^2(1-c^2),
\qquad c=\abs{\langle t,h\rangle},
\label{eq:weak-router-residual}
\end{equation}
so $\delta_{\rm wk}>0$ exactly when the projections are distinct and
nonorthogonal, in which case they generate $M_2$.
\end{proposition}

\begin{proof}
The colour statement is Theorem~\ref{thm:colour-bridge}.  For the weak
statement, choose a basis with $p=|e_1\rangle\langle e_1|$ and
$h=ce_1+\sqrt{1-c^2}e_2$ after absorbing a phase.  Direct multiplication gives
\eqref{eq:weak-router-residual}.  Its strict positivity is equivalent to
$0<c<1$, which is the Burnside condition of Lemma~\ref{lem:weak-M2}.
\end{proof}

\subsection{A terminating internal landing certificate}

The preceding results give a finite decision tree which does not fit an
internal algebra by dimension.  Start from a labelled physical word Gram and
its external isotypic census.

\begin{theorem}[Finite internal landing alternative]
\label{thm:finite-internal-landing-alternative}
At any finite word depth exactly one of the following outcomes applies to the
colour/weak landing problem.
\begin{enumerate}[label=\textup{(A\arabic*)},leftmargin=3em]
\item If the word hierarchy is flat and $m_{0,r}<3$, an external-trivial
      internal colour $M_3(\C)$ is impossible at every later depth.
\item If $m_{0,r}\ge3$ but every represented type--private bridge has
      $\delta_{\rm col}=0$, the colour block is reducible and remains
      $M_2\oplus\C$ on the displayed three-space.
\item If the weak multiplicity plane is present but every available pair has
      $\delta_{\rm wk}=0$, the represented weak multiplicity algebra is
      commutative.
\item If the multiplicity census is $(3,2,1,1)$ and both a positive colour
      residual and a positive weak-router residual occur, the represented
      internal word algebra equals the full relative commutant
      $M_3\oplus M_2\oplus\C\oplus\C$.
\end{enumerate}
On the fixed census $(3,2,1,1)$ the intermediate algebra dimensions are
exactly those of Theorem~\ref{thm:internal-deficit-alternatives}.  Therefore
every failure returns a finite algebraic witness and no failed condition is
silently completed to the terminal branch.
\end{theorem}

\begin{proof}
Item (A1) is Corollary~\ref{cor:flat-multiplicity-obstruction}.  Item (A2) is
Theorem~\ref{thm:colour-bridge}; item (A3) is
Lemma~\ref{lem:weak-M2}; and item (A4) is
Theorem~\ref{thm:internal-seed-saturation}.  The dimension statement is
Theorem~\ref{thm:internal-deficit-alternatives}.
\end{proof}

\begin{remark}[What this finite alternative does not claim]
The theorem classifies and terminates the internal landing once the labelled
source-word hierarchy is supplied.  It does not assert that an arbitrary
microscopic renewal law must realize the multiplicity birth, bridge, or router
rows.  Their occurrence is an operator-valued provenance question.  This is
precisely why the determinant line and matter incidence are tested only after
the relative commutant has landed.
\end{remark}

% =============================================================================

\section{Determinant incidence and anomaly bookkeeping}
\label{app:anomaly}
% =============================================================================

\subsection{The determinant line}

For $C\cong\C^3$ and $W_2\cong\C^2$, define
\[
\mathscr D=(\det C\otimes\det W_2)^*,
\qquad
\mathscr U=
\Hom(C\otimes W_2\otimes W_2,\Lambda^2C^*).
\]
For $\tau\in\mathscr D$, define
\[
\bigl(\Theta_\tau(c\otimes w_1\otimes w_2)\bigr)(c_1,c_2)
=
\tau(c\wedge c_1\wedge c_2\otimes w_1\wedge w_2).
\]
Complete alternation in the three colour and two weak arguments gives a map
$\mathfrak A_{\det}:\mathscr U\to\mathscr D$.

\begin{proposition}[Determinant-incidence line]
\label{prop:det-incidence}
The maps $\Theta$ and $\mathfrak A_{\det}$ restrict to inverse isomorphisms
between $\mathscr D$ and the fully alternating incidence line.  In oriented
orthonormal bases,
\[
\mathfrak A_{\det}(\Theta_\tau)=\tau,
\qquad
\norm{\Theta_\tau}_{\HS}^2=6\norm{\tau}^2.
\]
Thus a nonzero fully alternating incidence shadow is equivalent to a nonzero
determinant seed.
\end{proposition}

\begin{proof}
Complete alternation is the orthogonal projection onto the unique
$(\det C\,\det W_2)^{-1}$ symmetry type.  Evaluation on oriented orthonormal
bases gives the factor $3!\,2!/2=6$ in the Hilbert--Schmidt normalization.
\end{proof}

\subsection{Hypercharge in left-handed convention}

With $(a,b)=(-2,3)$, the physical right-handed convention gives
\[
Q:(\mathbf3,\mathbf2)_1,\qquad
u:(\mathbf3,\mathbf1)_4,\qquad
d:(\mathbf3,\mathbf1)_{-2},
\]
\[
L:(\mathbf1,\mathbf2)_{-3},\qquad
e:(\mathbf1,\mathbf1)_{-6},\qquad
H:(\mathbf1,\mathbf2)_3.
\]
Writing all fermions as left-handed Weyl fields replaces
\[
u\mapsto u^c:(\overline{\mathbf3},\mathbf1)_{-4},
\qquad
d\mapsto d^c:(\overline{\mathbf3},\mathbf1)_2,
\qquad
e\mapsto e^c:(\mathbf1,\mathbf1)_6.
\]
The common kernel of these representations of
$\SU(3)\times\SU(2)\times\U(1)_y$ is generated by
\[
\left(e^{2\pi i/3}I_3,-I_2,e^{i\pi/3}\right),
\]
hence is $\Z_6$.  This agrees with the quotient in
\eqref{eq:SM-gauge-group}.

% =============================================================================
\section{Finite common-action representative}
\label{app:finite-common-action}

The finite commutant duality is independent of a variational action.  This
appendix records one explicit local gauge--Higgs--Dirac representative on the
reconstructed spacetime complex, together with the Ward, BRST, stress, and
vacuum identities used only in the compatibility discussion of
Section~\ref{sec:coupled-action}.

Let $V_X,E_X,F_X$ be the finite spacetime vertices, oriented edges, and
plaquettes of a reconstructed regulator, with positive Hodge weights
$\star_{0,X},\star_{1,X},\star_{2,X}$.  A finite internal configuration is
\[
U_e\in G_{\rm SM},
\qquad
H_v\in W_2,
\qquad
\Psi_v\in\mathcal S_v\otimes G_{\rm gen}\otimes\mathcal H_F^{(1)},
\]
together with an independent dual amplitude $\overline\Psi_v$.
Let $F_{p,X}$ be the principal-logarithm plaquette curvature on its local
chart and let $D_UH$ be the covariant edge difference.

For positive gauge coefficients $g_3,g_2,g_1$, scalar coefficient
$\lambda_H>0$, and $v_H\ge0$, define
\begin{equation}
\boxed{
\begin{aligned}
\mathscr S_{{\rm SM},X}^{E}
={}&
\frac12\langle F_X,\star_{2,X}F_X\rangle_{g_3,g_2,g_1}
+\frac12\langle D_UH,\star_{1,X}D_UH\rangle\\
&+
\sum_{v}m_{v,X}\lambda_H
(\norm{H_v}^2-v_H^2)^2\\
&+
\operatorname{Re}
\left\langle
\overline\Psi,
\star_{0,X}
\bigl(
D_{{\rm sp},X}^{U}
+\Gamma_{\rm Cl}\otimes D_F(H)
\bigr)\Psi
\right\rangle .
\end{aligned}}
\label{eq:finite-SM-action}
\end{equation}
The coefficients in \eqref{eq:finite-SM-action} are physical parameters and
are not fixed by Theorem~\ref{thm:structural-SM}.

\begin{theorem}[Finite Ward, BRST, and stress identities]
\label{thm:finite-action-identities}
The action \eqref{eq:finite-SM-action} is exactly gauge invariant and has
cutoff-independent local support radius.  Its Euler covectors obey the finite
Ward identity
\begin{equation}
\boxed{
d_U^*E_U
+\mathcal R_H^*E_H
+\mathcal R_\Psi^*E_\Psi
+\mathcal R_{\overline\Psi}^*E_{\overline\Psi}
=0.}
\label{eq:ward}
\end{equation}
The standard ghost extension
\[
sc=-c^2,\qquad
sU=c_tU-Uc_s,\qquad
sH=cH,\qquad
s\Psi=c\Psi,\qquad
s\bar c=B,\qquad
sB=0
\]
is nilpotent and satisfies
\begin{equation}
\boxed{s^2=0,\qquad s\mathscr S_{{\rm SM},X}^{E}=0.}
\label{eq:brst}
\end{equation}
Metric/coframe variation defines the matter stress tensor, and common
relabeling invariance gives
\begin{equation}
\boxed{
\mathcal L_q^*\mathsf T
=
2\mathcal L_U^*E_U
+2\mathcal L_H^*E_H
+2\mathcal L_\Psi^*E_\Psi
+2\mathcal L_{\overline\Psi}^*E_{\overline\Psi}.}
\label{eq:stress-identity}
\end{equation}
\end{theorem}

\begin{proof}
Plaquette holonomy transforms by conjugation, the Higgs edge difference
covariantly in the target fibre, and the finite Dirac/Yukawa term intertwines
the left and right typed representations.  These identities prove gauge
invariance.  Differentiating gauge invariance gives \eqref{eq:ward}.
Nilpotence follows from the graded Jacobi identity and covariance.
Differentiating common relabeling invariance gives
\eqref{eq:stress-identity}.
\end{proof}

\begin{corollary}[Stationary internal vacuum]
\label{cor:internal-vacuum}
The configuration
\[
U_e=I,
\qquad
H_v=v_Hh_0,
\qquad
\Psi_v=\overline\Psi_v=0
\]
is a stationary internal vacuum.  For $v_H>0$ its gauge stabilizer is
\[
U(3)\cong
\frac{\SU(3)\times\U(1)_{\rm em}}{\Z_3}.
\]
\end{corollary}

\section{Perturbative stability of the commutant gap}
\label{app:stability}
% =============================================================================

This appendix supplies the elementary estimate behind
Theorem~\ref{thm:finite-anchor}.

Let
\[
\partial_{\bm c}:
\mathcal B_2(\Hh)\to
\mathcal B_2(\Hh)^{\oplus s},
\qquad
\partial_{\bm c}X=([c_1,X],\ldots,[c_s,X]).
\]
Then
\[
\mathscr L_{\bm c}
=
\partial_{\bm c}^*\partial_{\bm c}.
\]
If
\[
\varepsilon
=
\left(\sum_j\norm{c_j-d_j}_{\rm op}^2\right)^{1/2},
\qquad
C_{\bm c}
=
\left(\sum_j\norm{c_j}_{\rm op}^2\right)^{1/2},
\]
then
\[
\norm{\partial_{\bm c}-\partial_{\bm d}}_{2\to2}
\le2\varepsilon,
\qquad
\norm{\partial_{\bm c}}_{2\to2}
\le2C_{\bm c}.
\]
Therefore
\[
\begin{aligned}
\norm{\mathscr L_{\bm c}-\mathscr L_{\bm d}}
&=
\norm{
\partial_{\bm c}^*\partial_{\bm c}
-
\partial_{\bm d}^*\partial_{\bm d}
}\\
&\le
\norm{\partial_{\bm c}}\,
\norm{\partial_{\bm c}-\partial_{\bm d}}
+
\norm{\partial_{\bm c}-\partial_{\bm d}}\,
\norm{\partial_{\bm d}}\\
&\le
4(C_{\bm c}+C_{\bm d})\varepsilon.
\end{aligned}
\]
This proves the perturbative lower bound used in
\eqref{eq:anchor-lock}.

The perturbative estimate above is local in generator norm.  The remainder of
this appendix proves the compact-Mosco mechanism used in
Section~\ref{sec:compact-mosco}.

\subsection{Transported forms and resolvent minimizers}
\label{app:compact-mosco}

Retain the embeddings $J_X:\mathscr H_X\to\mathscr H$ and the closed
nonnegative forms $q_X$ of Definitions~\ref{def:transported-mosco}
and~\ref{def:asymptotic-compactness}.  For $f\in\mathscr H$ let
\begin{equation}
u_X=(\mathscr L_X+I)^{-1}J_X^*f.
\label{eq:resolvent-minimizer}
\end{equation}
Equivalently, $u_X$ is the unique minimizer on $D(q_X)$ of
\begin{equation}
\mathcal F_{X,f}(u)
=q_X(u)+\norm{u}^2-2\operatorname{Re}\ip{J_Xu}{f}_{\mathscr H}.
\label{eq:resolvent-functional}
\end{equation}
Writing $q_X(\cdot,\cdot)$ for the sesquilinear form obtained by polarizing
$q_X$, the Euler identity is
\begin{equation}
q_X(u_X,v)+\ip{u_X}{v}
=\ip{f}{J_Xv}_{\mathscr H},
\qquad v\in D(q_X).
\label{eq:resolvent-euler}
\end{equation}
Taking $v=u_X$ gives
\begin{equation}
\norm{u_X}^2+q_X(u_X)
\le \norm{f}\,\norm{u_X}
\le \norm{f}^2,
\qquad
\norm{u_X}\le\norm f.
\label{eq:resolvent-energy-bound}
\end{equation}

\begin{lemma}[Mosco convergence gives strong transported resolvents]
\label{lem:mosco-strong-resolvent}
Under Definition~\ref{def:transported-mosco}, for every fixed
$f\in\mathscr H$,
\begin{equation}
J_X(\mathscr L_X+I)^{-1}J_X^*f
\longrightarrow
J_\infty(\mathscr L_\infty+I)^{-1}J_\infty^*f
\label{eq:strong-resolvent}
\end{equation}
strongly in $\mathscr H$.
\end{lemma}

\begin{proof}
The estimate \eqref{eq:resolvent-energy-bound} makes $(J_Xu_X)$ bounded.
Every weakly convergent subsequence has, by the Mosco liminf inequality, a
limit $J_\infty u$ satisfying
\[
q_\infty(u)+\norm u^2-2\operatorname{Re}\ip{J_\infty u}{f}
\le\liminf_X\mathcal F_{X,f}(u_X).
\]
For every $v\in D(q_\infty)$ choose a Mosco recovery sequence $v_X$.
Minimality of $u_X$ gives
$\mathcal F_{X,f}(u_X)\le\mathcal F_{X,f}(v_X)$; passing to the limit shows
that $u$ minimizes $\mathcal F_{\infty,f}$.  Strict convexity identifies
$u=(\mathscr L_\infty+I)^{-1}J_\infty^*f$ and therefore identifies every weak
cluster point.  Comparing the minimum values and using the liminf inequality
gives convergence of $\norm{u_X}^2+q_X(u_X)$ to the limiting value.  Weak
convergence together with convergence of the Hilbert norms yields strong
convergence of $J_Xu_X$.  Thus the whole sequence converges.
\end{proof}

\begin{lemma}[Asymptotic compactness gives collective compactness]
\label{lem:collective-compactness}
Assume Definition~\ref{def:asymptotic-compactness}.  If $(f_X)$ is bounded in
$\mathscr H$, then every sequence
\[
J_X(\mathscr L_X+I)^{-1}J_X^*f_X
\]
has a strongly convergent subsequence.  Hence the transported resolvents
$(R_X)$ form a collectively compact family.  The limiting resolvent $R_\infty$
is compact.
\end{lemma}

\begin{proof}
Apply \eqref{eq:resolvent-energy-bound} with $f=f_X$.  The corresponding
resolvent vectors have uniformly bounded norms and form energies, so
asymptotic compactness gives a strongly convergent subsequence.  The same
argument for the limiting form, applied to an arbitrary bounded sequence of
right-hand sides, proves compactness of $R_\infty$.
\end{proof}

\begin{lemma}[Collectively compact self-adjoint convergence is norm convergence]
\label{lem:collective-norm}
Let $T_X$ be bounded self-adjoint operators on a Hilbert space.  If
$T_X\to T$ strongly and $(T_X)$ is collectively compact, then $T$ is compact
and
\begin{equation}
\norm{T_X-T}_{\rm op}\longrightarrow0.
\label{eq:collective-norm}
\end{equation}
\end{lemma}

\begin{proof}
First let $(x_n)$ be bounded.  For each $n$ choose a cofinal index $X_n$ so
large that $\norm{(T_{X_n}-T)x_n}<1/n$.  Collective compactness gives a
convergent subsequence of $T_{X_n}x_n$, and hence of $Tx_n$; thus $T$ is
compact.

If \eqref{eq:collective-norm} failed, there would be unit vectors $x_X$ and
$\varepsilon>0$ with $\norm{(T_X-T)x_X}\ge\varepsilon$.  Passing to a
subsequence, take $x_X\rightharpoonup x$, $T_Xx_X\to y$ by collective
compactness, and $Tx_X\to Tx$ by compactness of $T$.  For every fixed $z$,
self-adjointness and strong convergence give
\[
\ip{y}{z}
=\lim_X\ip{x_X}{T_Xz}
=\ip{x}{Tz}
=\ip{Tx}{z}.
\]
Thus $y=Tx$, so $(T_X-T)x_X\to0$, a contradiction.
\end{proof}

\begin{proposition}[Proof of compact Mosco resolvent convergence]
\label{prop:app-compact-mosco}
Under the hypotheses of
Proposition~\ref{prop:compact-mosco-resolvent},
\[
\norm{R_X-R_\infty}_{\rm op}\to0.
\]
\end{proposition}

\begin{proof}
Lemma~\ref{lem:mosco-strong-resolvent} gives strong convergence,
Lemma~\ref{lem:collective-compactness} gives collective compactness, and
Lemma~\ref{lem:collective-norm} upgrades the convergence to operator norm.
\end{proof}

\subsection{Eigenvalues, Riesz projections, and kernel locking}

Let $\mu_{1,X}\ge\mu_{2,X}\ge\cdots>0$ denote the nonzero eigenvalues of the
compact positive operator $R_X$, repeated with multiplicity.  On the screened
space they are related to the eigenvalues of $\mathscr L_X$ by
\begin{equation}
\mu_{k,X}=\frac{1}{1+\lambda_{k,X}}.
\label{eq:resolvent-spectrum-map}
\end{equation}

\begin{lemma}[Low-spectrum convergence]
\label{lem:low-spectrum-convergence}
If $\norm{R_X-R_\infty}\to0$, then for every fixed $k$,
\begin{equation}
\mu_{k,X}\to\mu_{k,\infty},
\qquad
\lambda_{k,X}\to\lambda_{k,\infty}
\label{eq:low-spectrum-convergence}
\end{equation}
whenever the limiting eigenvalue is finite.  If $\Sigma$ is an isolated finite
spectral cluster of $R_\infty$ and $\Gamma$ is a positively oriented contour
separating it from the remaining spectrum, then the transported Riesz
projections
\begin{equation}
\Pi_X(\Gamma)
=\frac{1}{2\pi i}\int_\Gamma(z-R_X)^{-1}\,\dd z
\label{eq:riesz-projection}
\end{equation}
converge to $\Pi_\infty(\Gamma)$ in operator norm.  Their ranks are eventually
equal.
\end{lemma}

\begin{proof}
For compact self-adjoint operators the min--max principle gives
$|\mu_{k,X}-\mu_{k,\infty}|\le\norm{R_X-R_\infty}$ for each fixed $k$.
Equation~\eqref{eq:resolvent-spectrum-map} gives the second convergence.
For the projection statement, norm convergence and the resolvent identity
imply uniform convergence of $(z-R_X)^{-1}$ on $\Gamma$; integrate around the
contour.  If two projections differ in operator norm by less than one, their
ranges have the same finite dimension, proving eventual rank stability.
\end{proof}

\begin{proposition}[Protected-kernel locking]
\label{prop:protected-kernel-locking}
Suppose each $\M_X$ is an $m$-dimensional subspace of
$\Ker\mathscr L_X$ and
$\Ker\mathscr L_\infty=\M_\infty$ has dimension $m$.  Under compact Mosco
convergence, for all sufficiently late $X$,
\[
\Ker\mathscr L_X=\M_X,
\]
the transported kernel projections converge in operator norm, and
\[
\lambda_{m+1}(\mathscr L_X)
\to\lambda_{m+1}(\mathscr L_\infty)>0.
\]
\end{proposition}

\begin{proof}
The eigenvalue $0$ of $\mathscr L_\infty$ corresponds to the eigenvalue $1$
of $R_\infty$, of multiplicity $m$.  Compactness of the limiting resolvent and
finite-dimensionality of the kernel separate this eigenvalue from the rest of
the spectrum.  Lemma~\ref{lem:low-spectrum-convergence} therefore gives a
nearby spectral cluster of rank exactly $m$ for all sufficiently late $X$.
Because $\M_X\subseteq\Ker\mathscr L_X$ already supplies $m$ exact copies of
the eigenvalue $1$ of $R_X$, that entire cluster is the kernel and no extra
zero mode can occur.  The same lemma gives projection and next-eigenvalue
convergence.
\end{proof}

This proposition is the spectral step used in
Theorem~\ref{thm:cofinal-duality}.  It shows explicitly that the fixed protected
zero modes and compact-Mosco convergence exclude spectral pollution at zero;
norm-resolvent convergence and kernel stability are conclusions of the form
hypotheses rather than additional assumptions.

% =============================================================================
\section{A one-generation incidence audit}
\label{app:one-generation}
% =============================================================================

The commutant language also gives a compact check of how Yukawa incidence
reduces residual scalar symmetries inside one generation.

Let
\[
\mathcal F_{15}=Q\oplus u\oplus d\oplus L\oplus e
\]
with the modules of \eqref{eq:SM-types}.  Ignoring generation multiplicity for
the moment, the gauge-generated algebra is
\begin{equation}
\A_{\rm gauge}
=
M_6(\C)\oplus M_3(\C)\oplus M_3(\C)\oplus M_2(\C)\oplus\C,
\label{eq:one-gen-gauge-algebra}
\end{equation}
because the five modules are irreducible and pairwise inequivalent.  Hence
\[
\A_{\rm gauge}'\cong\C^5.
\]

\begin{proposition}[Incidence reduces the one-generation commutant]
\label{prop:one-gen-incidence}
Assume the three Higgs-contracted incidence histories
\[
Q\leftrightarrow u,\qquad
Q\leftrightarrow d,\qquad
L\leftrightarrow e
\]
are nonzero.  Then
\[
\boxed{
\A_{\rm gauge+inc}
=
M_{12}(\C)\oplus M_3(\C),
\qquad
\A_{\rm gauge+inc}'\cong\C^2.}
\]
The two residual central directions are the quark and lepton connected
components.  Adding a neutral singlet and a nonzero $L$--neutral incidence
changes the second matrix block to $M_4(\C)$ while leaving the residual
commutant $\C^2$.
\end{proposition}

\begin{proof}
A nonzero block between two full matrix summands, together with multiplication
by endpoint matrix units, generates every rank-one map between the summands.
Products and adjoints therefore generate the full matrix algebra on each
connected component of the incidence graph.  The three displayed edges have
components $\{Q,u,d\}$ and $\{L,e\}$, of dimensions $12$ and $3$.
\end{proof}

\begin{remark}
This audit is intentionally distinct from the generation multiplicity theorem.
It describes the commutant inside one fixed generation packet.  The
rank-three generation carrier of Proposition~\ref{prop:generation-carrier}
is an additional multiplicity factor reconstructed before numerical flavour
breaking is specified.
\end{remark}



% =============================================================================
\section{Incidence determinant support and deletion margins}
\label{app:incidence-details}
% =============================================================================

The rank theorem in Section~\ref{sec:incidence-skeleton} has a determinant
refinement useful for finite audits.  Keep the notation there and introduce
nonnegative incidence weights $w=(w_e)_{e\in E}$:
\begin{equation}
 G(w)=\sum_{e\in E}w_eG_e,
 \qquad P_E(w)=\det G(w).
 \label{eq:weighted-incidence-Gram}
\end{equation}

\begin{theorem}[Positive incidence determinant and sufficient supports]
\label{thm:incidence-determinant}
The polynomial $P_E$ is homogeneous of degree $n$ and has nonnegative real
coefficients.  For any $S\subseteq E$, if $w_e>0$ on $S$ and $w_e=0$ off
$S$, then
\begin{equation}
 \boxed{P_E(w)>0\quad\Longleftrightarrow\quad f(S)=n
 \quad\Longleftrightarrow\quad \M(S)=\M(E).}
 \label{eq:incidence-face-exactness}
\end{equation}
If
$\kappa_e=n-f(E\setminus\{e\})$, then the smallest power of $w_e$ occurring
with nonzero coefficient in $P_E$ is exactly $\kappa_e$.  Thus $e$ is
essential if and only if $w_e$ divides $P_E$.
\end{theorem}

\begin{proof}
Stack the row blocks $\sqrt{w_e}D_e$.  Cauchy--Binet expresses
$\det G(w)$ as a sum of squared absolute values of $n\times n$ row minors,
each multiplied by the monomial recording how many rows were selected from
each edge block.  All coefficients are therefore nonnegative.  A point on a
coordinate face has positive determinant exactly when the row spaces of the
retained blocks span $V$, which is $f(S)=n$.  For the last assertion, any row
basis uses at most $f(E\setminus\{e\})$ rows from the other blocks and hence
at least $\kappa_e$ rows from block $e$; extending a basis of the other-block
span using rows from $e$ attains this number.
\end{proof}

For quantitative deletion, put $G=G_E\succ0$ and
\[
 B_e=G^{-1/2}G_eG^{-1/2},\qquad \sum_eB_e=I,
 \qquad B_R=\sum_{e\in R}B_e.
\]

\begin{proposition}[Exact simultaneous-deletion margin]
\label{prop:deletion-margin}
For every $R\subseteq E$,
\begin{equation}
 G_{E\setminus R}=G^{1/2}(I-B_R)G^{1/2},
 \label{eq:deletion-factor}
\end{equation}
and therefore
\begin{equation}
 \boxed{\M(E\setminus R)=\M(E)
 \quad\Longleftrightarrow\quad \|B_R\|_{\rm op}<1.}
 \label{eq:deletion-margin}
\end{equation}
If $\beta_R=\|B_R\|_{\rm op}<1$, then
\[
 G_{E\setminus R}\succeq(1-\beta_R)G,
 \qquad
 \lambda_{\min}(G_{E\setminus R})
 \ge(1-\beta_R)\lambda_{\min}(G).
\]
\end{proposition}

\begin{proof}
The factorization follows from the definitions.  Since
$0\preceq B_R\preceq I$, $I-B_R$ is positive definite exactly when its
largest eigenvalue is smaller than one.  Congruence by $G^{1/2}$ preserves
nullity and positivity, giving both the commutant criterion and the lower
bound.
\end{proof}

\begin{example}[Minimal incidence skeletons need not form matroid bases]
\label{ex:nonmatroid-skeleton}
On $\C^3$ with scalar base algebra let
\[
 Y_a=E_{12}+E_{23},\qquad Y_b=E_{12},\qquad Y_c=E_{23}.
\]
The singleton $\{a\}$ and the pair $\{b,c\}$ both generate $M_3(\C)$ after
adjoining adjoints, and hence both have scalar commutant.  Each is
inclusion-minimal: neither $Y_b$ nor $Y_c$ alone has scalar commutant.  Thus
minimal sufficient incidence sets can have unequal cardinality, so the
incidence labels do not carry a matroid-basis structure in general.
\end{example}

% =============================================================================
\section{Historical operation, reciprocal-return, and determinant-source details}
\label{app:historical-source-selection}
% =============================================================================

\subsection{Large invariant word multiplicity need not be physical colour}

Let $\Omega=\{(i,j):1\le i,j\le4,\ i\ne j\}$ and
$\Hh=\C^\Omega$, with the simultaneous relabeling action
$u_g|i,j\rangle=|g(i),g(j)\rangle$ of $S_4$.  Adjoin all diagonal point
projections and the two fibre averages
\[
 (Tf)(i,j)=\frac13\sum_{k\ne i}f(i,k),\qquad
 (Hf)(i,j)=\frac13\sum_{k\ne j}f(k,j).
\]
Their support graph is connected, so the generated system algebra is
$M_{12}(\C)$.

\begin{proposition}[Invariant word-space inflation without an internal $M_3$]
\label{prop:word-space-inflation}
For the preceding benchmark,
\[
 \A=M_{12}(\C),\qquad
 \A^{S_4}\cong\C\oplus M_2(\C)\oplus\C\oplus\C.
\]
In the irreducible order
$(\mathbf1,\operatorname{sgn},W,W\otimes\operatorname{sgn},V_2)$, the
multiplicities of the physical $S_4$ action on $\Hh$ are
$(1,0,2,1,1)$, whereas the multiplicities of the conjugation action on
$L^2(M_{12})$ are $(7,5,19,17,12)$.  Thus the latter commutant contains an
$M_7(\C)$ block although the word-induced invariant system algebra has no
$M_3(\C)$ corner.
\end{proposition}

\begin{proof}
The permutation character on ordered unequal pairs is
$(12,2,0,0,0)$ on the conjugacy classes
$1,(12),(12)(34),(123),(1234)$.  Character inner products with the five
irreducibles of $S_4$ give $(1,0,2,1,1)$ and hence the displayed physical
commutant.  Conjugation on $M_{12}\cong\Hh\otimes\overline\Hh$ has character
$|\chi_\Hh|^2=(144,4,0,0,0)$, whose character inner products are
$(7,5,19,17,12)$.  Theorem~\ref{thm:word-module-recognition} retains only
left multiplication by $\A^{S_4}$ as word-induced invariant system actions.
\end{proof}

\subsection{Subdirect products and central separation}

\begin{theorem}[Finite subdirect matrix-algebra classification]
\label{thm:subdirect-classification}
Let $\mathcal B\subseteq\bigoplus_{i=1}^rM_{n_i}(\C)$ be a unital
$C^*$-subalgebra whose coordinate projection onto every displayed factor is
surjective.  Then the coordinates partition into classes of equal matrix size,
and $\mathcal B$ contains one independent full matrix algebra per class,
embedded diagonally across that class up to inner automorphisms.  In
particular, distinct matrix sizes cannot be diagonally identified.
\end{theorem}

\begin{proof}
Write the intrinsic decomposition
$\mathcal B\cong\bigoplus_\beta M_{m_\beta}(\C)$.  A surjective
$*$-homomorphism to a simple factor $M_{n_i}(\C)$ annihilates every intrinsic
summand except one: the images of the intrinsic central units are orthogonal
central projections summing to the identity, so exactly one is nonzero.  On
that summand the map is a nonzero homomorphism between simple full matrix
algebras and hence an isomorphism, forcing $m_\beta=n_i$.  Coordinates seeing
the same intrinsic summand are related by automorphisms of a full complex
matrix algebra, hence by inner automorphisms.  Different intrinsic summands
are independent.
\end{proof}

Theorem~\ref{thm:central-separation} is the specialization to block sizes
$(3,2,1,1)$.  It explains why independent scalar supports, rather than
blockwise surjectivity alone, are needed before the four dimensions in
Theorem~\ref{thm:internal-deficit-alternatives} may be added.

\subsection{A route-bank stopping projection}

For the mediator notation of Theorem~\ref{thm:reciprocal-return}, set
\[
 \mathcal N_A=\Span\{h^k\Ran\rho_A:0\le k<n\},\qquad
 \mathcal N_B=\Span\{h^k\Ran\rho_B:0\le k<n\}.
\]
Cayley--Hamilton makes both spaces $h$-invariant, and self-adjointness of $h$
makes them reducing.

\begin{proposition}[All-excursion return alternative]
\label{prop:all-excursion-return}
If $\mathfrak m_{\rm ret}=0$, then
$\mathcal N_A\perp\mathcal N_B$ and
\[
 Q=I_T\oplus0_{H_{\rm priv}}\oplus(I_V\otimes P_{\mathcal N_A})
\]
commutes with every generator of the declared reciprocal route bank and its
adjoint.  Consequently every word $w$ in that bank has
$p_Hwp_T=0$.  Thus failure of the finite return panel is an exact obstruction
to arbitrarily long repeated excursions within the same bank.
\end{proposition}

\begin{proof}
The equality $\mathfrak m_{\rm ret}=d^{-1}\Tr(\rho_BW_A)=0$ and positivity
imply $\Ran\rho_B\perp\mathcal N_A$, hence
$\mathcal N_B\perp\mathcal N_A$.  Partial-trace positivity gives
$\Ran A\subseteq V\otimes\Ran\rho_A$ and
$\Ran B\subseteq V\otimes\Ran\rho_B$.  The stated projection therefore
reduces the entry, exit, mediator and adjoint generators, contains the type
line and annihilates the private plane.  Every generated word has zero
$H_{\rm priv}\leftarrow T$ corner.
\end{proof}

\subsection{Sharp arity of the determinant character}

Let a native finite carrier have defining central weights $(1,0)$ on
$C\cong\C^3$, $(0,1)$ on $W_2\cong\C^2$, and $(0,0)$ on any neutral sector.
The centre of $\U(3)\times\U(2)$ acts on a weight $(a,b)$ by
$e^{i(a\alpha+b\beta)}$.

\begin{theorem}[Sharp five-slot determinant-mode floor]
\label{thm:determinant-arity-floor}
The central-weight support of the single-system endomorphism algebra contains
only
\[
 (0,0),\qquad \pm(1,0),\qquad \pm(0,1),\qquad \pm(1,-1),
\]
and therefore does not contain the determinant-incidence weight $(-3,-2)$.
On a diagonal $k$-fold tensor carrier this character is absent for $k<5$ and
is present for $k=5$.
\end{theorem}

\begin{proof}
Single-system endomorphism weights are target-minus-source differences among
$(1,0),(0,1),(0,0)$, giving the displayed list.  In $k$ tensor slots every
vector weight $(a,b)$ has $a,b\ge0$ and $a+b\le k$, so a target-minus-source
weight with total negative charge $-5$ cannot occur for $k<5$.  At $k=5$ the
alternating line in $C^{\otimes3}\otimes W_2^{\otimes2}$ transforms by
$\det C\,\det W_2$; a rank-one map from that line to a neutral line therefore
has weight $(-3,-2)$.
\end{proof}

This result concerns tensor arity, not temporal word length.  It explains why
longer evolution inside a fixed low-arity memory cannot substitute for the
typed five-leg incidence of Proposition~\ref{prop:main-det-incidence}.

% =============================================================================
\section{Source identifiability and finite system tomography}
\label{app:source-identifiability}
% =============================================================================

This appendix gives a finite witness for
Theorem~\ref{thm:record-nonidentifiability} and an explicit panel realizing
Corollary~\ref{cor:finite-system-export}.  The construction is algebraic and
is not assigned particle-physics semantics.

Let $H_E$ be a finite edge-system space, let
$Q_a$ ($1\le a\le4$) be nonzero orthogonal projections on a record factor
$H_R$ with $\sum_aQ_a=I$, and let $K_1,\ldots,K_6\in\mathcal B(H_E)$ satisfy
\begin{equation}
 \sum_{e=1}^6K_e^*K_e=I,
 \qquad
 \sum_{e=1}^6K_e=\sqrt2\,I,
 \qquad
 K_1,\ldots,K_6\ \text{linearly independent}.
 \label{eq:six-edge-normalization}
\end{equation}
For $0<\vartheta<1$ set
\[
 A_{e,a}=K_e\otimes Q_a,
 \qquad L_0=\sqrt{1-\vartheta}\,I,
 \qquad L_{e,a}=\sqrt\vartheta\,A_{e,a}.
\]
These operators define a normalized resolved instrument.

Choose a real $6\times5$ matrix $O$ with
\begin{equation}
 O^{\mathsf T}O=I_5,
 \qquad O^{\mathsf T}\mathbf1=0,
 \qquad OO^{\mathsf T}=I_6-\frac16\mathbf1\mathbf1^{\mathsf T}.
 \label{eq:helmert-frame}
\end{equation}

\begin{proposition}[Five-Kraus freedom behind the six-edge normalization]
\label{prop:five-kraus-freedom}
Equation~\eqref{eq:six-edge-normalization} holds if and only if there are
five matrices $B_1,\ldots,B_5$ with
$\sum_\mu B_\mu^*B_\mu=I$ and
$I,B_1,\ldots,B_5$ linearly independent such that
\begin{equation}
 \boxed{K_e=\frac1{\sqrt{18}}I
       +\sqrt{\frac23}\sum_{\mu=1}^5O_{e\mu}B_\mu.}
 \label{eq:six-edge-normal-form}
\end{equation}
Consequently the coarse edge channel has the form
\[
 \mathcal E_* =\frac13\id+\frac23\Psi_*,
 \qquad \Psi_*(\rho)=\sum_\mu B_\mu\rho B_\mu^*,
\]
and $C^*(K_e:e)=C^*(B_\mu:\mu)$.
\end{proposition}

\begin{proof}
Apply the orthogonal coefficient change whose first column is
$\mathbf1/\sqrt6$ and whose remaining columns are $O$.  The coherent-sum
constraint fixes the first transformed coefficient to $I/\sqrt3$; the other
five, $D_\mu=\sum_eO_{e\mu}K_e$, satisfy
$\sum_\mu D_\mu^*D_\mu=2I/3$.  Put
$B_\mu=\sqrt{3/2}\,D_\mu$ and invert the orthogonal change.
\end{proof}

\subsection{An exact nonidentifiability pair}

Let $\mathscr S$ be the twenty three-element subsets of
$\{1,\ldots,6\}$.  On $\C^{\mathscr S}$ define commuting Hermitian
involutions
\[
 H_e^{\rm c}|S\rangle=(2\mathbf1_{e\in S}-1)|S\rangle.
\]
Let $\Gamma_1,\ldots,\Gamma_5$ be five anticommuting Hermitian Clifford
matrices on $\C^4$, and on $\C^4\otimes\C^5$ set
\[
 H_e^{\rm q}=\sqrt{\frac65}\sum_{\mu=1}^5O_{e\mu}\Gamma_\mu\otimes I_5.
\]
For $b\in\{{\rm c},{\rm q}\}$ put
\begin{equation}
 K_e^b=\frac1{\sqrt{18}}I_{20}+\frac{i}{3}H_e^b.
 \label{eq:two-system-completions}
\end{equation}

\begin{proposition}[Record-complete nonidentifiability witness]
\label{prop:record-witness}
Both lists in \eqref{eq:two-system-completions} satisfy
\eqref{eq:six-edge-normalization}, and
\begin{equation}
 (K_e^b)^*K_e^b=I_{20}/6,
 \qquad
 \Tr[(K_e^b)^*K_f^b]
 =\begin{cases}10/3,&e=f,\\2/3,&e\ne f.\end{cases}
 \label{eq:matching-record-data}
\end{equation}
Yet
\begin{equation}
 \boxed{
 C^*(K_e^{\rm c})=\C^{20},
 \qquad
 C^*(K_e^{\rm q})=M_4(\C)\otimes I_5.}
 \label{eq:inequivalent-system-algebras}
\end{equation}
For the resolved instrument above, every uninterrupted word in the original
classical alphabet has the same effect in the two realizations.  A word with
$m\ge1$ accepted letters, $n_0$ no-response letters and one common accepted
record type $a$ has effect
\begin{equation}
 (1-\vartheta)^{n_0}(\vartheta/6)^m I_{20}\otimes Q_a;
 \label{eq:record-word-effect}
\end{equation}
if accepted record types disagree the effect is zero.  Moreover the two one-step opportunity channels have the same nonzero Choi
spectrum,
\begin{equation}
 \boxed{
 \spec_+(J_{\rm opp})=
 \left\{
 \left[\frac{20}{3}(12-11\vartheta)\right]^{(1)},
 \left(\frac{20\vartheta}{3}\right)^{(3)},
 \left(\frac{8\vartheta}{3}\right)^{(20)}
 \right\}.}
 \label{eq:matching-Choi-spectrum}
\end{equation}
Thus the complete classical record channel at every word length is identical
although the system algebras differ.
\end{proposition}

\begin{proof}
For the commuting family, counting three-subsets gives zero mean and pair
correlation $-1/5$.  For the Clifford family the row vectors of
$\sqrt{6/5}\,O$ are unit vectors with the same pair correlations, and
Clifford anticommutation gives $(H_e^{\rm q})^2=I$.  Hence both families
satisfy the same normalization and Gram relations, and each $\sqrt6K_e^b$
is unitary.  The twenty sign patterns distinguish all basis vectors in the
commuting realization, giving $\C^{20}$.  The five Clifford generators span
and generate $M_4(\C)\otimes I_5$.  Finally, every classical word effect is
$L_w^*L_w$; because each accepted system factor is $6^{-1/2}$ times a
unitary and the record projectors are orthogonal, this gives
\eqref{eq:record-word-effect} in both realizations.

For the one-step Choi spectrum, the Gram of the twenty-four accepted Choi
vectors is $G_{24}=G_6\otimes I_4$ with
$G_6=(8/3)I_6+(2/3)\mathbf1\mathbf1^{\mathsf T}$.  It has eigenvalue
$20/3$ on the uniform edge direction with four record directions and
$8/3$ on the twenty edge-contrast directions.  The coefficient matrix
$M_\vartheta=\vartheta I+(1-\vartheta)cc^*$ acts by
$12-11\vartheta$ on the global uniform direction and by $\vartheta$ on its
orthogonal complement.  The nonzero Choi eigenvalues are therefore those of
$M_\vartheta^{1/2}G_{24}M_\vartheta^{1/2}$, giving
\eqref{eq:matching-Choi-spectrum}.
\end{proof}

\subsection{Ten coherent pullbacks reconstruct the six coefficients}

Enumerate the twenty-four accepted operators $A_{e,a}$ as $A_j$ and let
$\mathsf B$ synthesize their Choi vectors $|A_j\rangle\!\rangle$.  In
$\C^{24}$ put
\[
 c=\mathbf1/\sqrt2,
 \qquad
 M_\vartheta=\vartheta I+(1-\vartheta)cc^*.
\]
Let $C_\nu$ be the matrices synthesized by the columns of
$\mathsf B M_\vartheta^{1/2}$ and define the minimal Stinespring isometry
\[
 V_\vartheta=\sum_{\nu=1}^{24}C_\nu\otimes|\nu\rangle,
 \qquad
 \mathcal P_{\vartheta,*}(\rho)
 =\Tr_{\rm sys}(V_\vartheta\rho V_\vartheta^*).
\]
For $z\in\C^{24}$ write
$C(z)=(I\otimes\langle z|)V_\vartheta$.  Then
$\mathcal P_\vartheta^*(|z\rangle\langle w|)=C(z)^*C(w)$.

For the twenty-four coefficient labels $(e,a)$ define
\[
 u_{e,a}=1/\sqrt{24},
 \qquad (r_\mu)_{e,a}=O_{e\mu}/2,
 \qquad \langle u,r_\mu\rangle=0,
\]
and the fixed Hermitian environment quadratures
\begin{equation}
 X_\mu=|u\rangle\langle r_\mu|+|r_\mu\rangle\langle u|,
 \qquad
 Y_\mu=-i|u\rangle\langle r_\mu|+i|r_\mu\rangle\langle u|.
 \label{eq:ten-quadratures}
\end{equation}
Let $\mathcal P_\vartheta^*$ denote the Heisenberg complementary pullback of
the normalized coefficient frame and write
$\mathsf X_\mu=\mathcal P_\vartheta^*(X_\mu)$,
$\mathsf Y_\mu=\mathcal P_\vartheta^*(Y_\mu)$.  Set
$h=12-11\vartheta$ and $D_\mu=\sum_eO_{e\mu}K_e$.

\begin{proposition}[Explicit ten-panel reconstruction]
\label{prop:ten-panel-reconstruction}
The ten Hermitian operator-valued pullbacks determine all six system
coefficients through
\begin{equation}
 \boxed{
 D_\mu\otimes I_{H_R}
 =\frac{2\sqrt3}{\sqrt{\vartheta h}}
   (\mathsf X_\mu+i\mathsf Y_\mu),
 \qquad
 K_e=\frac{I}{\sqrt{18}}+\sum_\mu O_{e\mu}D_\mu.}
 \label{eq:ten-panel-reconstruction}
\end{equation}
Thus this finite operator-valued export determines the represented system
algebra and commutant.
\end{proposition}

\begin{proof}
In the normalized coefficient frame the uniform environment vector is an
eigenvector of the coefficient Gram with eigenvalue $h$, while each
$r_\mu$ has eigenvalue $\vartheta$.  Hence
$C(u)=\sqrt{h/12}\,I$ and
$C(r_\mu)=\sqrt\vartheta D_\mu\otimes I_{H_R}/2$.
Since $X_\mu+iY_\mu=2|u\rangle\langle r_\mu|$,
Equation~\eqref{eq:complementary-pullback} gives the first reconstruction
formula; orthogonality of $O$ gives the second.
\end{proof}

\begin{proposition}[Robust commutant stopping test]
\label{prop:reconstruction-error}
If the ten pullbacks have Hilbert--Schmidt errors
$\epsilon_{X,\mu}$ and $\epsilon_{Y,\mu}$, then
\begin{equation}
 \left(\sum_e\|\widehat K_e-K_e\|_{\HS}^2\right)^{1/2}
 \le e_K:=\frac{2\sqrt3}{\sqrt{\vartheta h}}
 \left[\sum_\mu(\epsilon_{X,\mu}+\epsilon_{Y,\mu})^2\right]^{1/2}.
 \label{eq:appendix-reconstruction-error}
\end{equation}
For the star-closed stacked commutator map,
\begin{equation}
 \|\widehat{\mathscr D}-\mathscr D\|_{2\to2}\le2\sqrt2\,e_K.
 \label{eq:appendix-commutator-error}
\end{equation}
Consequently a measured least singular value
$\widehat s>2\sqrt2\,e_K$ on an independently verified protected complement
certifies absence of any further commutant direction there.
\end{proposition}

\begin{proof}
Equation~\eqref{eq:ten-panel-reconstruction} and the triangle inequality give
the $D_\mu$ errors.  Orthogonality of $O$ preserves the sum of squared
coefficient errors.  For unit Hilbert--Schmidt $X$, each coefficient error
changes a commutator by at most twice its operator norm; retaining both a
coefficient and its adjoint gives the factor $2\sqrt2$.  The last statement
is the standard least-singular-value perturbation inequality.
\end{proof}

\begin{remark}
The ten-panel theorem is a sufficient finite tomography statement, not an
assertion that these coherent quadratures are automatically present in every
microscopic source.  If only the classical resolved record is supplied,
Proposition~\ref{prop:record-witness} shows that the acted-on system algebra
can remain underdetermined.  Conversely, once the operator-valued pullbacks
are physically available on a common frame, the reconstruction is exact and
feeds directly into the commutant audit of the main text.
\end{remark}

\end{document}
